% !TEX program = xelatex
% ============================================================================
% One Math Book — 高中数学（十至十二年级）
% Simplified Chinese edition of one_math_book_2_high_school.tex.
%
% Build:  latexmk one_math_book_2_high_school_zh.tex
% Output: build/one_math_book_2_high_school_zh.pdf
% ============================================================================
\documentclass[11pt,a4paper,twoside,openany]{book}

\newcommand{\booklang}{zh}
\usepackage{styles/onemath}
% Give TeX room around unbreakable Latin and math runs in Chinese lines.
\setlength{\emergencystretch}{3em}

% Print, at the end of every exercise and weekend problem, a hyperlinked
% "Solution p. <n>" pointer to its solution in the appendix.
\omsolutionlinks

% Brand the running header: rosette + ONE-COURSE.COM (linking to the site)
% above the chapter mark and the book title, on even pages; page number at the
% foot. See \ombrandheader in styles/onemath.sty.
\ombrandheader

\newcommand{\bookdate}{\today}
\newcommand{\bookversion}{}
\IfFileExists{version.tex}{\input{version.tex}}{}

\newcommand{\bookline}{第 2 册：高中数学}

\title{One Math Book}
\author{One Math Book 贡献者}
\date{\bookdate}

\begin{document}

\frontmatter
% One Course cover (mockup 81): pale blue-gray header band with the rosette
% sitting on its lower edge, ONE-COURSE.COM wordmark, four-color quad lines.
% Drawn as a full-page overlay: "remember picture" needs two LaTeX passes
% (latexmk handles the reruns). Coordinates in cm from the page's south west
% corner; the page is A4, 21 x 29.7.
%
% Text strings come from styles/lang/<lang>.tex (\bookmaintitle, …) and from
% the book entry file (\bookline).
\begin{titlepage}
  \begin{tikzpicture}[remember picture, overlay]
    \begin{scope}[shift={(current page.south west)}]
      % header band and rosette astride its lower edge
      \fill[ocPaleBlue] (0,24.6) rectangle (21,29.7);
      \ocQuadLine{0}{21}{24.36}{0.24}
      \ocRosette{10.5}{24.6}{2.4}{58}
      % wordmark, linking out to the site (\omSiteURL, styles/onemath.sty).
      % \href would recolor it (hyperref's urlcolor), so the ink is restated
      % inside the link.
      \node[text=ocInk, font=\sffamily\bfseries\LARGE] at (10.5,20.6)
        {\href{\omSiteURL}{\textcolor{ocInk}{\textls[210]{ONE-COURSE.COM}}}};
      % title block
      \node[text=ocDarkBlue, font=\fontsize{40}{44}\selectfont\bfseries]
        at (10.5,16.4) {\bookmaintitle};
      \node[text=ocInk, font=\Large\itshape] at (10.5,14.75)
        {\booktagline};
      \node[text=ocInk, font=\normalsize, align=center, text width=13cm]
        at (10.5,13.4)
        {\bookdescription};
      \ocQuadLine{6.6}{14.4}{12.0}{0.22}
      % this volume's name (defined by each book's entry file)
      \ifdefined\bookline
        \node[text=ocDarkBlue, font=\Large\bfseries] at (10.5,11.15)
          {\bookline};
      \fi
      % credits, version and date
      \node[text=ocInk, font=\normalsize] at (10.5,6.2)
        {\bookcontributors};
      \node[text=ocInk, font=\small] at (10.5,5.55)
        {Benjamin Virrion \;\textperiodcentered\; One Course};
      \node[text=ocInk, font=\small] at (10.5,4.7)
        {\ifdefempty{\bookversion}{}%
          {\bookversionlabel\ \bookversion\;\textperiodcentered\;}\bookdate};
      % bottom band
      \ocQuadLine{0}{21}{0}{1.0}
    \end{scope}
  \end{tikzpicture}
  \null
\end{titlepage}

% Reference page (verso of the title page): where the book lives and its licence.
\thispagestyle{empty}
\vspace*{\fill}
\begin{center}
\omcolophonsource\\[2pt]
\url{https://github.com/bvirrion/one-math-book}

\bigskip
\omcolophonlicense\\[2pt]
\url{https://creativecommons.org/licenses/by-nc-sa/4.0/}\\[6pt]
\omcolophoncredit
\end{center}
\vspace*{\fill}
\clearpage

\chapter{前言}

本书的愿望很简单：用一部前后连贯的数学书，涵盖从幼儿园到本科毕业的全部内容，
为世界各地的读者而写。

本书力求简洁而严谨：课程由定义、例子、命题和定理构成，只要在相应阶段可以理解，
就给出证明，随后是精心挑选的练习。不加证明给出的结果都明确标为“承认”。
所有练习的完整解答集中在书末。

每个学年构成本书的一个部分，每个部分再分为若干章。本项目目前涵盖一年级至十二年级
（6--18 岁），从小学一年级到中学最后一年，其中最后两年达到数学强化方向的水平。
学年越早，面向的读者越年幼：插图更多，演算步骤更详细，练习也更平缓。

\section*{如何阅读本书}

各陈述在每章内编号，并以颜色区分：定义为蓝色，定理为红色，命题及其同类为橙色，
方法（处理常见任务的简短步骤）为绿色。练习按难度分级，从直接应用的（$\star$）
到较难的（$\star\star\star$）。

\section*{参与贡献}

本书的源代码存放在一个公开的 git 仓库中：
\begin{center}
\url{https://github.com/bvirrion/one-math-book}
\end{center}
欢迎各种贡献，包括新的章节、勘误、更好的证明和更多的练习；
详见仓库根目录下的 \texttt{CONTRIBUTING.md}。

\tableofcontents

\mainmatter

% ============================================================================
% Grade 10 — shared structure (language-agnostic).
% Bodies: parts/grade-10/<file>.tex (en) or parts/grade-10/<lang>/<file>.tex
% ============================================================================
\part{\omstr{part.grade10}}

\ominput{grade-10}{01-numbers-and-sets}
\ominput{grade-10}{02-algebra-equations-inequalities}
\ominput{grade-10}{03-functions}
\ominput{grade-10}{04-reference-functions}
\ominput{grade-10}{05-coordinate-geometry}
\ominput{grade-10}{06-vectors}
\ominput{grade-10}{07-lines-and-systems}
\ominput{grade-10}{08-statistics}
\ominput{grade-10}{09-probability-and-sampling}

% ============================================================================
% Grade 11 — shared structure (language-agnostic).
% Bodies: parts/grade-11/<file>.tex (en) or parts/grade-11/<lang>/<file>.tex
% ============================================================================
\part{\omstr{part.grade11}}

\ominput{grade-11}{01-quadratic-functions}
\ominput{grade-11}{02-functions-variations}
\ominput{grade-11}{03-differentiation}
\ominput{grade-11}{04-sequences}
\ominput{grade-11}{05-trigonometry}
\ominput{grade-11}{06-vectors-lines}
\ominput{grade-11}{07-scalar-product}
\ominput{grade-11}{08-descriptive-statistics}
\ominput{grade-11}{09-probability-random-variables}
\ominput{grade-11}{10-binomial-distribution}

% ============================================================================
% Grade 12 — shared structure (language-agnostic).
% Bodies: parts/grade-12/<file>.tex (en) or parts/grade-12/<lang>/<file>.tex
% ============================================================================
\part{\omstr{part.grade12}}

\ominput{grade-12}{01-sequences}
\ominput{grade-12}{02-limits-continuity}
\ominput{grade-12}{03-differentiation-convexity}
\ominput{grade-12}{04-exponential-logarithm}
\ominput{grade-12}{05-trigonometric-functions}
\ominput{grade-12}{06-integration}
\ominput{grade-12}{07-differential-equations}
\ominput{grade-12}{08-combinatorics}
\ominput{grade-12}{09-complex-numbers}
\ominput{grade-12}{10-arithmetic}
\ominput{grade-12}{11-matrices-graphs}
\ominput{grade-12}{12-space-geometry}
\ominput{grade-12}{13-conditional-probability}
\ominput{grade-12}{14-random-variables-binomial}
\ominput{grade-12}{15-sums-lln}
\ominput{grade-12}{16-continuous-distributions}


\appendix
\part{\omsolutionspart}
% Solutions to the Grade 10 exercises — one file per chapter.
\chapter{Grade 10 (Age 15--16)}

\section*{Chapter \ref{ch:g10:numbers} --- Numbers and Sets of Numbers}

\begin{solution}{exo:g10:numbers:1}
$\frac{15}{3} = 5 \in \N$. \quad $-7 \in \Z$. \quad
$\frac{22}{7} \in \Q$ (it is not an integer: $22 = 7 \times 3 + 1$).
\quad $\sqrt 9 = 3 \in \N$. \quad $\sqrt{10} \in \R$ (irrational, since
$10$ is not the square of a rational --- admitted here, in the spirit of
\cref{thm:g10:numbers:sqrt2}). \quad $-2.4 = -\frac{24}{10} \in \Q$.
\quad $0 \in \N$.
\end{solution}

\begin{solution}{exo:g10:numbers:2}
(a) $\intco{-2}{5}$: filled dot at $-2$, hollow dot at $5$.
(b) $\intoo{3}{+\infty}$: hollow dot at $3$, shading to the right.
(c) $\intoc{-\infty}{-1}$: shading from the left up to a filled dot at
$-1$.
(d) ``distance from $x$ to $2$ at most $3$'' means $\abs{x - 2} \leq 3$,
i.e.\ $x \in \intcc{-1}{5}$: filled dots at $-1$ and $5$.
\end{solution}

\begin{solution}{exo:g10:numbers:3}
(a) The two intervals overlap between $0$ and $2$:
$I \cap J = \intcc{0}{2}$ ($0 \in J$ and $2 \in I$, and both belong to the
other set as well), and $I \cup J = \intco{-3}{4}$.

(b) $I \cap J$ is the set of $x$ with $-2 \leq x$ and $x < 1$:
$\intco{-2}{1}$. The union covers everything: $I \cup J = \R$.
\end{solution}

\begin{solution}{exo:g10:numbers:4}
$\abs{-6} = 6$; \quad $\abs{4 - 9} = \abs{-5} = 5$; \quad
$\abs{-3 - 5} = \abs{-8} = 8$; \quad
$\sqrt 2 > 1$ so $\abs{\sqrt2 - 1} = \sqrt2 - 1$; \quad
$1 - \sqrt2 < 0$ so $\abs{1 - \sqrt2} = \sqrt2 - 1$ as well (a number and
its opposite have the same absolute value).
\end{solution}

\begin{solution}{exo:g10:numbers:5}
$\abs{x} = 5$: distance $5$ from $0$, so $x = 5$ or $x = -5$.

$\abs{x - 1} = 3$: distance $3$ from $1$, so $x = 4$ or $x = -2$.

$\abs{x - 4} \leq 1$: distance at most $1$ from $4$, so
$x \in \intcc{3}{5}$.

$\abs{x + 2} < 3$: distance strictly less than $3$ from $-2$, so
$x \in \intoo{-5}{1}$.
\end{solution}

\begin{solution}{exo:g10:numbers:6}
Each interval is described by its center $a$ (midpoint of the bounds) and
radius $r$ (half the length).

$\intcc{2}{8}$: $a = 5$, $r = 3$, so $\abs{x - 5} \leq 3$.

$\intoo{-5}{1}$: $a = -2$, $r = 3$, open bounds, so $\abs{x + 2} < 3$.

$\intcc{-7}{-3}$: $a = -5$, $r = 2$, so $\abs{x + 5} \leq 2$.
\end{solution}

\begin{solution}{exo:g10:numbers:7}
Let $x = 0.272727\dots$ Then $100x = 27.2727\dots$, and subtracting:
\[
100x - x = 27.2727\dots - 0.2727\dots = 27,
\]
so $99x = 27$ and $x = \frac{27}{99} = \frac{3}{11}$, a quotient of
integers: $x$ is rational.
\end{solution}

\begin{solution}{exo:g10:numbers:8}
\emph{1. True:} adding integers (positive or negative whole numbers)
always produces an integer.

\emph{2. False:} $\frac{1}{2}$ is a quotient of the integers $1$ and $2$
and is not an integer.

\emph{3. True:} $\frac pq + \frac{p'}{q'} = \frac{pq' + p'q}{qq'}$ is
again a quotient of integers (with nonzero denominator).

\emph{4. True:} suppose $r$ is rational, $t$ irrational, and $r + t = s$
were rational. Then $t = s - r$ would be a difference of two rationals,
hence rational (by 3, applied with $-r$) --- contradiction. So $r + t$ is
irrational.
\end{solution}

\begin{solution}{exo:g10:numbers:9}
Multiplying $1.414 \leq \sqrt2 \leq 1.415$ by $2 > 0$:
$2.828 \leq 2\sqrt2 \leq 2.830$.

Adding $3$: $4.414 \leq \sqrt2 + 3 \leq 4.415$.

Multiplying by $-1 < 0$ reverses the inequalities:
$-1.415 \leq -\sqrt2 \leq -1.414$.
\end{solution}

\begin{solution}{exo:g10:numbers:10}
First the auxiliary fact. Divide $p$ by $3$: the remainder is $0$, $1$ or
$2$, i.e.\ $p = 3k$, $p = 3k+1$ or $p = 3k+2$. Squaring:
\[
(3k)^2 = 3(3k^2), \quad
(3k+1)^2 = 3(3k^2 + 2k) + 1, \quad
(3k+2)^2 = 3(3k^2 + 4k + 1) + 1 .
\]
Only the first is a multiple of $3$: if $p^2$ is a multiple of $3$, then
$p$ is too.

Now suppose $\sqrt3 = \frac pq$ fully simplified. Squaring gives
$p^2 = 3q^2$, so $p^2$ is a multiple of $3$, so $p = 3k$. Then
$9k^2 = 3q^2$, so $q^2 = 3k^2$ and $q$ is also a multiple of $3$ --- but
then the fraction $\frac pq$ was not fully simplified: contradiction. So
$\sqrt3$ is irrational.
\end{solution}

\section*{Chapter \ref{ch:g10:algebra} --- Algebra: Equations and Inequalities}

\begin{solution}{exo:g10:algebra:1}
$3(2x-5) - 2(x-7) = 6x - 15 - 2x + 14 = 4x - 1$.

$(x+4)(2x-1) = 2x^2 - x + 8x - 4 = 2x^2 + 7x - 4$.

$(3x-2)^2 = 9x^2 - 12x + 4$.

$(5-x)(5+x) = 25 - x^2$.
\end{solution}

\begin{solution}{exo:g10:algebra:2}
$x^2 - 49 = (x+7)(x-7)$ (difference of squares).

$x^2 + 10x + 25 = (x+5)^2$.

$7x^2 - 14x = 7x(x - 2)$ (common factor).

$16x^2 - 8x + 1 = (4x - 1)^2$.
\end{solution}

\begin{solution}{exo:g10:algebra:3}
$4x + 9 = x - 3$: $3x = -12$, so $x = -4$.

$\frac{2x-1}{3} = 5$: $2x - 1 = 15$, so $x = 8$.

$2(x-3) = 2x + 1$ expands to $2x - 6 = 2x + 1$, i.e.\ $-6 = 1$: false for
every $x$. No solution --- the two sides define parallel lines that never
meet.
\end{solution}

\begin{solution}{exo:g10:algebra:4}
$(x+5)(2x-8) = 0$: $x = -5$ or $x = 4$.

$x^2 - 16 = (x-4)(x+4) = 0$: $x = 4$ or $x = -4$.

$x^2 + 6x + 9 = (x+3)^2 = 0$: the single solution $x = -3$.
\end{solution}

\begin{solution}{exo:g10:algebra:5}
$3x - 5 \leq 7$: $3x \leq 12$, so $x \leq 4$: solution set
$\intoc{-\infty}{4}$.

$-4x + 1 > 9$: $-4x > 8$, dividing by $-4$ reverses the inequality:
$x < -2$: solution set $\intoo{-\infty}{-2}$.

$2(x-1) \geq 5x + 4$: $2x - 2 \geq 5x + 4$, so $-3x \geq 6$, so
$x \leq -2$: solution set $\intoc{-\infty}{-2}$.
\end{solution}

\begin{solution}{exo:g10:algebra:6}
The common factor is $(x-1)$:
\[
E = (x-1)\bigl[(x+3) - (2x - 5)\bigr] = (x-1)(x + 3 - 2x + 5)
= (x-1)(8 - x).
\]
Then $E = 0$ when $x = 1$ or $x = 8$ (zero-product rule).
\end{solution}

\begin{solution}{exo:g10:algebra:7}
$x^2 = 5x$: $x^2 - 5x = x(x - 5) = 0$, so $x = 0$ or $x = 5$ (never
divide by $x$: the solution $0$ would be lost).

$(2x+1)^2 = (x-4)^2$: bring to the left and factor the difference of
squares:
\[
(2x+1)^2 - (x-4)^2
= \bigl[(2x+1) + (x-4)\bigr]\bigl[(2x+1) - (x-4)\bigr]
= (3x - 3)(x + 5).
\]
Zero-product rule: $x = 1$ or $x = -5$.
\end{solution}

\begin{solution}{exo:g10:algebra:8}
$(x+2)(4-x) > 0$: the factor $x + 2$ vanishes at $-2$ (positive after),
$4 - x$ vanishes at $4$ (positive before). The product is positive
exactly when both factors are positive, i.e.\ on $\intoo{-2}{4}$.

$\frac{2x-6}{x+1} \geq 0$: numerator zero at $3$, denominator zero at
$-1$ (forbidden). \omterm{def:g10:algebra:signtable}{Sign table}: the quotient is positive when both parts
have the same sign, i.e.\ for $x < -1$ or $x > 3$; it vanishes at $x = 3$.
Solution set: $\intoo{-\infty}{-1} \cup \intco{3}{+\infty}$.
\end{solution}

\begin{solution}{exo:g10:algebra:9}
After $x$ months the first service costs $9x + 15$ and the second $12x$.
The first is cheaper when
\[
9x + 15 < 12x
\ \Longleftrightarrow\
15 < 3x
\ \Longleftrightarrow\
x > 5 .
\]
From the $6$th month on, the first service is the cheaper choice overall.
\end{solution}

\begin{solution}{exo:g10:algebra:10}
Forbidden value: $x = 1$. For $x \neq 1$, multiply both sides by $x - 1$:
\[
x + 3 = 2(x - 1) = 2x - 2
\ \Longleftrightarrow\
x = 5 .
\]
Since $5 \neq 1$, the solution is $x = 5$. Check:
$\frac{5+3}{5-1} = \frac84 = 2$.
\end{solution}

\begin{solution}{exo:g10:algebra:11}
$(n+1)^2 - n^2 = n^2 + 2n + 1 - n^2 = 2n + 1$, which is odd for every
\omterm{def:g10:numbers:sets}{integer} $n$.

Two consecutive odd numbers can be written $2n - 1$ and $2n + 1$. Then
\[
(2n+1)^2 - (2n-1)^2
= \bigl[(2n+1)+(2n-1)\bigr]\bigl[(2n+1)-(2n-1)\bigr]
= 4n \times 2 = 8n,
\]
a multiple of $8$.
\end{solution}

\begin{solution}{pb:g10:algebra:1}
\textbf{1.} $x^2 - 10x + 25 = (x - 5)^2$;
$9x^2 - 49 = (3x - 7)(3x + 7)$;
$x^2 + 6x + 9 - 4 = (x + 3)^2 - 2^2 = (x + 1)(x + 5)$.

\textbf{2.} $x = -1$ or $x = -5$; \quad $x = \frac73$ or
$x = -\frac73$.

\textbf{3.} $(5 - t)(5 + t) = 25 - t^2 = 21$ gives $t^2 = 4$,
$t = 2$: the numbers are $3$ and $7$ (sum $10$, product $21$:
check).

\textbf{4.} $(7 - t)(7 + t) = 49 - t^2 = 33$: $t = 4$, numbers
$3$ and $11$.

\textbf{5.} $\left(\frac s2 - t\right)\left(\frac s2 + t\right)
= \left(\frac s2\right)^2 - t^2 = p$ forces
$t^2 = \left(\frac s2\right)^2 - p$. A solution exists exactly
when this is $\geq 0$, i.e.\ when
$p \leq \left(\frac s2\right)^2$ --- the product of two numbers
of given sum can never beat the square of the half-sum
(queen Dido's fence problem, in the Middle School volume knew).

\textbf{6.} $(x - a)(x - b) = x^2 - (a + b)x + ab$: the
\omterm{def:g10:algebra:equation}{equation} $x^2 - sx + p = 0$ says precisely ``$x$ is one of two
numbers with sum $s$ and product $p$''.

\textbf{7.} $(x-5)^2 - 4 = x^2 - 10x + 25 - 4 = x^2 - 10x +
21$. Difference of squares:
$(x - 5 - 2)(x - 5 + 2) = (x - 7)(x - 3) = 0$: solutions $3$
and $7$ --- the scribe's numbers of question 3.

\textbf{8.} $x^2 + 6x - 7 = (x + 3)^2 - 16 = 0$ gives
$(x + 3)^2 = 16$, $x + 3 = \pm 4$: $x = 1$ or $x = -7$.

\textbf{9.} $x^2 + 4x + 5 = (x + 2)^2 + 1 \geq 1 > 0$ for every
$x$: no real solution. In the form $(x + h)^2 + k$: solutions
exist exactly when $k \leq 0$ (then $(x+h)^2 = -k$ has the two
roots $-h \pm \sqrt{-k}$); when $k > 0$ the expression stays
positive.

\textbf{10.} The square of side $x$ (area $x^2$) with two
$3 \times x$ rectangles ($6x$ in all) forms an L of area
$x^2 + 6x = 7$; the missing $3 \times 3$ corner (area $9$)
completes a big square of side $x + 3$. Adding $9$ to both
sides: $(x + 3)^2 = 16$, so $x + 3 = 4$ (lengths are positive)
and $x = 1$. The algebra and the picture are the same act.

\textbf{11.} $(x-3)(x+2)$: negative between the roots. Solution
of $< 0$: $x \in \intoo{-2}{3}$.

\textbf{12.} $x^2 - 10x + 21 \leq 0$, i.e.\
$(x - 3)(x - 7) \leq 0$: between the roots, ends included:
$x \in \intcc{3}{7}$.

\textbf{13.} Zero at $x = 1$, forbidden at $x = -2$; the
quotient is positive outside the roots:
$x \in \intoo{-\infty}{-2} \cup \intco{1}{+\infty}$.

\textbf{14.} $25 - (x - 5)^2 = 25 - x^2 + 10x - 25 =
x(10 - x)$: identity verified. Since $(x-5)^2 \geq 0$ always,
$A(x) \leq 25$, with equality exactly when $x = 5$: the
$5 \times 5$ square pen, maximum $25$ m$^2$ --- the grade-6
table, now a two-line proof.

\textbf{15.} For $x > 0$, multiply the claim by $x$:
$x^2 + 1 \geq 2x$, i.e.\ $x^2 - 2x + 1 = (x - 1)^2 \geq 0$ ---
true, with equality only at $x = 1$. Dividing back by $x > 0$
preserves everything (\cref{prop:g10:algebra:ineqrules}).

\textbf{16.} Cross-multiplying $\frac1x = \frac{x}{1 - x}$
(all lengths positive): $x^2 = 1 - x$, i.e.\
$x^2 + x - 1 = 0$.

\textbf{17.} $x^2 + x - 1 = \left(x + \frac12\right)^2 -
\frac54 = 0$ gives $x + \frac12 = \pm\frac{\sqrt5}{2}$. The
positive root: $x = \frac{\sqrt5 - 1}{2} \approx 0.618$.

\textbf{18.} $\varphi = \frac{2}{\sqrt5 - 1} =
\frac{2(\sqrt5 + 1)}{(\sqrt5 - 1)(\sqrt5 + 1)}
= \frac{2(\sqrt5 + 1)}{5 - 1} = \frac{\sqrt5 + 1}{2}
\approx 1.618$.

\textbf{19.} $\varphi^2 = \frac{(\sqrt5 + 1)^2}{4} =
\frac{6 + 2\sqrt5}{4} = \frac{3 + \sqrt5}{2}$, and
$\varphi + 1 = \frac{\sqrt5 + 3}{2}$: equal. And
$\varphi - 1 = \frac{\sqrt5 - 1}{2} = x = \frac1\varphi$ (that
is how $\varphi$ was defined). The golden ratio is the number
whose square adds one and whose reciprocal subtracts one.

\textbf{20.} $\sqrt5$ is \omterm{ex:g10:numbers:classify}{irrational} (the irrationality weekend problem of the Middle School volume: $5$
is no perfect square), and adding $1$ then halving preserves
irrationality (rational operations, \cref{pb:g10:numbers:1}).
Ratios: $\frac53 \approx 1.667$, $\frac85 = 1.600$,
$\frac{13}{8} = 1.625$, $\frac{21}{13} \approx 1.615$:
alternately above and below, closing in on
$\varphi = 1.618\ldots$ --- the rhythm-counting numbers of
the rhythm-counting weekend problem of the Middle School volume secretly chase the golden section; the
chapters on sequences will catch them in the act.
\end{solution}

\section*{Chapter \ref{ch:g10:functions} --- Functions}

\begin{solution}{exo:g10:functions:1}
$f(0) = 1$; \quad $f(2) = 8 - 6 + 1 = 3$; \quad
$f(-1) = 2 + 3 + 1 = 6$; \quad
$f\!\left(\tfrac12\right) = 2 \times \tfrac14 - \tfrac32 + 1 = \tfrac12 -
\tfrac32 + 1 = 0$.
\end{solution}

\begin{solution}{exo:g10:functions:2}
$f$: no constraint, \omterm{def:g10:functions:function}{domain} $\R$.

$g$: requires $x + 4 \neq 0$, \omterm{def:g10:functions:function}{domain} all reals except $-4$.

$h$: requires $2x - 6 \geq 0$, i.e.\ $x \geq 3$: \omterm{def:g10:functions:function}{domain}
$\intco{3}{+\infty}$.

$k$: $x^2 + 1 \geq 1 > 0$ never vanishes, \omterm{def:g10:functions:function}{domain} $\R$.
\end{solution}

\begin{solution}{exo:g10:functions:3}
\emph{\omterm{def:g10:functions:function}{Preimages} of $0$:} $x^2 - 2x = x(x-2) = 0$, so $0$ and $2$.

\emph{\omterm{def:g10:functions:function}{Preimages} of $3$:} $x^2 - 2x = 3$, i.e.\ $x^2 - 2x - 3 = 0$, i.e.\
$(x-3)(x+1) = 0$ (check by \omterm{def:g10:algebra:expand}{expanding}), so $3$ and $-1$.

\emph{\omterm{def:g10:functions:function}{Preimages} of $-1$:} $x^2 - 2x + 1 = (x - 1)^2 = 0$, so the single
\omterm{def:g10:functions:function}{preimage} $1$.
\end{solution}

\begin{solution}{exo:g10:functions:4}
$f(-1) = 1 + 2 + 1 = 4$: yes, $A(-1, 4)$ is on the \omterm{def:g10:functions:graph}{graph}.
$f(2) = 4 - 4 + 1 = 1$: yes, $B(2, 1)$ is on the \omterm{def:g10:functions:graph}{graph} too.
\end{solution}

\begin{solution}{exo:g10:functions:5}
\emph{1.} The largest value in the table is $4 = g(3)$ (\omterm{def:g10:functions:extrema}{maximum}), the
smallest is $-2 = g(0)$ (\omterm{def:g10:functions:extrema}{minimum}).

\emph{2.} On $\intcc{-3}{0}$ the \omterm{def:g10:functions:function}{function} decreases, and
$-1 < -0.5$, so $g(-1) \geq g(-0.5)$.

\emph{3.} On each \omterm{def:g10:numbers:interval}{interval} of monotonicity, $g$ takes each value at most
once, so $g(x) = 0$ has at most one solution per \omterm{def:g10:numbers:interval}{interval}: at most $3$
solutions in total. (Here $0$ lies between $-2$ and $1$, between $-2$ and
$4$, and between $0$ and $4$, so there are exactly three.)
\end{solution}

\begin{solution}{exo:g10:functions:6}
Let $u < v$. Then
\[
f(v) - f(u) = (-2v + 7) - (-2u + 7) = -2(v - u) < 0,
\]
since $v - u > 0$. So $f(v) < f(u)$: $f$ is strictly \omterm{def:g10:functions:variations}{decreasing} on $\R$.
\end{solution}

\begin{solution}{exo:g10:functions:7}
Complete the square: $(x+3)^2 = x^2 + 6x + 9$, so
\[
f(x) = x^2 + 6x + 11 = (x + 3)^2 + 2 .
\]
For every $x$, $(x+3)^2 \geq 0$, so $f(x) \geq 2$; and
$f(-3) = 0 + 2 = 2$. The \omterm{def:g10:functions:extrema}{minimum} of $f$ is $2$, attained at $x = -3$.
\end{solution}

\begin{solution}{exo:g10:functions:8}
\emph{1.} Length $+$ width $= 20$, so the length is $20 - x$ and
\[
A(x) = x(20 - x).
\]
Both dimensions must be positive: $0 < x < 20$.

\emph{2.} $A(4) = 64$, $A(8) = 96$, $A(10) = 100$, $A(12) = 96$,
$A(16) = 64$. The values rise until $x = 10$ then fall symmetrically:
the area seems largest for $x = 10$, i.e.\ for a \emph{square} garden.
(\cref{ch:g11:quad} proves this conjecture.)
\end{solution}

\begin{solution}{exo:g10:functions:9}
\emph{1.} For every $x$: $x^2 + 1 \geq 1 > 0$, so $f(x) = \frac{1}{x^2+1}$
is positive, and since $x^2 + 1 \geq 1$, taking inverses (both sides
positive) gives $f(x) \leq 1$.

\emph{2.} \omterm{def:g10:functions:extrema}{Maximum}: $f(0) = 1$ and $f(x) \leq 1$ for all $x$, so $1$ is the
\omterm{def:g10:functions:extrema}{maximum}, attained at $0$. \omterm{def:g10:functions:extrema}{Minimum}: $f(x) > 0$ always, but no value of $x$
achieves $0$ (a quotient of nonzero numbers is nonzero), and $f$ takes
values as close to $0$ as we like for large $x$; so $f$ has no \omterm{def:g10:functions:extrema}{minimum}.
\end{solution}

\begin{solution}{exo:g10:functions:10}
\emph{1.} $f(v) - f(u) = v^2 - u^2 = (v-u)(v+u)$. If $0 \leq u < v$, then
$v - u > 0$ and $v + u > 0$ (sum of a nonnegative and a positive number),
so $f(v) - f(u) > 0$: $f$ is strictly \omterm{def:g10:functions:variations}{increasing} on $\intco{0}{+\infty}$.

\emph{2.} If $u < v \leq 0$, then $v - u > 0$ still, but now
$v + u < 0$ (sum of a nonpositive and a negative number), so
$f(v) - f(u) < 0$: $f$ is strictly \omterm{def:g10:functions:variations}{decreasing} on $\intoc{-\infty}{0}$.
\end{solution}

\section*{Chapter \ref{ch:g10:reffunc} --- Reference Functions}

\begin{solution}{exo:g10:reffunc:1}
$f$: \omterm{def:g10:reffunc:affine}{slope} $3$, intercept $-2$, \omterm{def:g10:functions:variations}{increasing}. \quad
$g$: \omterm{def:g10:reffunc:affine}{slope} $-\frac12$, intercept $4$, \omterm{def:g10:functions:variations}{decreasing}. \quad
$h$: \omterm{def:g10:reffunc:affine}{slope} $0$, intercept $7$, constant. \quad
$k$: \omterm{def:g10:reffunc:affine}{slope} $-1$, intercept $0$, \omterm{def:g10:functions:variations}{decreasing} (and linear).
\end{solution}

\begin{solution}{exo:g10:reffunc:2}
Through $(2,1)$ and $(5,10)$: \omterm{def:g10:reffunc:affine}{slope} $m = \frac{10 - 1}{5 - 2} = 3$; then
$1 = 3 \times 2 + p$ gives $p = -5$: $f(x) = 3x - 5$.

Through $(-1,4)$ and $(3,-4)$: \omterm{def:g10:reffunc:affine}{slope} $m = \frac{-4 - 4}{3 - (-1)} =
\frac{-8}{4} = -2$; then $4 = -2 \times (-1) + p$ gives $p = 2$:
$f(x) = -2x + 2$.
\end{solution}

\begin{solution}{exo:g10:reffunc:3}
$(3.1)^2 < (3.2)^2$: the square increases on positives.

$(-2.7)^2 < (-2.8)^2$: the square decreases on negatives and
$-2.8 < -2.7$.

$\frac{1}{5.1} > \frac{1}{5.2}$: the inverse decreases on positives.

$\sqrt{17} > \sqrt{15}$: the square root increases.
\end{solution}

\begin{solution}{exo:g10:reffunc:4}
$x^2 = 16$: two solutions, $x = 4$ and $x = -4$ (the horizontal line
$y = 16$ cuts the parabola twice).

$x^2 \leq 16$: the parabola is below the line between the two
\omterm{def:g10:numbers:interunion}{intersection} points: $x \in \intcc{-4}{4}$.

$x^2 > 9$: the parabola is strictly above $y = 9$ outside
$\intcc{-3}{3}$: $x \in \intoo{-\infty}{-3} \cup \intoo{3}{+\infty}$.
\end{solution}

\begin{solution}{exo:g10:reffunc:5}
For $0 < x < 1$: multiplying $x < 1$ repeatedly by $x$ gives
$x^3 < x^2 < x$, and \cref{prop:g10:reffunc:compare} gives $x < \sqrt x$,
so
\[
x^3 < x^2 < x < \sqrt x .
\]
Check with $x = 0.25$: $x^3 = 0.015625$, $x^2 = 0.0625$, $x = 0.25$,
$\sqrt x = 0.5$.
\end{solution}

\begin{solution}{exo:g10:reffunc:6}
$\frac1x = 2$ has the unique solution $x = \frac12$ (the hyperbola meets
the horizontal line $y = 2$ once, on the positive branch).

For $x > 0$: $\frac1x < 2$. Multiplying by $x > 0$: $1 < 2x$, so
$x > \frac12$: solution set $\intoo{\frac12}{+\infty}$.

If $x < 0$ is allowed: every negative $x$ satisfies $\frac1x < 0 < 2$,
so the full solution set is
$\intoo{-\infty}{0} \cup \intoo{\frac12}{+\infty}$.
\end{solution}

\begin{solution}{exo:g10:reffunc:7}
(a) On $\intcc{2}{5}$, all inputs are positive and squaring increases:
$4 \leq x^2 \leq 25$.

(b) On $\intcc{-3}{-1}$, squaring decreases:
the largest square comes from $-3$: $1 \leq x^2 \leq 9$.

(c) On $\intcc{-2}{3}$, the \omterm{def:g10:numbers:interval}{interval} contains $0$, where the square
reaches its \omterm{def:g10:functions:extrema}{minimum} $0$; the \omterm{def:g10:functions:extrema}{maximum} is the larger of $(-2)^2 = 4$ and
$3^2 = 9$. So $0 \leq x^2 \leq 9$.
\end{solution}

\begin{solution}{exo:g10:reffunc:8}
First company: $f(x) = 1.5x + 4$; second: $g(x) = 2x$. The first is
cheaper when
\[
1.5x + 4 < 2x
\ \Longleftrightarrow\
4 < 0.5x
\ \Longleftrightarrow\
x > 8 .
\]
Beyond $8$ km, the first company is cheaper; at exactly $8$ km both charge
$16$.
\end{solution}

\begin{solution}{exo:g10:reffunc:9}
Both sides are nonnegative, so the inequality is equivalent to the one
between their squares:
\[
\left(\sqrt{a+b}\right)^2 = a + b
\qquad\text{and}\qquad
\left(\sqrt a + \sqrt b\right)^2 = a + 2\sqrt a \sqrt b + b .
\]
Since $2\sqrt a\sqrt b \geq 0$, the second square is at least the first,
which proves $\sqrt{a+b} \leq \sqrt a + \sqrt b$. Equality holds exactly
when $\sqrt a \sqrt b = 0$, i.e.\ when $a = 0$ or $b = 0$.
\end{solution}

\begin{solution}{exo:g10:reffunc:10}
\emph{1.} Put the right-hand side over a common denominator:
\[
2 + \frac{3}{x-1} = \frac{2(x-1) + 3}{x - 1} = \frac{2x + 1}{x - 1}
= f(x).
\]

\emph{2.} On $\intoo{1}{+\infty}$, as $x$ increases, $x - 1$ increases
and stays positive, so $\frac{3}{x-1}$ decreases (inverse \omterm{def:g10:functions:function}{function}, times
the positive constant $3$), so $f(x) = 2 + \frac{3}{x-1}$ decreases:
$f$ is strictly \omterm{def:g10:functions:variations}{decreasing} on $\intoo{1}{+\infty}$.
\end{solution}

\begin{solution}{pb:g10:reffunc:1}
\textbf{1.} $d(30) = 3^2 = 9$ m; $d(50) = 25$ m; $d(90) =
81$ m; $d(130) = 169$ m.

\textbf{2.} Doubling $v$ doubles $\frac v{10}$ and multiplies
its square by $4$: $d(100) = 100$ m $= 4 \times d(50)$. Twice
the speed, four times the distance --- the square \omterm{def:g10:functions:function}{function}'s
scaling.

\textbf{3.} $\left(\frac{v}{10}\right)^2 = 50$ gives
$\frac{v}{10} = \sqrt{50}$, so $v = 10\sqrt{50} \approx 71$
km/h. The square root answered; the reading is unambiguous
because speeds are positive, and on
$\intco{0}{+\infty}$ the square \omterm{def:g10:functions:function}{function} climbs strictly
(\cref{prop:g10:reffunc:sqrt}: one positive \omterm{def:g10:functions:function}{preimage}).

\textbf{4.} $d(31) - d(30) = 9.61 - 9 = 0.61$ m, while
$d(91) - d(90) = 82.81 - 81 = 1.81$ m: the same $+1$ km/h costs
three times more distance at high speed. The parabola not only
rises, it \emph{steepens} --- its rate of climb grows with $x$.

\textbf{5.} $D(50) = 14 + 25 = 39$ m; $D(130) = 36.4 + 169 =
205.4$ m. In town the affine (reaction) term is the bigger
share; on the highway the square term crushes it --- affine
grows steadily, squares run away.

\textbf{6.} $F(0.5) = 240$, $F(1) = 120$, $F(3) = 40$.

\textbf{7.} The inverse \omterm{def:g10:functions:function}{function} decreases: the further from
the pivot, the smaller the needed force --- with a long enough
lever arm ($d$ huge), any force, however small, balances any
load: Archimedes' boast lives in the hyperbola's tail, which
approaches $0$ without reaching it. And $d = 0$ is the
forbidden value: no arm, no lever --- the vertical asymptote.

\textbf{8.} $\frac{100}{4} = 25$ at $2$ m,
$\frac{100}{9} \approx 11$ at $3$ m, $1$ at $10$ m. At distance
$d$ the light has spread over a sphere of area $4\pi d^2$
(Archimedes' tombstone problem, in the Middle School volume): the same energy divided by an area
growing like $d^2$ --- hence the inverse \emph{square}.

\textbf{9.} $60$, $30$, $20$ euros per head for $n = 10, 20,
30$. For $\frac{600}{n} \leq 25$: $n \geq 24$ participants. The
hyperbola promises ever-cheaper shares as $n$ grows --- and
refuses ever to reach $0$: the bus is never free.

\textbf{10.} $a(x) = \frac{2x + 3}{x} = 2 + \frac3x$:
\omterm{def:g10:functions:variations}{decreasing} on $\intoo{0}{+\infty}$ (inverse \omterm{def:g10:functions:function}{function} shifted),
approaching the asymptote $y = 2$. Economically: spreading the
fixed $3$ euros over more posters pushes the unit cost down
towards the incompressible $2$ euros of materials --- economies
of scale, with a hard floor.

\textbf{11.} Earth: $a^3 = 1$, $T^2 = 1$. Mars:
$1.52^3 \approx 3.51$ and $1.88^2 \approx 3.53$. Jupiter:
$5.20^3 \approx 140.6$ and $11.86^2 \approx 140.7$. Kepler's
notice: $T^2 = a^3$ for every planet --- one law for the whole
sky.

\textbf{12.} $T = \sqrt{a^3}$: for Saturn
$\sqrt{9.54^3} = \sqrt{868} \approx 29.5$ years, against the
observed $29.4$: the law holds to a tenth.

\textbf{13.} Asteroid: $T = \sqrt{4^3} = \sqrt{64} = 8$ years.
Comet: $a^3 = 27^2 = 729 = 9^3$, so $a = 9$ AU.

\textbf{14.} Square ($T^2$), cube ($a^3$) and square root (to
extract $T$). Scaling: multiplying $a$ by $4$ multiplies $T$ by
$4\sqrt4 = 8$ --- as the asteroid confirms ($a$ four times
Earth's, $T$ eight times).

\textbf{15.} A table of numbers, matched against the shapes of
a few reference \omterm{def:g10:functions:function}{functions}, yielded a law of the universe
seventy years before anyone knew \emph{why} it held ---
recognizing the \omterm{def:g10:functions:function}{function} is half of science.

\textbf{16.} On $\intoo{0}{1}$: $x^2 < x < \sqrt x$; on
$\intoo{1}{+\infty}$: $\sqrt x < x < x^2$. Checks: at
$x = 0.25$: $0.0625 < 0.25 < 0.5$; at $x = 4$:
$2 < 4 < 16$.

\textbf{17.} For $x \geq 0$, both sides are $\geq 0$, so
$\sqrt x > x \iff x > x^2 \iff x(1 - x) > 0 \iff
x \in \intoo{0}{1}$.

\textbf{18.} At $x = 0.5$: $x^3 = 0.125 < x^2 = 0.25 < x =
0.5$. At $x = 2$: $2 < 4 < 8$: the order reverses. Crossings of
$x^2$ and $x^3$: $x^3 - x^2 = x^2(x - 1) = 0$: at $x = 0$ and
$x = 1$.

\textbf{19.} $10t^2 = 100\sqrt t$ gives $t^2 = 10\sqrt t$;
squaring, $t^4 = 100t$, so $t^3 = 100$ and
$t = \sqrt[3]{100} \approx 4.64$ years --- about $4$ years and
$8$ months. (Both schemes then pay about $215$ euros.) Moral:
roots sprint early, powers win every marathon.

\textbf{20.} Braking: the parabola's \emph{steepening} is the
danger of speed. Lever: the \emph{vertical asymptote} at $0$
and the long tail are the mechanic's advantage. Lamp: the
$d^2$ of the \emph{spreading sphere} sets the fading. Kepler:
the \emph{exponent} $\frac32$ is the solar system's tempo.
Savings A: the root's \emph{flattening} is the slow saver's
fate. Five \omterm{def:g10:functions:graph}{graphs}, five laws: nature's alphabet indeed.
\end{solution}

\section*{Chapter \ref{ch:g10:coordgeom} --- Coordinate Geometry}

\begin{solution}{exo:g10:coordgeom:1}
Midpoint: $\left(\frac{2 + (-4)}{2}, \frac{5 + 1}{2}\right) = (-1, 3)$.
Distance:
\[
AB = \sqrt{(-4 - 2)^2 + (1 - 5)^2} = \sqrt{36 + 16} = \sqrt{52}
= 2\sqrt{13}.
\]
\end{solution}

\begin{solution}{exo:g10:coordgeom:2}
$I$ is the midpoint of $[AB]$ means $\frac{3 + x_B}{2} = 1$ and
$\frac{-2 + y_B}{2} = 2$, so $x_B = -1$ and $y_B = 6$: $B(-1, 6)$.
\end{solution}

\begin{solution}{exo:g10:coordgeom:3}
$OP = \sqrt{16 + 1} = \sqrt{17}$; $OQ = \sqrt{9 + 4} = \sqrt{13}$;
$OR = \sqrt{0 + 25} = 5 = \sqrt{25}$. Since $13 < 17 < 25$, the square
root being increasing, $Q$ is the closest to the origin.
\end{solution}

\begin{solution}{exo:g10:coordgeom:4}
Squared lengths:
\begin{align*}
AB^2 &= (5-1)^2 + (4-2)^2 = 16 + 4 = 20, \\
BC^2 &= (7-5)^2 + (0-4)^2 = 4 + 16 = 20, \\
AC^2 &= (7-1)^2 + (0-2)^2 = 36 + 4 = 40 .
\end{align*}
$AB = BC$: the triangle is isosceles at $B$. Moreover
$AB^2 + BC^2 = 40 = AC^2$: by the converse of the Pythagorean theorem it
is right-angled at $B$ as well.
\end{solution}

\begin{solution}{exo:g10:coordgeom:5}
\emph{1.} Midpoint of $[AC]$:
$\left(\frac{-2+4}{2}, \frac{1+1}{2}\right) = (1, 1)$; midpoint of
$[BD]$: $\left(\frac{1+1}{2}, \frac{3+(-1)}{2}\right) = (1, 1)$. Same
midpoint: $ABCD$ is a parallelogram.

\emph{2.} $AB^2 = 3^2 + 2^2 = 13$ and $BC^2 = 3^2 + (-4)^2 = 25$:
$AB \neq BC$, not a rhombus.

\emph{3.} $AC^2 = 6^2 + 0^2 = 36$ and $BD^2 = 0^2 + (-4)^2 = 16$:
$AC \neq BD$, not a rectangle. $ABCD$ is a plain parallelogram.
\end{solution}

\begin{solution}{exo:g10:coordgeom:6}
Compute squared distances to $\Omega(2,1)$ and compare with $5^2 = 25$:
\[
\Omega A^2 = 3^2 + 4^2 = 25, \qquad
\Omega B^2 = (-3)^2 + (-4)^2 = 25, \qquad
\Omega C^2 = 4^2 + 1^2 = 17 .
\]
$A$ and $B$ are on the circle; $C$ is inside it ($17 < 25$).
\end{solution}

\begin{solution}{exo:g10:coordgeom:7}
$MA^2 = (x - 1)^2 + 1$ and $MB^2 = (x - 7)^2 + 25$. Setting
$MA^2 = MB^2$:
\[
x^2 - 2x + 2 = x^2 - 14x + 74
\ \Longleftrightarrow\
12x = 72
\ \Longleftrightarrow\
x = 6 .
\]
The unique such point is $M(6, 0)$.
\end{solution}

\begin{solution}{exo:g10:coordgeom:8}
\emph{1.} $I = \left(\frac{0+4}{2}, \frac{3+1}{2}\right) = (2, 2)$.

\emph{2.} $M(2,2)$ is indeed $I$. And
$MA^2 = (0-2)^2 + (3-2)^2 = 5$, $MB^2 = (4-2)^2 + (1-2)^2 = 5$: the
midpoint is equidistant from the two endpoints, as expected.

\emph{3.} Look for $N(0, y)$ with $NA^2 = NB^2$:
\[
(3 - y)^2 = 16 + (1 - y)^2
\ \Longleftrightarrow\
9 - 6y = 16 + 1 - 2y
\ \Longleftrightarrow\
-4y = 8,
\]
so $y = -2$: the point $N(0, -2)$.
\end{solution}

\begin{solution}{exo:g10:coordgeom:9}
\begin{align*}
AB &= \sqrt{4^2 + 2^2} = \sqrt{20} = 2\sqrt5, \\
BC &= \sqrt{8^2 + 4^2} = \sqrt{80} = 4\sqrt5, \\
AC &= \sqrt{12^2 + 6^2} = \sqrt{180} = 6\sqrt5 .
\end{align*}
Then $AB + BC = 2\sqrt5 + 4\sqrt5 = 6\sqrt5 = AC$: the triangle
inequality is an equality, so $B$ lies on the segment $[AC]$ --- the
three points are aligned.
\end{solution}

\begin{solution}{exo:g10:coordgeom:10}
The midpoint of $[AB]$ is $I\left(\frac a2, \frac b2\right)$. Then
\[
OI^2 = \frac{a^2}{4} + \frac{b^2}{4},
\qquad
IA^2 = \left(a - \frac a2\right)^2 + \left(0 - \frac b2\right)^2
= \frac{a^2}{4} + \frac{b^2}{4},
\]
and $IB^2 = \left(0 - \frac a2\right)^2 + \left(b - \frac b2\right)^2 =
\frac{a^2}{4} + \frac{b^2}{4}$ as well. The three squared distances are
equal, so $OI = IA = IB$: the midpoint of the hypotenuse of the right
triangle $OAB$ is equidistant from all three vertices (it is the center
of its circumscribed circle).
\end{solution}

\section*{Chapter \ref{ch:g10:vectors} --- Vectors}

\begin{solution}{exo:g10:vectors:1}
$\vect{AB} + \vect{BD} = \vect{AD}$ (Chasles).

$\vect{KL} + \vect{LM} + \vect{MK} = \vect{KK} = \vec 0$.

$\vect{AC} - \vect{BC} = \vect{AC} + \vect{CB} = \vect{AB}$.

$\vect{AB} + \vect{CA} = \vect{CA} + \vect{AB} = \vect{CB}$.
\end{solution}

\begin{solution}{exo:g10:vectors:2}
$\vect{AB}\,(5-2,\ 3-1) = (3, 2)$; \quad
$\vect{BC}\,(-1-5,\ 4-3) = (-6, 1)$; \quad
$\vect{AC}\,(-1-2,\ 4-1) = (-3, 3)$.
Check: $(3, 2) + (-6, 1) = (-3, 3)$, coordinate by coordinate.
\end{solution}

\begin{solution}{exo:g10:vectors:3}
$\vec u + \vec v = (2 + (-1),\ -3 + 4) = (1, 1)$.

$3\vec u = (6, -9)$.

$2\vec u - \vec v = (4 - (-1),\ -6 - 4) = (5, -10)$.

$\norm{\vec u} = \sqrt{2^2 + (-3)^2} = \sqrt{13}$.
\end{solution}

\begin{solution}{exo:g10:vectors:4}
(a) $4 \times 9 - 6 \times 6 = 36 - 36 = 0$: \omterm{def:g10:vectors:collinear}{collinear}
($\vec v = \frac32 \vec u$).

(b) $2 \times 10 - (-5) \times (-4) = 20 - 20 = 0$: \omterm{def:g10:vectors:collinear}{collinear}
($\vec v = -2\vec u$).

(c) $3 \times 3 - 1 \times 1 = 8 \neq 0$: not \omterm{def:g10:vectors:collinear}{collinear}.
\end{solution}

\begin{solution}{exo:g10:vectors:5}
$\vect{AB}\,(4, -2)$, and $\vect{DC}\,(6 - x_D,\ 3 - y_D)$. The condition
$\vect{AB} = \vect{DC}$ gives $6 - x_D = 4$ and $3 - y_D = -2$, so
$D(2, 5)$.
\end{solution}

\begin{solution}{exo:g10:vectors:6}
\emph{1.} $\vect{AB}\,(3, 2)$ and $\vect{AC}\,(6, -2)$, so
$\vect{AB} + \vect{AC} = (9, 0)$ and
$M = (-2 + 9,\ 1 + 0) = (7, 1)$.

\emph{2.} By the parallelogram rule, $\vect{AM} = \vect{AB} + \vect{AC}$
means that $ABMC$ is a parallelogram ($M$ is the fourth vertex opposite
$A$). One can check: $\vect{AB}\,(3,2)$ and $\vect{CM}\,(7-4,\ 1-(-1)) =
(3, 2)$ are equal.
\end{solution}

\begin{solution}{exo:g10:vectors:7}
First triple: $\vect{AB}\,(3, 2)$, $\vect{AC}\,(9, 6)$;
$3 \times 6 - 2 \times 9 = 0$: aligned ($\vect{AC} = 3\vect{AB}$).

Second triple: $\vect{AB}\,(2, -1)$, $\vect{AC}\,(8, -4)$;
$2 \times (-4) - (-1) \times 8 = -8 + 8 = 0$: aligned as well
($\vect{AC} = 4\vect{AB}$).
\end{solution}

\begin{solution}{exo:g10:vectors:8}
$\vect{AB}\,(4, 2)$ and $\vect{DC}\,(6-2,\ 1-(-1)) = (4, 2)$: the two
\omterm{def:g10:vectors:vector}{vectors} are \emph{equal}, so $ABCD$ is a parallelogram, and in particular
$(AB) \parallel (CD)$. The parallelogram property is the stronger one: it
implies the parallelism (collinearity of the \omterm{def:g10:vectors:vector}{vectors}), not conversely.
\end{solution}

\begin{solution}{exo:g10:vectors:9}
\emph{1.} $\vect{AB}\,(6, 3)$, so
$\vect{AI} = \frac23\vect{AB} = (4, 2)$ and $I(4, 2)$.

\emph{2.} $\vect{CB}\,(6-2,\ 3-5) = (4, -2)$, so
$\vect{CJ} = \frac23\vect{CB} = \left(\frac83, -\frac43\right)$ and
$J\left(2 + \frac83,\ 5 - \frac43\right) = \left(\frac{14}{3},
\frac{11}{3}\right)$.

\emph{3.} $\vect{IJ}\,\left(\frac{14}{3} - 4,\ \frac{11}{3} - 2\right) =
\left(\frac23, \frac53\right)$ and $\vect{AC}\,(2, 5)$. Collinearity
criterion:
\[
\frac23 \times 5 - \frac53 \times 2 = \frac{10}{3} - \frac{10}{3} = 0 :
\]
the \omterm{def:g10:vectors:vector}{vectors} are \omterm{def:g10:vectors:collinear}{collinear}, so $(IJ)$ is parallel to $(AC)$. (Both $I$ and
$J$ sit one third of the way from the side towards $B$: the small
triangle $IBJ$ is a reduction of $ABC$.)
\end{solution}

\begin{solution}{exo:g10:vectors:10}
The \omterm{prop:g10:coordgeom:midpoint}{midpoints} are
\[
I\left(\tfrac{x_A + x_B}{2}, \tfrac{y_A + y_B}{2}\right), \quad
J\left(\tfrac{x_B + x_C}{2}, \tfrac{y_B + y_C}{2}\right), \quad
K\left(\tfrac{x_C + x_D}{2}, \tfrac{y_C + y_D}{2}\right), \quad
L\left(\tfrac{x_D + x_A}{2}, \tfrac{y_D + y_A}{2}\right).
\]
First \omterm{def:g10:coordgeom:system}{coordinates} of the two \omterm{def:g10:vectors:vector}{vectors}:
\[
x_J - x_I = \frac{x_B + x_C}{2} - \frac{x_A + x_B}{2}
= \frac{x_C - x_A}{2},
\qquad
x_K - x_L = \frac{x_C + x_D}{2} - \frac{x_D + x_A}{2}
= \frac{x_C - x_A}{2},
\]
and the same computation works for the second \omterm{def:g10:coordgeom:system}{coordinates}. So
$\vect{IJ} = \vect{LK}$ (both equal $\frac12\vect{AC}$), and $IJKL$ is a
parallelogram, whatever the quadrilateral $ABCD$ looks like.
\end{solution}

\begin{solution}{pb:g10:vectors:1}
\textbf{1.} $\vect{AB} + \vect{BC} + \vect{CD} = \vect{AD}$;
\quad $\vect{MA} - \vect{MB} = \vect{MA} + \vect{BM} =
\vect{BA}$; \quad $\vect{AB} + \vect{BC} + \vect{CD} +
\vect{DA} = \vect{AA} = \vec 0$ (reordering the third \omterm{def:g10:vectors:sum}{sum}).

\textbf{2.} $\vect{OB} - \vect{OA} = \vect{AO} + \vect{OB} =
\vect{AB}$ by Chasles.

\textbf{3.} $I$ is the \omterm{prop:g10:coordgeom:midpoint}{midpoint} of $[AB]$ exactly when
$\vect{IA} = -\vect{IB}$, i.e.\
$\vect{IA} + \vect{IB} = \vec 0$. Then, for any $O$:
$\vec 0 = \vect{IO} + \vect{OA} + \vect{IO} + \vect{OB}$, so
$2\,\vect{OI} = \vect{OA} + \vect{OB}$ --- and conversely.

\textbf{4.} $\vect{AB} = (4, 2)$, and $D$ satisfies
$\vect{DC} = \vect{AB}$: $D = C - (4, 2) = (2, -2)$ --- the
same $D$ as \cref{pb:g10:coordgeom:1}, one line faster.

\textbf{5.} The translation of \omterm{def:g10:vectors:vector}{vector} $\vec u$ --- exactly the
transformation produced by two parallel mirrors or two
half-turns. \omterm{def:g10:functions:function}{Image} of $(1, 2)$: $(4, 1)$.

\textbf{6.} \omterm{def:g10:algebra:expand}{Expanding} each $\vect{GX}$ as
$\vect{GO} + \vect{OX}$:
$3\,\vect{GO} + \vect{OA} + \vect{OB} + \vect{OC} = \vec 0$,
so $\vect{OG} = \frac13(\vect{OA} + \vect{OB} + \vect{OC})$:
one and only one such point, whose \omterm{def:g10:coordgeom:system}{coordinates} are the averages
of the three vertices'.

\textbf{7.} $\vect{AB} + \vect{AC} =
(\vect{AG} + \vect{GB}) + (\vect{AG} + \vect{GC})
= 2\vect{AG} + (\vect{GB} + \vect{GC})
= 2\vect{AG} - \vect{GA} = 3\vect{AG}$, whence
$\vect{AG} = \frac13(\vect{AB} + \vect{AC})$. The \omterm{prop:g10:coordgeom:midpoint}{midpoint}
formula (question 3, seen from $A$) gives
$\vect{AK} = \frac12(\vect{AB} + \vect{AC})$. Comparing:
$\vect{AG} = \frac23\vect{AK}$ --- $G$ lies on the median
$[AK]$, two thirds of the way from $A$.

\textbf{8.} The defining condition treats $A$, $B$, $C$
identically, so the same computation from $B$ or from $C$
places the \emph{same} $G$ at two thirds of those medians:
concurrency costs nothing. The parallelogram proof was
ingenious but ad hoc; \omterm{def:g10:coordgeom:system}{coordinates} compute but obscure;
\omterm{def:g10:vectors:vector}{vectors} make the symmetry visible and the proof two lines.

\textbf{9.} $G\left(\frac{1+5+3}{3}, \frac{1+2+6}{3}\right) =
(3, 3)$; $K(4, 4)$; $\vect{AG} = (2, 2)$ and
$\vect{AK} = (3, 3)$: determinant
$2 \times 3 - 2 \times 3 = 0$, \omterm{def:g10:vectors:collinear}{collinear}, and indeed
$\vect{AG} = \frac23\vect{AK}$.

\textbf{10.} $G'\left(\frac{2 + 5 + 3}{4}, \frac{2 + 2 +
6}{4}\right) = (2.5,\ 2.5)$: pulled from $G(3,3)$ towards the
doubled vertex $A(1,1)$, as a balance point should be.

\textbf{11.} $3 \times 4 - (-2)(-6) = 12 - 12 = 0$: \omterm{def:g10:vectors:collinear}{collinear}
($\vec v = -2\vec u$). $2 \times 9 - 5 \times 4 = -2 \neq 0$:
not \omterm{def:g10:vectors:collinear}{collinear}.

\textbf{12.} $\vect{AB} = (2, 3)$, $\vect{AC} = (6, 9)$:
determinant $2 \times 9 - 3 \times 6 = 0$: aligned.

\textbf{13.} $\vect{OI} = \frac12(\vect{OA} + \vect{OB})$ and
$\vect{OK} = \frac12(\vect{OC} + \vect{OD})$, so the \omterm{prop:g10:coordgeom:midpoint}{midpoint}
of $[IK]$ has position \omterm{def:g10:vectors:vector}{vector}
$\frac12(\vect{OI} + \vect{OK}) = \frac14(\vect{OA} +
\vect{OB} + \vect{OC} + \vect{OD})$ --- and the computation
for $[JL]$ gives the very same \omterm{def:g10:vectors:vector}{vector}. The two bimedians (and,
by Varignon, the parallelogram $IJKL$'s diagonals) all share
this center: the quadrilateral's centroid.

\textbf{14.} $\vect{MN} = \vect{AN} - \vect{AM} =
\frac13\vect{AC} - \frac13\vect{AB} = \frac13\vect{BC}$:
so $(MN) \parallel (BC)$ and $MN = \frac13 BC$.

\textbf{15.} With $\vect{AM} = k\,\vect{AB}$ and
$\vect{AN} = k\,\vect{AC}$: $\vect{MN} = k\,\vect{BC}$ --- one
line containing the \omterm{prop:g10:coordgeom:midpoint}{midpoint} theorem ($k = \frac12$) and
Thales (any $k$): the ratio passes through the subtraction
untouched.

\textbf{16.} Resultant $(3, 4)$, magnitude
$\sqrt{9 + 16} = 5$ km/h: the boat actually travels
north-east-by-north along the $3$--$4$--$5$ triangle's
hypotenuse.

\textbf{17.} $\vec u + \vec v + \vec w = (0, 0)$: the ring
stays put --- equilibrium. Along a closed polygon
$A_1 A_2 \dots A_n A_1$, Chasles telescopes the \omterm{def:g10:vectors:sum}{sum} to
$\vect{A_1 A_1} = \vec 0$: going around brings you back.

\textbf{18.} $\vec w = -(\vec u + \vec v) = -(3, 2) =
(-3, -2)$.

\textbf{19.} Resultant $(1.5,\ 2)$, magnitude
$\sqrt{2.25 + 4} = 2.5$ m/s. Crossing: $50$ m at $2$ m/s
across: $25$ s; drift: $1.5 \times 25 = 37.5$ m downstream.

\textbf{20.} Re-proved this weekend: the centroid's two-thirds
theorem (question 7), the bimedian/Varignon center (question
13), Thales--midpoint (questions 14--15). What arrows still
cannot do: measure \emph{angles} (and hence lengths in
arbitrary directions) --- the scalar product of grade 11 is
precisely the missing instrument.
\end{solution}

\section*{Chapter \ref{ch:g10:lines} --- Equations of Lines and Linear Systems}

\begin{solution}{exo:g10:lines:1}
$y = 2x - 3$: slope $2$, intercept $-3$.
$y = -\frac13 x + 1$: slope $-\frac13$, intercept $1$.
$y = 4$: slope $0$, intercept $4$ (horizontal line).
$x = -2$: vertical line, no slope and no reduced equation.
\end{solution}

\begin{solution}{exo:g10:lines:2}
$3 \times 2 + 1 = 7$: yes, $A$ is on the line.
$3 \times (-1) + 1 = -2 \neq -1$: $B$ is not.
$3 \times 4 + 1 = 13$: yes, $C$ is on the line.
\end{solution}

\begin{solution}{exo:g10:lines:3}
(a) $m = \frac{8 - 2}{3 - 0} = 2$ and $p = y_A = 2$ (the point $A$ is on
the $y$-axis): $y = 2x + 2$.

(b) $m = \frac{-4 - 5}{2 - (-1)} = \frac{-9}{3} = -3$; then
$5 = -3 \times (-1) + p$ gives $p = 2$: $y = -3x + 2$.

(c) $x_A = x_B = 4$: vertical line $x = 4$.
\end{solution}

\begin{solution}{exo:g10:lines:4}
$y = \frac{6x+2}{2} = 3x + 1$. The lines with slope $3$ are $y = 3x - 1$,
$y = 3x + 5$ and $y = 3x + 1$: those three are parallel to each other;
no two of them are equal (different intercepts), and none is parallel to
$y = -3x + 2$.
\end{solution}

\begin{solution}{exo:g10:lines:5}
First system: substitute $y = 2x - 1$ into $3x + y = 9$:
$3x + 2x - 1 = 9$, so $5x = 10$, $x = 2$, then $y = 3$. Solution:
$(2, 3)$.

Second system: from $x - y = 4$, $x = y + 4$. Then
$2(y + 4) + 3y = 3$, so $5y = -5$, $y = -1$, then $x = 3$. Solution:
$(3, -1)$.
\end{solution}

\begin{solution}{exo:g10:lines:6}
First system: adding the equations eliminates $y$: $2x = 14$, $x = 7$;
then $7 + y = 10$ gives $y = 3$. Solution: $(7, 3)$.

Second system: multiply the first equation by $5$ and the second by $-2$:
\[
\begin{cases}
15x + 10y = 60 \\
-4x - 10y = -16
\end{cases}
\]
Adding: $11x = 44$, so $x = 4$; then $3 \times 4 + 2y = 12$ gives
$y = 0$. Solution: $(4, 0)$.
\end{solution}

\begin{solution}{exo:g10:lines:7}
Compute $ab' - a'b$ in each case.

First: $2 \times 2 - (-4) \times (-1) = 4 - 4 = 0$, and the second
equation is $-2$ times the first: twice the same line, infinitely many
solutions.

Second: $ab' - a'b = 0$ again, but the right-hand sides are not in the
ratio $-2$ ($1 \neq -6$): two distinct parallel lines, no solution.

Third: $2 \times 1 - 1 \times (-1) = 3 \neq 0$: exactly one solution.
\end{solution}

\begin{solution}{exo:g10:lines:8}
Let $x$ be the number of child tickets and $y$ the number of adult
tickets:
\[
\begin{cases}
x + y = 200 \\
5x + 9y = 1352 .
\end{cases}
\]
From the first equation $x = 200 - y$; substituting,
$5(200 - y) + 9y = 1352$, so $1000 + 4y = 1352$, $y = 88$, then
$x = 112$. So $112$ child tickets and $88$ adult tickets. Check:
$5 \times 112 + 9 \times 88 = 560 + 792 = 1352$.
\end{solution}

\begin{solution}{exo:g10:lines:9}
\emph{1.} Parallel means same slope $2$: $y = 2x + p$ with
$4 = 2 \times 1 + p$, so $p = 2$: $d'\colon y = 2x + 2$.

\emph{2.} On the $x$-axis, $y = 0$: $2x + 2 = 0$ gives $x = -1$. The
intersection point is $(-1, 0)$.
\end{solution}

\begin{solution}{exo:g10:lines:10}
\emph{1.} Midpoint of $[BC]$: $A' = \left(\frac{6+2}{2},
\frac{2+4}{2}\right) = (4, 3)$. The median from $A(0,0)$ has slope
$\frac{3 - 0}{4 - 0} = \frac34$: equation $y = \frac34 x$.

Midpoint of $[AC]$: $B' = (1, 2)$. The median from $B(6, 2)$ has slope
$\frac{2 - 2}{1 - 6} = 0$: it is the horizontal line $y = 2$.

\emph{2.} Intersection: $\frac34 x = 2$ gives $x = \frac83$, so
$G\left(\frac83, 2\right)$. The third median joins $C(2,4)$ to the
midpoint of $[AB]$, $C' = (3, 1)$; its slope is
$\frac{1 - 4}{3 - 2} = -3$, equation $y = -3x + 10$. At $x = \frac83$:
$y = -8 + 10 = 2$: $G$ is on it. And indeed
$G = \left(\frac{0 + 6 + 2}{3}, \frac{0 + 2 + 4}{3}\right)$: the
centroid averages the coordinates of the vertices.
\end{solution}

\section*{Chapter \ref{ch:g10:stats} --- Descriptive Statistics}

\begin{solution}{exo:g10:stats:1}
Sum: $0+1+1+2+2+2+3+3+4+4+5+6+7+8+12 = 60$, so the \omterm{def:g10:stats:mean}{mean} is
$\frac{60}{15} = 4$. The series is already sorted and $N = 15$ is odd:
the \omterm{def:g10:stats:median}{median} is the $8$th value, $3$. \omterm{def:g10:stats:quartiles}{Range}: $12 - 0 = 12$.
\end{solution}

\begin{solution}{exo:g10:stats:2}
Frequencies: $\frac{7}{50} = 0.14$, $\frac{9}{50} = 0.18$,
$\frac{8}{50} = 0.16$, $\frac{10}{50} = 0.20$, $\frac{6}{50} = 0.12$,
$\frac{10}{50} = 0.20$ (sum $1$). \omterm{def:g10:stats:mean}{Mean} outcome:
\[
\bar x = \frac{7 \times 1 + 9 \times 2 + 8 \times 3 + 10 \times 4
+ 6 \times 5 + 10 \times 6}{50}
= \frac{7 + 18 + 24 + 40 + 30 + 60}{50}
= \frac{179}{50} = 3.58 .
\]
\end{solution}

\begin{solution}{exo:g10:stats:3}
$N = 12$ (even): the \omterm{def:g10:stats:median}{median} is the \omterm{prop:g10:coordgeom:midpoint}{midpoint} of the $6$th and $7$th
values, both equal to $6$: \omterm{def:g10:stats:median}{median} $6$.

$Q_1$: $\frac{12}{4} = 3$, so $Q_1$ is the $3$rd value: $4$.

$Q_3$: $\frac{3 \times 12}{4} = 9$, so $Q_3$ is the $9$th value: $8$.

Box plot: box from $4$ to $8$ with the \omterm{def:g10:stats:median}{median} bar at $6$, whiskers from
$1$ to $11$.
\end{solution}

\begin{solution}{exo:g10:stats:4}
The $24$ scores total $24 \times 11.5 = 276$. With the new score the
total is $276 + 19 = 295$ for $25$ students: new \omterm{def:g10:stats:mean}{mean}
$\frac{295}{25} = 11.8$.
\end{solution}

\begin{solution}{exo:g10:stats:5}
Current total: $5 \times 12 = 60$. A \omterm{def:g10:stats:mean}{mean} of $13$ over $6$ tests requires
a total of $78$, so the $6$th score must be $78 - 60 = 18$. A \omterm{def:g10:stats:mean}{mean} of
$14$ would require $84 - 60 = 24 > 20$: not reachable.
\end{solution}

\begin{solution}{exo:g10:stats:6}
\omterm{def:g10:stats:mean}{Mean} larger than $20$: for instance
$1,\ 2,\ 3,\ 10,\ 40,\ 50,\ 60$ --- \omterm{def:g10:stats:median}{median} $10$ (the $4$th sorted value),
\omterm{def:g10:stats:mean}{mean} $\frac{166}{7} \approx 23.7$.

\omterm{def:g10:stats:mean}{Mean} smaller than $6$: for instance
$0,\ 0,\ 0,\ 10,\ 10,\ 10,\ 10$ --- \omterm{def:g10:stats:median}{median} $10$, \omterm{def:g10:stats:mean}{mean}
$\frac{40}{7} \approx 5.7$.

The \omterm{def:g10:stats:mean}{mean} cannot go below $\frac{40}{7}$: for the \omterm{def:g10:stats:median}{median} (the $4$th sorted
value) to be $10$, the values in positions $4$ to $7$ must all be
$\geq 10$, so the total is at least $40$ and the \omterm{def:g10:stats:mean}{mean} at least
$\frac{40}{7}$ --- the second example is extremal.

These examples show that the \omterm{def:g10:stats:mean}{mean} and the \omterm{def:g10:stats:median}{median} are largely independent:
knowing one says little about the other, which is why a good summary
reports both.
\end{solution}

\begin{solution}{exo:g10:stats:7}
\emph{1.} \omterm{def:g10:stats:mean}{Mean}:
$\frac{9 \times 1800 + 12000}{10} = \frac{16200 + 12000}{10} = 2820$.
\omterm{def:g10:stats:median}{Median}: sorted, the $10$ salaries are nine values $1800$ then $12000$;
the \omterm{def:g10:stats:median}{median} is the \omterm{prop:g10:coordgeom:midpoint}{midpoint} of the $5$th and $6$th values, both $1800$:
\omterm{def:g10:stats:median}{median} $1800$.

\emph{2.} The \omterm{def:g10:stats:median}{median} ($1800$) describes the typical salary: it is what
$9$ employees out of $10$ actually earn. The \omterm{def:g10:stats:mean}{mean} ($2820$) is pulled up
by the single high salary and matches nobody's payslip.
\end{solution}

\begin{solution}{exo:g10:stats:8}
Total height: $12 \times 168 + 18 \times 163 = 2016 + 2934 = 4950$ cm for
$30$ students, so the \omterm{def:g10:stats:mean}{mean} is $\frac{4950}{30} = 165$ cm. It is closer to
the girls' \omterm{def:g10:stats:mean}{mean} because they are more numerous (weighted \omterm{def:g10:stats:mean}{mean}, weights
$12$ and $18$).
\end{solution}

\begin{solution}{exo:g10:stats:9}
\emph{1.} The new \omterm{def:g10:stats:mean}{mean} is
\[
\bar y = \frac{y_1 + \dots + y_N}{N}
= \frac{(a x_1 + b) + \dots + (a x_N + b)}{N}
= \frac{a(x_1 + \dots + x_N) + Nb}{N}
= a \bar x + b .
\]

\emph{2.} With $a = 1.8$ and $b = 32$: \omterm{def:g10:stats:mean}{mean}
$1.8 \times 20 + 32 = 68$ degrees Fahrenheit.
\end{solution}

\begin{solution}{pb:g10:stats:1}
\textbf{1.} \omterm{def:g10:stats:mean}{Mean}: $\frac{2\,550}{7} \approx 364$ thousand
euros; \omterm{def:g10:stats:median}{median}: the fourth of the seven ordered prices, $230$.
The \omterm{def:g10:stats:median}{median} describes the street --- six of the seven houses
cost nowhere near $364$.

\textbf{2.} Without the mansion: \omterm{def:g10:stats:mean}{mean} $\frac{1\,350}{6} = 225$,
\omterm{def:g10:stats:median}{median} $\frac{220 + 230}{2} = 225$. The \omterm{def:g10:stats:mean}{mean} fell by about
$139$; the \omterm{def:g10:stats:median}{median} by $5$. One extreme value can drag the \omterm{def:g10:stats:mean}{mean}
anywhere; the \omterm{def:g10:stats:median}{median} barely notices --- it is \emph{robust}.

\textbf{3.} \omterm{def:g10:stats:mean}{Mean} $= \frac{0 \times 20 + 1 \times 25 + 2 \times
30 + 3 \times 15 + 4 \times 10}{100} = \frac{170}{100} = 1.7$.
No family has $1.7$ children, but $1.7 \times 100 = 170$
\emph{is} the exact total number of children: the \omterm{def:g10:stats:mean}{mean} encodes
the total, redistributed equally --- its true talent.

\textbf{4.} $N = 15$: \omterm{def:g10:stats:median}{median} $=$ 8th value $= 11$;
$Q_1 =$ 4th value $= 8$; $Q_3 =$ 12th value $= 14$.

\textbf{5.} \omterm{def:g10:stats:quartiles}{Range} $19 - 4 = 15$: the full spread, extremes
included. \omterm{def:g10:stats:quartiles}{Interquartile range} $14 - 8 = 6$: the width of the
middle half --- spread without the extremes' noise.

\textbf{6.} A shift $+b$ moves \emph{every} value, so
differences of values --- \omterm{def:g10:stats:quartiles}{range}, IQR --- are unchanged:
$(x_i + b) - (x_j + b) = x_i - x_j$. A scaling by $a$
multiplies all differences by $a$: \omterm{def:g10:stats:quartiles}{range} and IQR are scaled by
$\abs a$.

\textbf{7.} Additive curve: $6 \to 9$ and $18 \to 21$ --- the
weak student gains $50\,\%$, the strong one $17\,\%$, and $21$
bursts the $20$-point scale only mildly. Multiplicative:
$6 \to 8$ and $18 \to 24$ --- ratios preserved, but $24$ is
impossible on a $20$-scale: the ceiling problem. Additive
curves flatter the bottom; multiplicative ones explode the
top.

\textbf{8.} Total marks: $20 \times 12 + 30 \times 10 = 540$,
so the merged \omterm{def:g10:stats:mean}{mean} is $\frac{540}{50} = 10.8$: class B's thirty
students outweigh class A's twenty --- means merge by
\emph{weighted} average, never by simple average.

\textbf{9.} \omterm{def:g10:stats:median}{Median} of $A$: $2$; of $B$: $5$. Merged list
$\{0, 1, 2, 3, 10\}$: \omterm{def:g10:stats:median}{median} $2$. The size-weighted average of
the \omterm{def:g10:stats:median}{medians} would be $\frac{3 \times 2 + 2 \times 5}{5} = 3.2
\neq 2$: \omterm{def:g10:stats:median}{medians} carry no total, so no merging formula exists
--- one must re-sort the whole list.

\textbf{10.} \omterm{def:g10:stats:mean}{Mean}: $\frac{31.3}{5} = 6.26$. \omterm{def:g10:stats:median}{Median}: $5.4$.
Trimmed \omterm{def:g10:stats:mean}{mean}: $\frac{5.3 + 5.4 + 5.5}{3} = 5.4$. One
enthusiastic (or corrupt) judge moved the \omterm{def:g10:stats:mean}{mean} by almost a
point; the \omterm{def:g10:stats:median}{median} and the trimmed \omterm{def:g10:stats:mean}{mean} shrugged --- hence the
trimming rules of judged sports.

\textbf{11.} Factory A: middle half within $0.2$ mm; factory B:
within $4$ mm. Same \omterm{def:g10:stats:mean}{mean}, opposite reliability: buy from A. The
\omterm{def:g10:stats:quartiles}{interquartile range} decided --- means alone compare centers,
never consistency.

\textbf{12.} \omterm{def:g10:stats:median}{Median} far below \omterm{def:g10:stats:mean}{mean} betrays a \emph{long right
tail}: most waits short, a few catastrophic. Example:
$\{5, 5, 10, 10, 95\}$ --- \omterm{def:g10:stats:median}{median} $10$, \omterm{def:g10:stats:mean}{mean}
$\frac{125}{5} = 25$.

\textbf{13.} It means $75\,\%$ of babies of that age are
shorter (and $25\,\%$ taller): a rank, not a size. ``Every
child above the \omterm{def:g10:stats:mean}{mean}'' is impossible --- if all values exceeded
the \omterm{def:g10:stats:mean}{mean}, their average would exceed the \omterm{def:g10:stats:mean}{mean}, which is itself:
contradiction. But \emph{almost} all can be: one very short
child can hold the \omterm{def:g10:stats:mean}{mean} below everyone else. Above the
\emph{\omterm{def:g10:stats:median}{median}}: never more than half, by definition. Ranks
cannot be flattered; averages can.

\textbf{14.} Typical mark: same, \omterm{def:g10:stats:median}{median} $11$ for both.
Homogeneity of the middle: B tighter (IQR $3$ against $5$).
Extremes: B wilder (\omterm{def:g10:stats:quartiles}{range} $17$ against $13$). B is a class of
similar students plus a few outliers; A is evenly spread.

\textbf{15.} (1) Is $5\,000$ the \omterm{def:g10:stats:mean}{mean} or the \omterm{def:g10:stats:median}{median} --- and may
I have the \omterm{def:g10:stats:median}{median}? (2) What is the spread (\omterm{def:g10:stats:quartiles}{quartiles}), so I
know how typical ``typical'' is? (3) Who was counted --- sample
size and selection (the how-data-lies weekend problem of the Middle School volume's poll lesson)?

\textbf{16.} For instance $100, 150, 200, 250, 800$: \omterm{def:g10:stats:median}{median}
$200$, \omterm{def:g10:stats:mean}{mean} $\frac{1\,500}{5} = 300$.

\textbf{17.} $\sum (x_i - \bar x) = \sum x_i - N \bar x =
N\bar x - N\bar x = 0$. Check: $-200 - 150 - 100 - 50 + 500 =
0$. The \omterm{def:g10:stats:mean}{mean} is the point where the deviations balance --- the
data's center of gravity (the centroid weekend problem of the Middle School volume's $G$, in
one dimension).

\textbf{18.} At the \omterm{def:g10:stats:median}{median} $4$: $3 + 1 + 0 + 5 + 9 = 18$ km. At
the \omterm{def:g10:stats:mean}{mean} $6$: $5 + 3 + 2 + 3 + 7 = 20$ km. At $5$:
$4 + 2 + 1 + 4 + 8 = 19$ km. The \omterm{def:g10:stats:median}{median} wins --- moving away
from it always passes more shops than it approaches, so the
total can only grow.

\textbf{19.} $0.3^{10} \approx 6 \times 10^{-6}$: about six
pilots \emph{per million}. In $4\,000$ pilots the expected
count is $0.02$ --- zero, as lieutenant Daniels found. The Air
Force stopped designing for the average man and made
everything adjustable: seats, pedals, straps --- design for the
\emph{spread}, not the center.

\textbf{20.} Totals: the \emph{\omterm{def:g10:stats:mean}{mean}}. Typical individual: the
\emph{\omterm{def:g10:stats:median}{median}}. Fairness of spread: the \emph{\omterm{def:g10:stats:quartiles}{quartiles}} (and
IQR). Contamination: the \emph{trimmed \omterm{def:g10:stats:mean}{mean}}. And the closing
line: the average person does not exist --- but the \omterm{def:g10:stats:median}{median}
person almost does.
\end{solution}

\section*{Chapter \ref{ch:g10:proba} --- Probability and Sampling}

\begin{solution}{exo:g10:proba:1}
Each of the $6$ outcomes has \omterm{def:g10:proba:distribution}{probability} $\frac16$.
``Getting a $6$'': $\frac16$.
``Odd number'' $= \{1, 3, 5\}$: $\frac36 = \frac12$.
``At most $4$'' $= \{1,2,3,4\}$: $\frac46 = \frac23$.
``Getting a $7$'': impossible \omterm{def:g10:proba:events}{event}, \omterm{def:g10:proba:distribution}{probability} $0$.
\end{solution}

\begin{solution}{exo:g10:proba:2}
There are $10$ equally likely marbles.
$\P(\text{green}) = \frac{5}{10} = \frac12$.
$\P(\text{not black}) = 1 - \frac{2}{10} = \frac{8}{10} = \frac45$.
Green and yellow are \omterm{def:g10:proba:operations}{incompatible}:
$\P(\text{green or yellow}) = \frac{5}{10} + \frac{3}{10} = \frac{8}{10}
= \frac45$.
\end{solution}

\begin{solution}{exo:g10:proba:3}
Let $S$: ``studies Spanish'', $G$: ``studies German''. Then
$\P(S) = \frac{18}{30}$, $\P(G) = \frac{10}{30}$,
$\P(S \cap G) = \frac{4}{30}$. Addition rule:
\[
\P(S \cup G) = \frac{18}{30} + \frac{10}{30} - \frac{4}{30}
= \frac{24}{30} = \frac45 .
\]
\omterm{def:g10:proba:operations}{Complement}: $\P(\text{neither}) = 1 - \frac45 = \frac15$.
\end{solution}

\begin{solution}{exo:g10:proba:4}
The \omterm{met:g10:proba:tree}{tree} has four paths HH, HT, TH, TT, each of \omterm{def:g10:proba:distribution}{probability}
$\frac12 \times \frac12 = \frac14$.
Two heads: $\frac14$. Exactly one head (HT or TH): $\frac24 = \frac12$.
At least one head: $1 - \P(\text{TT}) = 1 - \frac14 = \frac34$.
\end{solution}

\begin{solution}{exo:g10:proba:5}
First draw: red with \omterm{def:g10:proba:distribution}{probability} $\frac{4}{10}$, white $\frac{6}{10}$.
Second draw \emph{without} replacement: after a red, $3$ red and $6$
white remain among $9$; after a white, $4$ red and $5$ white among $9$.

Two reds: $\frac{4}{10} \times \frac39 = \frac{12}{90} = \frac{2}{15}$.

Different colors: paths RW and WR:
\[
\frac{4}{10} \times \frac69 + \frac{6}{10} \times \frac49
= \frac{24}{90} + \frac{24}{90} = \frac{48}{90} = \frac{8}{15}.
\]
\end{solution}

\begin{solution}{exo:g10:proba:6}
\emph{1.} Each die independently shows each face with \omterm{def:g10:proba:distribution}{probability}
$\frac16$, so the $36$ ordered pairs are equally likely; the sum is
determined by the pair.

\emph{2.} Sum $7$: the pairs $(1,6), (2,5), (3,4), (4,3), (5,2), (6,1)$,
so $\P = \frac{6}{36} = \frac16$. Sum $12$: only $(6,6)$,
$\P = \frac{1}{36}$. Sum at most $4$: pairs summing to $2$, $3$ or $4$:
$(1,1)$; $(1,2), (2,1)$; $(1,3), (2,2), (3,1)$ --- six pairs, so
$\P = \frac{6}{36} = \frac16$.

\emph{3.} The number of pairs achieving each sum grows from $1$ (sum $2$)
to $6$ (sum $7$) then decreases: the most likely sum is $7$.
\end{solution}

\begin{solution}{exo:g10:proba:7}
Each question is answered correctly with \omterm{def:g10:proba:distribution}{probability} $\frac14$,
independently. \omterm{met:g10:proba:tree}{Tree} with two levels:

Both correct: $\frac14 \times \frac14 = \frac{1}{16}$.

Exactly one correct: paths correct--wrong and wrong--correct:
$\frac14 \times \frac34 + \frac34 \times \frac14 = \frac{6}{16}
= \frac38$.

None correct: $\frac34 \times \frac34 = \frac{9}{16}$.

Check: $\frac{1}{16} + \frac{6}{16} + \frac{9}{16} = 1$.
\end{solution}

\begin{solution}{exo:g10:proba:8}
\emph{1.} Observed \omterm{def:g10:stats:series}{frequency}: $\frac{396}{900} = 0.44$. With $n = 900$:
$\frac{1}{\sqrt{900}} = \frac{1}{30} \approx 0.033$, so the fluctuation
\omterm{def:g10:numbers:interval}{interval} is about $\intcc{0.467}{0.533}$.

\emph{2.} The observed $0.44$ lies outside the \omterm{def:g10:numbers:interval}{interval}: \omterm{def:g10:proba:sample}{samples} of $900$
voters from a population with $50\%$ approval rarely stray that far. The
poll casts serious doubt on the claimed $50\%$.
\end{solution}

\begin{solution}{exo:g10:proba:9}
\emph{1.} $\P(\text{win } 6) = \frac16$ (roll a $6$),
$\P(\text{win } 3) = \frac16$ (roll a $5$),
$\P(\text{win } 0) = \frac46 = \frac23$.

\emph{2.} Over $600$ games, expect about $100$ sixes, $100$ fives, $400$
other rolls. Winnings: about $100 \times 6 + 100 \times 3 = 900$. Cost:
$600 \times 2 = 1200$. Expected net result: $900 - 1200 = -300$, i.e.\ a
loss of about $0.5$ per game on average --- the game is not worth
playing.
\end{solution}


% Solutions — shared structure; bodies under solutions/ or solutions/<lang>/
\chapter{\omstr{part.grade11}}

\ominputsol{grade-11}{01-quadratic-functions}
\ominputsol{grade-11}{02-functions-variations}
\ominputsol{grade-11}{03-differentiation}
\ominputsol{grade-11}{04-sequences}
\ominputsol{grade-11}{05-trigonometry}
\ominputsol{grade-11}{06-vectors-lines}
\ominputsol{grade-11}{07-scalar-product}
\ominputsol{grade-11}{08-descriptive-statistics}
\ominputsol{grade-11}{09-probability-random-variables}
\ominputsol{grade-11}{10-binomial-distribution}

% Solutions to the Grade 12 exercises — one file per chapter.
\chapter{Grade 12 (Age 17--18)}

\section*{Chapter \ref{ch:g12:seq} --- Sequences}

\begin{solution}{exo:g12:seq:1}
Let $P(n)$ be the statement $\sum_{k=1}^{n} k^2 = \frac{n(n+1)(2n+1)}{6}$.
\emph{Base case:} for $n = 0$ both sides are $0$ (empty sum).
\emph{Inductive step:} assume $P(n)$. Then
\[
\sum_{k=1}^{n+1} k^2 = \frac{n(n+1)(2n+1)}{6} + (n+1)^2
= \frac{(n+1)\bigl(n(2n+1) + 6(n+1)\bigr)}{6}
= \frac{(n+1)(2n^2 + 7n + 6)}{6}.
\]
Since $2n^2 + 7n + 6 = (n+2)(2n+3)$, this is
$\frac{(n+1)(n+2)(2(n+1)+1)}{6}$, which is $P(n+1)$. By induction, $P(n)$
holds for all $n$.
\end{solution}

\begin{solution}{exo:g12:seq:2}
$a_{n+1} - a_n = \frac{n+2}{n+1} - \frac{n+1}{n}
= \frac{n(n+2) - (n+1)^2}{n(n+1)} = \frac{-1}{n(n+1)} < 0$: $(a_n)$ is
strictly decreasing.

$(b_n)$ has positive terms and
$\frac{b_{n+1}}{b_n} = \frac{2^{n+1}}{n+1}\cdot\frac{n}{2^n} =
\frac{2n}{n+1} \geq 1 \iff 2n \geq n+1 \iff n \geq 1$: $(b_n)$ is increasing
(strictly for $n \geq 2$).

$c_{n+1} - c_n = (n+1)^2 - 10(n+1) - n^2 + 10n = 2n - 9$, which is negative
for $n \leq 4$ and positive for $n \geq 5$: $(c_n)$ decreases up to
$c_5 = -25$, its minimum, then increases. It is not monotonic.
\end{solution}

\begin{solution}{exo:g12:seq:3}
Factor dominant terms:
\[
u_n = \frac{n^2(3 - 1/n + 1/n^2)}{n^2(2 + 5/n^2)} \xrightarrow[n\to+\infty]{} \frac{3}{2}.
\]
Multiply by the conjugate:
\[
v_n = \frac{(n+1) - n}{\sqrt{n+1} + \sqrt{n}} = \frac{1}{\sqrt{n+1}+\sqrt{n}}
\xrightarrow[n\to+\infty]{} 0.
\]
Divide numerator and denominator by $3^n$:
\[
w_n = \frac{(2/3)^n - 1}{1 + (1/3)^n} \xrightarrow[n\to+\infty]{} \frac{0-1}{1+0} = -1,
\]
using $\lim q^n = 0$ for $\abs{q} < 1$.
\end{solution}

\begin{solution}{exo:g12:seq:4}
$u_n = 5 + 3n \to +\infty$ and $v_n = 8 \cdot (1/2)^n = 2^{3-n} \to 0$.
The sums are
\[
\sum_{k=0}^{n} u_k = (n+1)\,\frac{5 + (5+3n)}{2} = \frac{(n+1)(10+3n)}{2}
\xrightarrow[n\to+\infty]{} +\infty,
\]
\[
\sum_{k=0}^{n} v_k = 8\,\frac{1 - (1/2)^{n+1}}{1 - 1/2}
= 16\left(1 - \left(\tfrac{1}{2}\right)^{n+1}\right)
\xrightarrow[n\to+\infty]{} 16 .
\]
\end{solution}

\begin{solution}{exo:g12:seq:5}
Since $-1 \leq \cos n \leq 1$,
\[
\frac{n-1}{n+1} \leq \frac{n + \cos n}{n+1} \leq 1,
\]
and $\frac{n-1}{n+1} \to 1$, so the limit is $1$ by the squeeze theorem.

For the second limit, write
\[
0 \leq \frac{n!}{n^n}
= \frac{1}{n}\cdot\frac{2}{n}\cdots\frac{n}{n}
\leq \frac{1}{n},
\]
because every factor $\frac{k}{n}$ with $2 \leq k \leq n$ is at most $1$.
Since $\frac1n \to 0$, the squeeze theorem gives $\frac{n!}{n^n} \to 0$. (The
suggested geometric bound also works: each factor with $k \leq n/2$ is at most
$\frac12$, giving the stronger bound $(1/2)^{\floor{n/2}}$.)
\end{solution}

\begin{solution}{exo:g12:seq:6}
\emph{1.} $u_0 = 0 \in \intcc{0}{2}$. If $0 \leq u_n \leq 2$, then
$2 \leq u_n + 2 \leq 4$, so $\sqrt{2} \leq u_{n+1} \leq 2$; in particular
$0 \leq u_{n+1} \leq 2$. By induction the property holds for all $n$.

\emph{2.} $u_{n+1} - u_n = \sqrt{u_n + 2} - u_n$. For $x \in \intcc{0}{2}$,
$\sqrt{x+2} \geq x \iff x + 2 \geq x^2 \iff (2-x)(x+1) \geq 0$, which is true.
Hence $(u_n)$ is increasing.

\emph{3.} Increasing and bounded above by $2$, $(u_n)$ converges to some
$\ell \in \intcc{0}{2}$. Passing to the limit in $u_{n+1} = \sqrt{u_n + 2}$
(the map $x \mapsto \sqrt{x+2}$ is continuous) gives $\ell = \sqrt{\ell + 2}$,
so $\ell^2 - \ell - 2 = 0$, i.e.\ $\ell \in \{-1, 2\}$. Since $\ell \geq 0$,
$\lim u_n = 2$.
\end{solution}

\begin{solution}{exo:g12:seq:7}
\emph{1.} Between two doses, $40\%$ of the drug is eliminated, so the
quantity $u_n$ becomes $0.6\,u_n$; the next dose adds $1$ unit:
$u_{n+1} = 0.6\,u_n + 1$.

\emph{2.} $v_{n+1} = u_{n+1} - 2.5 = 0.6\,u_n + 1 - 2.5 = 0.6(u_n - 2.5)
= 0.6\,v_n$: $(v_n)$ is geometric with ratio $0.6$ and first term
$v_0 = 1 - 2.5 = -1.5$. Hence $v_n = -1.5 \times 0.6^n$ and
\[
u_n = 2.5 - 1.5 \times 0.6^n .
\]

\emph{3.} Since $0.6^n \to 0$, $u_n \to 2.5$: the quantity of drug stabilizes
at $2.5$ units.
\end{solution}

\begin{solution}{exo:g12:seq:8}
\emph{1.} $u_0 = 3 > 1$. If $u_n > 1$, then $u_n + 2 > 0$ and
\[
u_{n+1} - 1 = \frac{4u_n - 1 - u_n - 2}{u_n + 2} = \frac{3(u_n - 1)}{u_n + 2} > 0 .
\]
By induction, $u_n > 1$ for all $n$ (and in particular $u_n + 2 \neq 0$, so
the sequence is well defined).

\emph{2.} Using the identity above,
\[
v_{n+1} = \frac{1}{u_{n+1} - 1} = \frac{u_n + 2}{3(u_n - 1)}
= \frac{(u_n - 1) + 3}{3(u_n - 1)} = \frac{1}{3} + v_n .
\]
So $(v_n)$ is arithmetic with common difference $\frac13$ and
$v_0 = \frac{1}{u_0 - 1} = \frac12$.

\emph{3.} $v_n = \frac12 + \frac{n}{3}$, hence
$u_n = 1 + \frac{1}{v_n} = 1 + \frac{6}{3 + 2n}$. Since $v_n \to +\infty$,
$u_n \to 1$.
\end{solution}

\begin{solution}{exo:g12:seq:9}
\emph{1.} $H_{2n} - H_n = \sum_{k=n+1}^{2n} \frac{1}{k}$ is a sum of $n$
terms, each at least $\frac{1}{2n}$; hence
$H_{2n} - H_n \geq n \cdot \frac{1}{2n} = \frac12$.

\emph{2.} By induction on $k$: $H_{2^0} = H_1 = 1 \geq 1$. If
$H_{2^k} \geq 1 + \frac{k}{2}$, then applying point 1 with $n = 2^k$,
\[
H_{2^{k+1}} \geq H_{2^k} + \frac12 \geq 1 + \frac{k+1}{2}.
\]
The sequence $(H_n)$ is increasing (each step adds $\frac{1}{n+1} > 0$) and
the subsequence $H_{2^k}$ is unbounded, so $(H_n)$ is not bounded above.
Increasing and unbounded, it tends to $+\infty$
(\cref{thm:g12:seq:monotone}).
\end{solution}

\begin{solution}{exo:g12:seq:10}
\emph{1.} The sequence $d_n = b_n - a_n$ satisfies
$d_{n+1} - d_n = (b_{n+1} - b_n) - (a_{n+1} - a_n) \leq 0$, so $(d_n)$ is
decreasing; since $d_n \to 0$, we get $d_n \geq 0$ for all $n$ (a decreasing
sequence with a negative term would stay below it forever, preventing the
limit $0$). Hence $a_n \leq b_n$.

\emph{2.} From $a_n \leq b_n \leq b_0$, the increasing sequence $(a_n)$ is
bounded above, so it converges to some $\ell$. Similarly $(b_n)$, decreasing
and bounded below by $a_0$, converges to some $\ell'$. Then
$\ell' - \ell = \lim (b_n - a_n) = 0$, so $\ell = \ell'$.

\emph{3.} $(a_n)$ is (strictly) increasing since
$a_{n+1} - a_n = \frac{1}{(n+1)!} > 0$. For $(b_n)$,
\[
b_{n+1} - b_n = \frac{1}{(n+1)!} + \frac{1}{(n+1)(n+1)!} - \frac{1}{n\,n!}
= \frac{n(n+1) + n - (n+1)^2}{n(n+1)(n+1)!}
= \frac{-1}{n(n+1)(n+1)!} < 0 ,
\]
so $(b_n)$ is decreasing. Finally $b_n - a_n = \frac{1}{n\,n!} \to 0$. The
two sequences are adjacent, hence converge to a common limit.
\end{solution}

\section*{Chapter \ref{ch:g12:limcont} --- Limits and Continuity}

\begin{solution}{exo:g12:limcont:1}
Dividing by $x^3$:
$\dfrac{2x^3 - x + 1}{5x^3 + x^2} = \dfrac{2 - 1/x^2 + 1/x^3}{5 + 1/x}
\to \dfrac{2}{5}$.

$x^2 + 3x - 1 = x^2(1 + 3/x - 1/x^2) \to +\infty$ as $x \to -\infty$, since
$x^2 \to +\infty$ and the bracket tends to $1$.

As $x \to 2^+$, the numerator tends to $3 > 0$ and the denominator to $0^+$,
so $\dfrac{x+1}{x-2} \to +\infty$.

For $x > 0$, $-\dfrac1x \leq \dfrac{\cos x}{x} \leq \dfrac1x$, so the limit
is $0$ by the squeeze theorem.
\end{solution}

\begin{solution}{exo:g12:limcont:2}
\emph{1.} $2x^2 - 3x + 1 = (x-1)(2x-1)$ (the \omterm{def:g11:quad:discriminant}{root} $x=1$ is apparent from
$2 - 3 + 1 = 0$). For $x \neq 1$, cancel $x - 1$: $f(x) = 2x - 1$.

\emph{2.} Hence $\lim_{x\to1} f(x) = 1$. Setting $f(1) = 1$ extends $f$ to
the polynomial \omterm{def:g11:func:function}{function} $x \mapsto 2x-1$ on $\R$, which is \omterm{def:g12:limcont:continuity}{continuous} at $1$.
\end{solution}

\begin{solution}{exo:g12:limcont:3}
Write $f(x) = \dfrac{3x - 1}{x + 2} = \dfrac{3(x+2) - 7}{x+2}
= 3 - \dfrac{7}{x+2}$. As $x \to \pm\infty$, $f(x) \to 3$: the line $y = 3$
is a \omterm{def:g12:limcont:asymptote}{horizontal asymptote} at both infinities. As $x \to (-2)^\pm$,
$f(x) \to \mp\infty$: the line $x = -2$ is a \omterm{def:g12:limcont:asymptote}{vertical asymptote}. The curve is
a \omterm{def:g11:func:reference}{hyperbola}, below $y=3$ for $x > -2$ and above it for $x < -2$.
\end{solution}

\begin{solution}{exo:g12:limcont:4}
Multiply by the conjugate:
\[
\sqrt{x^2 + x} - x = \frac{x^2 + x - x^2}{\sqrt{x^2+x} + x}
= \frac{x}{x\left(\sqrt{1 + 1/x} + 1\right)}
= \frac{1}{\sqrt{1 + 1/x} + 1}
\xrightarrow[x\to+\infty]{} \frac{1}{2}.
\]
Similarly,
\[
\frac{\sqrt{1+x} - 1}{x} = \frac{(1+x) - 1}{x\left(\sqrt{1+x} + 1\right)}
= \frac{1}{\sqrt{1+x} + 1} \xrightarrow[x\to0]{} \frac{1}{2}.
\]
\end{solution}

\begin{solution}{exo:g12:limcont:5}
\emph{1.} Since $\sin x \geq -1$, $f(x) \geq x - 1 \to +\infty$, and we
conclude by comparison. Yet $f'(x) = 1 + \cos x$ vanishes at every odd
multiple of $\pi$, and on every \omterm{def:g10:numbers:interval}{interval} $\intco{A}{+\infty}$ the \omterm{def:g11:func:function}{function}
alternates between \omterm{def:g10:numbers:interval}{intervals} of strict increase; it is \omterm{def:g12:seq:monotonic}{increasing} but not
strictly, and more importantly its growth is not needed: comparison suffices.

\emph{2.} For $x > 1$, divide by $x$:
\[
g(x) = \frac{1 + \frac{\sin x}{x}}{1 - \frac{\sin x}{x}},
\]
and $\frac{\sin x}{x} \to 0$ by the squeeze theorem
($\abs{\sin x} \leq 1$). Hence $g(x) \to \frac{1+0}{1-0} = 1$.
\end{solution}

\begin{solution}{exo:g12:limcont:6}
$f(x) = x^3 - 3x + 1$ is \omterm{def:g12:limcont:continuity}{continuous}, with $f'(x) = 3x^2 - 3 = 3(x-1)(x+1)$:
$f$ increases on $\intoc{-\infty}{-1}$, decreases on $\intcc{-1}{1}$,
increases on $\intco{1}{+\infty}$, with local \omterm{def:g10:functions:extrema}{maximum} $f(-1) = 3 > 0$ and
local \omterm{def:g10:functions:extrema}{minimum} $f(1) = -1 < 0$, $\lim_{-\infty} f = -\infty$,
$\lim_{+\infty} f = +\infty$.

On each of the three monotonicity \omterm{def:g10:numbers:interval}{intervals}, $f$ is \omterm{def:g12:limcont:continuity}{continuous}, strictly
\omterm{def:g12:seq:monotonic}{monotonic}, and $0$ lies strictly between the endpoint values/limits, so the
bijection theorem gives exactly one zero per \omterm{def:g10:numbers:interval}{interval}: three zeros in total.
Sign checks: $f(-2) = -1 < 0 < f(-1) = 3$, so one zero in
$\intoo{-2}{-1}$; $f(0) = 1 > 0 > f(1) = -1$, one zero in $\intoo{0}{1}$;
$f(1) = -1 < 0 < f(2) = 3$, one zero in $\intoo{1}{2}$.
\end{solution}

\begin{solution}{exo:g12:limcont:7}
Let $g(x) = f(x) - x$, \omterm{def:g12:limcont:continuity}{continuous} on $\intcc{0}{1}$. Since
$f(0) \in \intcc{0}{1}$, $g(0) = f(0) \geq 0$; since $f(1) \in \intcc{0}{1}$,
$g(1) = f(1) - 1 \leq 0$. By the intermediate value theorem, $g$ vanishes at
some $c \in \intcc{0}{1}$, i.e.\ $f(c) = c$.
\end{solution}

\begin{solution}{exo:g12:limcont:8}
Let $u(t)$ and $d(t)$ be the positions (measured as distance from the bottom
of the trail) at time $t \in \intcc{8}{20}$ on the way up and on the way
down. Both are \omterm{def:g12:limcont:continuity}{continuous}, $u(8) = 0$, $u(20) = L$ (the length of the trail),
$d(8) = L$, $d(20) = 0$. The \omterm{def:g11:func:function}{function} $g = u - d$ is \omterm{def:g12:limcont:continuity}{continuous} with
$g(8) = -L \leq 0$ and $g(20) = L \geq 0$, so $g(c) = 0$ for some time $c$ by
the intermediate value theorem: at time $c$, both days' positions coincide.
\end{solution}

\begin{solution}{exo:g12:limcont:9}
\emph{1.} $f'(x) = 5x^4 + 3x^2 + 1 > 0$, so $f$ is strictly \omterm{def:g12:seq:monotonic}{increasing} on
$\R$; with $\lim_{\pm\infty} f = \pm\infty$ and \omterm{def:g12:limcont:continuity}{continuity}, the bijection
theorem gives a unique real zero $\alpha$. Since $f(0) = -1 < 0$ and
$f(1) = 2 > 0$, $\alpha \in \intoo{0}{1}$.

\emph{2.} Maintain an \omterm{def:g10:numbers:interval}{interval} $\intcc{a_n}{b_n}$ with
$f(a_n) < 0 < f(b_n)$, starting from $\intcc{0}{1}$; at each step evaluate
$f$ at the \omterm{prop:g10:coordgeom:midpoint}{midpoint} $m$ and keep the \omterm{def:g10:numbers:interval}{half-interval} on which the sign change
persists. After $n$ steps the \omterm{def:g10:numbers:interval}{interval} has length $2^{-n}$, so the \omterm{prop:g10:coordgeom:midpoint}{midpoint}
approximates $\alpha$ to within $2^{-(n+1)}$. We need
$2^{-(n+1)} \leq 10^{-6}$, i.e.\ $n + 1 \geq \frac{6}{\log_{10} 2} \approx 19.93$:
$n = 20$ iterations suffice.
\end{solution}

\section*{Chapter \ref{ch:g12:deriv} --- Differentiation and Convexity}

\begin{solution}{exo:g12:deriv:1}
$f'(x) = 3x^2 - 5$.

Quotient rule:
$g'(x) = \dfrac{(x^2+1) - x\cdot 2x}{(x^2+1)^2} = \dfrac{1 - x^2}{(x^2+1)^2}$.

Chain rule with $u = 2x+1$: $h'(x) = 5 \cdot 2 \cdot (2x+1)^4 = 10(2x+1)^4$.

Chain rule with $u = x^2 + 1$:
$k'(x) = \dfrac{2x}{2\sqrt{x^2+1}} = \dfrac{x}{\sqrt{x^2+1}}$.
\end{solution}

\begin{solution}{exo:g12:deriv:2}
$f'(x) = 2x - 3$, so $f(2) = -1$ and $f'(2) = 1$: the tangent is
$y = -1 + (x - 2) = x - 3$. The gap is
\[
f(x) - (x - 3) = x^2 - 4x + 4 = (x-2)^2 \geq 0 ,
\]
so the curve lies above the tangent everywhere, touching it only at $x = 2$
(as expected: $f$ is convex).
\end{solution}

\begin{solution}{exo:g12:deriv:3}
Both are difference quotients. With $f(x) = (1+x)^{100}$ at $0$:
$f'(x) = 100(1+x)^{99}$, so the limit is $f'(0) = 100$.
With $g(x) = \sqrt{x}$ at $4$: $g'(x) = \frac{1}{2\sqrt{x}}$, so the limit is
$g'(4) = \frac14$.
\end{solution}

\begin{solution}{exo:g12:deriv:4}
Domain $\R \setminus \{1\}$. Limits: at $\pm\infty$, $f(x) \sim x \to
\pm\infty$; at $1^\pm$, numerator $\to 4 > 0$ and denominator $\to 0^\pm$,
so $f \to \pm\infty$; the line $x = 1$ is a vertical asymptote.

\[
f'(x) = \frac{2x(x-1) - (x^2+3)}{(x-1)^2} = \frac{x^2 - 2x - 3}{(x-1)^2}
= \frac{(x-3)(x+1)}{(x-1)^2}.
\]
Hence $f' > 0$ on $\intoo{-\infty}{-1}$ and $\intoo{3}{+\infty}$, $f' < 0$ on
$\intoo{-1}{1}$ and $\intoo{1}{3}$. The function increases to a local maximum
$f(-1) = -2$, decreases to $-\infty$, jumps to $+\infty$ after the asymptote,
decreases to a local minimum $f(3) = 6$, then increases.
\end{solution}

\begin{solution}{exo:g12:deriv:5}
For $x \in \intoo{0}{15}$, the box has a square base of side $30 - 2x$ and
height $x$, so
\[
V(x) = x(30 - 2x)^2 .
\]
$V'(x) = (30-2x)^2 + x \cdot 2(30-2x)(-2) = (30-2x)(30 - 2x - 4x)
= (30-2x)(30-6x)$. On $\intoo{0}{15}$, $30 - 2x > 0$, so $V'$ has the sign of
$30 - 6x$: positive before $x = 5$, negative after. The volume is maximal for
$x = 5$ cm, giving $V(5) = 5 \times 20^2 = 2000\ \text{cm}^3$.
\end{solution}

\begin{solution}{exo:g12:deriv:6}
\emph{1.} $f'(x) = 4x^3 - 12x^2$ and $f''(x) = 12x^2 - 24x = 12x(x-2)$.
Thus $f'' > 0$ on $\intoo{-\infty}{0}$ and on $\intoo{2}{+\infty}$ (convex),
$f'' < 0$ on $\intoo{0}{2}$ (concave). $f''$ changes sign at $0$ and at $2$:
both are inflection points.

\emph{2.} At $a = 2$: $f(2) = 16 - 32 + 10 = -6$, $f'(2) = 32 - 48 = -16$,
tangent $y = -6 - 16(x-2) = -16x + 26$. The gap is
\[
g(x) = f(x) + 16x - 26 = x^4 - 4x^3 + 16x - 16 .
\]
Since $g(2) = g'(2) = g''(2) = 0$, $(x-2)^3$ divides $g$; division gives
$g(x) = (x-2)^3(x+2)$. Near $x = 2$, $x + 2 > 0$, so $g$ has the sign of
$(x-2)^3$: negative before $2$, positive after. The curve crosses the
tangent: $2$ is indeed an inflection point.
\end{solution}

\begin{solution}{exo:g12:deriv:7}
The function $f(x) = \sqrt{x}$ satisfies
$f''(x) = -\frac{1}{4}x^{-3/2} < 0$ on $\intoo{0}{+\infty}$: it is concave,
so its graph lies \emph{below} each tangent. The tangent at $a = 1$ is
\[
y = f(1) + f'(1)(x - 1) = 1 + \tfrac12 (x-1) = \frac{x+1}{2},
\]
whence $\sqrt{x} \leq \frac{x+1}{2}$ for all $x > 0$. Equality means the
curve touches the tangent, which happens only at the point of tangency
$x = 1$. (Equivalently: $\left(\sqrt x - 1\right)^2 \geq 0$.)
\end{solution}

\begin{solution}{exo:g12:deriv:8}
By \cref{thm:g12:deriv:convexity}, the graph of a differentiable convex
function lies above each of its tangents. The tangent at $a$ is horizontal
($f'(a) = 0$), of equation $y = f(a)$. Hence $f(x) \geq f(a)$ for all
$x \in \R$: $f(a)$ is the global minimum.
\end{solution}

\begin{solution}{exo:g12:deriv:9}
\emph{Fixed perimeter.} A rectangle with sides $x$ and $\frac{p}{2} - x$
($0 < x < \frac p2$) has area $S(x) = x\left(\frac p2 - x\right)$. Then
$S'(x) = \frac p2 - 2x$ vanishes at $x = \frac p4$, positive before and
negative after: the maximum is at $x = \frac{p}{4}$, where both sides equal
$\frac p4$ --- a square.

\emph{Fixed area.} Sides $x$ and $\frac Ax$ give perimeter
$P(x) = 2\left(x + \frac Ax\right)$, with
$P'(x) = 2\left(1 - \frac{A}{x^2}\right)$, negative for $x < \sqrt A$ and
positive after: minimum at $x = \sqrt{A}$, a square again.

\emph{Relation.} The two statements are dual. Suppose a non-square rectangle
minimized the perimeter at fixed area $A$; the square of the same perimeter
would have area $> A$ by the first statement, and shrinking it homothetically
to area $A$ would strictly decrease its perimeter --- contradicting
minimality. Each statement thus implies the other.
\end{solution}

\section*{Chapter \ref{ch:g12:exp} --- Exponential and Logarithm}

\begin{solution}{exo:g12:exp:1}
$\dfrac{\eu^{3x}\,\eu^{-x+1}}{\eu^{x}} = \eu^{3x - x + 1 - x} = \eu^{x+1}$.

$\ln\!\left(\dfrac{\eu^2\sqrt{\eu}}{\eu^{-3}}\right)
= 2 + \tfrac12 + 3 = \tfrac{11}{2}$.
\end{solution}

\begin{solution}{exo:g12:exp:2}
\emph{(a)} Set $X = \eu^x > 0$: $X^2 - 3X + 2 = 0$ gives $X = 1$ or $X = 2$,
both positive, so $x = 0$ or $x = \ln 2$.

\emph{(b)} \omterm{def:g11:func:function}{Domain}: $x - 1 > 0$ and $x + 2 > 0$, so $x > 1$. The \omterm{def:g10:algebra:equation}{equation}
becomes $\ln\bigl((x-1)(x+2)\bigr) = \ln 4$, hence $(x-1)(x+2) = 4$,
\emph{i.e.} $x^2 + x - 6 = 0$, so $x = 2$ or $x = -3$. Only $x = 2$ is in the
\omterm{def:g11:func:function}{domain}: $S = \{2\}$.
\end{solution}

\begin{solution}{exo:g12:exp:3}
Divide by $\eu^x$:
$\dfrac{1 - x^2\eu^{-x}}{1 + \eu^{-x}} \to \dfrac{1-0}{1+0} = 1$, using
$x^2 \eu^{-x} \to 0$ (\cref{thm:g12:exp:growth}).

$\dfrac{\ln x}{\sqrt x} \to 0$ by \cref{prop:g12:exp:lnanalytic} (case
$n = 2$).

$x^2 \ln x = x \cdot (x \ln x) \to 0 \times 0 = 0$.

$x - \ln x = x\left(1 - \frac{\ln x}{x}\right) \to +\infty$ since
$\frac{\ln x}{x} \to 0$.
\end{solution}

\begin{solution}{exo:g12:exp:4}
$f'(x) = \eu^{-x} - x\eu^{-x} = (1 - x)\eu^{-x}$, of the sign of $1 - x$:
$f$ increases on $\intoc{-\infty}{1}$, decreases on $\intco{1}{+\infty}$,
with a global \omterm{def:g10:functions:extrema}{maximum} $f(1) = \eu^{-1}$.

Limits: as $x \to +\infty$, $f(x) = \frac{x}{\eu^x} \to 0$
(\cref{thm:g12:exp:growth}); as $x \to -\infty$, $x \to -\infty$ and
$\eu^{-x} \to +\infty$, so $f(x) \to -\infty$.

$f''(x) = -\eu^{-x} - (1-x)\eu^{-x} = (x - 2)\eu^{-x}$, which changes sign at
$x = 2$: \omterm{def:g12:deriv:inflection}{inflection point} at $\bigl(2,\, 2\eu^{-2}\bigr)$.
\end{solution}

\begin{solution}{exo:g12:exp:5}
Let $g(x) = x - \ln(1+x)$ on $\intoo{-1}{+\infty}$. Then
$g'(x) = 1 - \frac{1}{1+x} = \frac{x}{1+x}$, negative on $\intoo{-1}{0}$ and
positive on $\intoo{0}{+\infty}$: $g$ attains its \omterm{def:g10:functions:extrema}{minimum} $g(0) = 0$, so
$g \geq 0$ with equality only at $0$. Hence $\ln(1+x) \leq x$.

Apply this with $x = \frac1n$:
$n\ln\left(1 + \frac1n\right) \leq 1$, so
$\left(1+\frac1n\right)^n = \eu^{\,n\ln(1+1/n)} \leq \eu$.
Apply it with $x = -\frac{1}{n+1} > -1$:
$\ln\left(1 - \frac{1}{n+1}\right) \leq -\frac{1}{n+1}$, so
$-(n+1)\ln\left(1 - \frac{1}{n+1}\right) \geq 1$ and
$\left(1 - \frac{1}{n+1}\right)^{-(n+1)} \geq \eu$.
\end{solution}

\begin{solution}{exo:g12:exp:6}
\emph{1.} Each year multiplies the capital by $1.03$: $C_n = C_0 (1.03)^n$.

\emph{2.} $C_n \geq 2C_0 \iff (1.03)^n \geq 2 \iff n \ln 1.03 \geq \ln 2
\iff n \geq \frac{\ln 2}{\ln 1.03} \approx 23.45$: the capital doubles after
$24$ years.

\emph{3.} $\frac{70}{3} \approx 23.3$, close to the exact
$\frac{\ln 2}{\ln 1.03}$. The rule works because
$\ln(1 + r) \approx r$ for small $r$, so
$\frac{\ln 2}{\ln(1+r)} \approx \frac{0.693}{r} \approx \frac{70}{100r}$.
\end{solution}

\begin{solution}{exo:g12:exp:7}
$\eu^{2x} > \eu^{x+1} \iff 2x > x + 1 \iff x > 1$ (the exponential is
strictly \omterm{def:g12:seq:monotonic}{increasing}): $S = \intoo{1}{+\infty}$.

\omterm{def:g11:func:function}{Domain} of the second inequality: $x^2 - 1 > 0$ and $x + 5 > 0$, \emph{i.e.}
$x \in \intoo{-5}{-1} \cup \intoo{1}{+\infty}$. On this \omterm{def:g11:func:function}{domain}, $\ln$ being
strictly \omterm{def:g12:seq:monotonic}{increasing},
\[
\ln(x^2-1) \leq \ln(x+5) \iff x^2 - 1 \leq x + 5 \iff x^2 - x - 6 \leq 0
\iff x \in \intcc{-2}{3}.
\]
Intersecting with the \omterm{def:g11:func:function}{domain}: $S = \intco{-2}{-1} \cup \intoc{1}{3}$.
\end{solution}

\begin{solution}{exo:g12:exp:8}
\emph{1.} $f'(x) = \dfrac{\eu^x x - \eu^x}{x^2} = \dfrac{(x-1)\eu^x}{x^2}$,
negative on $\intoo{0}{1}$, positive on $\intoo{1}{+\infty}$: \omterm{def:g10:functions:extrema}{minimum}
$f(1) = \eu$. Limits: $f \to +\infty$ at $0^+$ (numerator $\to 1$,
denominator $\to 0^+$) and at $+\infty$ (\cref{thm:g12:exp:growth}).

\emph{2.} For $x > 0$, $\eu^x = kx \iff f(x) = k$. From the \omterm{def:g10:functions:table}{variation table}
($+\infty \searrow \eu \nearrow +\infty$, \omterm{def:g12:limcont:continuity}{continuous} and strictly \omterm{def:g12:seq:monotonic}{monotonic}
on each side): no solution for $k < \eu$; exactly one ($x = 1$) for
$k = \eu$; exactly two for $k > \eu$ (one in $\intoo{0}{1}$, one in
$\intoo{1}{+\infty}$, by the bijection theorem on each \omterm{def:g10:numbers:interval}{interval}).
\end{solution}

\begin{solution}{exo:g12:exp:9}
\emph{1.} Direct from the first inequality of \cref{exo:g12:exp:5}:
$u_n = \eu^{\,n\ln(1+1/n)} \leq \eu^1 = \eu$.

\emph{2.} Write
\[
\ln u_n = n \ln\!\left(1 + \frac1n\right)
= \frac{\ln\!\left(1 + \frac1n\right) - \ln 1}{\frac1n},
\]
a difference quotient of $\ln$ at the point $1$ with increment
$h = \frac1n \to 0$. Since $\ln$ is \omterm{def:g12:deriv:derivative}{differentiable} at $1$ with \omterm{def:g12:deriv:derivative}{derivative}
$1$, $\ln u_n \to 1$. By \omterm{def:g12:limcont:continuity}{continuity} of $\exp$
(\cref{prop:g12:limcont:seqcont}), $u_n = \eu^{\ln u_n} \to \eu^1 = \eu$.
\end{solution}

\begin{solution}{pb:g12:exp:1}
\textbf{1.} $x = \frac{\ln 5}{2}$. \quad
$3x - 1 = \eu^2$: $x = \frac{\eu^2 + 1}{3}$. \quad With
$u = \eu^x$: $u^2 - 3u + 2 = (u - 1)(u - 2) = 0$: $x = 0$ or
$x = \ln 2$.

\textbf{2.} $\frac{3\ln 2}{\ln 2} = 3$; \quad
$3 + \frac12 = \frac72$; \quad $\frac52$.

\textbf{3.} $f'(x) = \eu^x - 1$: negative before $0$, positive
after: \omterm{def:g10:functions:extrema}{minimum} $f(0) = 0$, so $f \geq 0$ everywhere:
$\eu^x \geq 1 + x$. Substituting $x = \ln(1 + u)$:
$1 + u \geq 1 + \ln(1 + u)$, i.e.\ $\ln(1 + u) \leq u$.

\textbf{4.} $x^2\eu^{-x} = \frac{x^2}{\eu^x} \to 0$ and
$\frac{\ln x}{x} \to 0$: exponentials crush powers, powers
crush \omterm{def:g12:exp:ln}{logarithms}. And $\frac{\eu^x - 1}{x} =
\frac{\eu^x - \eu^0}{x - 0} \to \exp'(0) = 1$.

\textbf{5.} $f'(x) = \ln x + 1$: zero at $\eu^{-1}$, negative
before, positive after: \omterm{def:g10:functions:extrema}{minimum}
$f\!\left(\frac1\eu\right) = -\frac1\eu$; and
$x \ln x \to 0$ at $0^+$: the curve leaves the origin, dips to
$-\frac1\eu$, and climbs away.

\textbf{6.} $n = \frac{\ln 2}{\ln 1.03} \approx 23.4$ years.

\textbf{7.} Exact: $23.4$, $11.9$, $8.0$ years; rule of 72:
$24$, $12$, $8$. Since $\ln(1+r) \approx r$,
$\frac{\ln 2}{\ln(1+r)} \approx \frac{0.693}{r}$, i.e.\
$\frac{69.3}{\text{rate in }\%}$; bankers round up to $72$
because it divides beautifully by $2, 3, 4, 6, 8, 9, 12$ ---
mental arithmetic beats a decimal of accuracy.

\textbf{8.} $2^{t/20} = \eu^{t \ln 2 / 20}$:
$\lambda = \frac{\ln 2}{20}$ per minute. After
$24 \times 60 = 1440$ minutes: $2^{72} \approx 4.7 \times
10^{21}$ bacteria --- more than the grains of sand on Earth,
from one cell in one day. The model is honest only while food
and space last: real growth bends into the logistic S-curve
(the epidemiologist's curve of \cref{pb:g12:deriv:1}).

\textbf{9.} At 23:00 ($8$ hours):
$100 \times 2^{-8/5} \approx 33$ mg. Below $10$ mg:
$2^{-t/5} < 0.1$ gives $t > 5\,\frac{\ln 10}{\ln 2} \approx
16.6$ h: around 7:40 the next morning --- the espresso at
three has a long tail.

\textbf{10.} Yearly: $2$. Monthly:
$\left(1 + \frac{1}{12}\right)^{12} \approx 2.613$. Daily:
$\approx 2.715$. The ceiling: $\eu = 2.71828\ldots$
(\cref{exo:g12:exp:9}) --- compounding \omterm{def:g12:limcont:continuity}{continuously}, greed
\omterm{def:g12:seq:limit}{converges}.

\textbf{11.} If cases grow exponentially, $c(t) = c_0
\eu^{\lambda t}$, then $\ln c(t) = \ln c_0 + \lambda t$: a
straight line of \omterm{def:g10:reffunc:affine}{slope} $\lambda$ --- the growth rate. A
straight stretch on the log plot \emph{is} the exponential
phase, and its steepness is the epidemic's tempo.

\textbf{12.} $120 - 60 = 60$ dB means
$10\log_{10}(I_2/I_1) = 60$: the concert is
$10^6$ --- a million times --- more intense than the
conversation.

\textbf{13.} Amplitude: $10^{7-5} = 100$. Energy:
$100^{3/2} = 1000$: two magnitude points hide three orders of
magnitude in energy.

\textbf{14.} $10^{6.5 - 2} = 10^{4.5} \approx 32\,000$: lemon
juice is thirty thousand times more acidic than milk --- pH
compresses chemistry's chasms into a pocket scale.

\textbf{15.} $\log_2\!\frac{4186}{27.5} =
\frac{\ln(4186/27.5)}{\ln 2} \approx 7.25$: a piano spans just
over seven octaves. Senses respond to \emph{ratios} of stimuli
--- doubling the intensity feels like one step, whatever the
starting level --- so perception is built on a logarithmic
scale, and so are the units we invented for it.

\textbf{16.} Placing the stick for $a$ end-to-end with the
stick position for $b$ adds the lengths $\ln a + \ln b$, and
the graduation sitting at that total length reads
$\eu^{\ln a + \ln b} = ab$: the slide rule computes products
by adding \omterm{def:g12:exp:ln}{logarithms} --- $\ln(ab) = \ln a + \ln b$
(\cref{prop:g12:exp:lnalg}) carved in boxwood.

\textbf{17.} The \omterm{def:g10:functions:extrema}{minimum} of $\frac{\eu^x}{x}$ on
$\intoo{0}{+\infty}$ is $\eu$, at $x = 1$: no solution for
$k < \eu$, exactly one for $k = \eu$, two for $k > \eu$.
\omterm{def:g12:deriv:tangent}{Tangent} check: at $a = 1$ the \omterm{def:g12:deriv:tangent}{tangent} to $\eu^x$ is
$y = \eu + \eu(x - 1) = \eu x$: it passes through the origin
--- the threshold line, grazing the curve at $(1, \eu)$.

\textbf{18.} $n \ln\!\left(1 - \frac1n\right) =
\frac{\ln(1 - 1/n)}{-1/n} \times (-1) \to -1$, so
$\left(1 - \frac1n\right)^n \to \eu^{-1} \approx 0.368$: the
$0.999^{1000} \approx 0.368$ of \cref{pb:g11:binom:1},
explained --- the lottery's near-miss constant is
$\frac1\eu$.

\textbf{19.} Both are the limit shape of ``many independent
\omterm{def:g11:prob:model}{events}, each individually unlikely'': the chance that
\emph{none} of $n$ chances of size $\frac1n$ fires tends to
$\eu^{-1}$, and the secretary rule tunes its rejection window
so that success concentrates at that same constant.
$\frac1\eu \approx 0.368$.

\textbf{20.} The ceiling of compounding; the unique base whose
\omterm{def:g12:deriv:tangent}{tangent} at $0$ has \omterm{def:g10:reffunc:affine}{slope} $1$ (which is why calculus loves it);
the growth that outruns every power; the constant where
near-misses and optimal stopping settle; and the \omterm{def:g12:exp:ln}{logarithm} ---
multiplication become addition, the world's wildest ranges
folded onto a ruler.
\end{solution}

\section*{Chapter \ref{ch:g12:trigo} --- Trigonometric Functions}

\begin{solution}{exo:g12:trigo:1}
$\cos\frac{2\pi}{3} = \cos\left(\pi - \frac\pi3\right) = -\cos\frac\pi3
= -\frac12$.

$\sin\frac{5\pi}{6} = \sin\left(\pi - \frac\pi6\right) = \sin\frac\pi6
= \frac12$.

$\cos\frac{7\pi}{4} = \cos\left(2\pi - \frac{\pi}{4}\right)
= \cos\left(-\frac\pi4\right) = \frac{\sqrt2}{2}$.

$\sin\left(-\frac\pi3\right) = -\sin\frac\pi3 = -\frac{\sqrt3}{2}$.
\end{solution}

\begin{solution}{exo:g12:trigo:2}
$f'(x) = 2\sin x\cos x = \sin 2x$ (chain rule).

$g'(x) = -3\sin(3x) + \sin x + x\cos x$ (chain and product rules).

Quotient rule:
\[
h'(x) = \frac{\cos x\,(2 + \cos x) - \sin x\,(-\sin x)}{(2+\cos x)^2}
= \frac{2\cos x + \cos^2 x + \sin^2 x}{(2+\cos x)^2}
= \frac{2\cos x + 1}{(2+\cos x)^2}.
\]
\end{solution}

\begin{solution}{exo:g12:trigo:3}
\emph{(a)} A particular solution of $\sin x = \frac{\sqrt3}{2}$ is
$\frac\pi3$. General solutions: $x = \frac\pi3 + 2k\pi$ or
$x = \pi - \frac\pi3 + 2k\pi = \frac{2\pi}{3} + 2k\pi$. In
$\intoc{-\pi}{\pi}$: $x \in \left\{\frac\pi3, \frac{2\pi}{3}\right\}$.

\emph{(b)} $\cos 2x = \frac{\sqrt2}{2} = \cos\frac\pi4$ gives
$2x = \pm\frac\pi4 + 2k\pi$, so $x = \pm\frac\pi8 + k\pi$. In
$\intoc{-\pi}{\pi}$:
$x \in \left\{-\frac{7\pi}{8}, -\frac\pi8, \frac\pi8, \frac{7\pi}{8}\right\}$.
\end{solution}

\begin{solution}{exo:g12:trigo:4}
With $a = \frac\pi3$, $b = \frac\pi4$:
\[
\cos\frac{\pi}{12} = \cos a\cos b + \sin a \sin b
= \frac12\cdot\frac{\sqrt2}{2} + \frac{\sqrt3}{2}\cdot\frac{\sqrt2}{2}
= \frac{\sqrt2 + \sqrt6}{4},
\]
\[
\sin\frac{\pi}{12} = \sin a\cos b - \cos a \sin b
= \frac{\sqrt3}{2}\cdot\frac{\sqrt2}{2} - \frac12\cdot\frac{\sqrt2}{2}
= \frac{\sqrt6 - \sqrt2}{4}.
\]
\end{solution}

\begin{solution}{exo:g12:trigo:5}
Set $s = \sin x$: $2s^2 - s - 1 = (2s + 1)(s - 1) \geq 0$ iff
$s \leq -\frac12$ or $s = 1$.

$\sin x = 1$ gives $x = \frac\pi2$. $\sin x \leq -\frac12$: on the unit
circle, the arc below the horizontal line $y = -\frac12$, which in
$\intcc{0}{2\pi}$ is $x \in \intcc{\frac{7\pi}{6}}{\frac{11\pi}{6}}$.
Hence
\[
S = \left\{\frac{\pi}{2}\right\} \cup
\intcc{\frac{7\pi}{6}}{\frac{11\pi}{6}} .
\]
\end{solution}

\begin{solution}{exo:g12:trigo:6}
$\dfrac{\sin 3x}{x} = 3\,\dfrac{\sin 3x}{3x} \to 3 \cdot 1 = 3$
(composition with $u = 3x \to 0$).

Using $1 - \cos x = \dfrac{\sin^2 x}{1 + \cos x}$:
\[
\frac{1 - \cos x}{x^2}
= \left(\frac{\sin x}{x}\right)^{\!2} \frac{1}{1 + \cos x}
\xrightarrow[x\to0]{} 1 \cdot \frac12 = \frac12 .
\]

Substituting $u = \frac1x \to 0^+$:
$x \sin\frac1x = \frac{\sin u}{u} \to 1$.
\end{solution}

\begin{solution}{exo:g12:trigo:7}
$f'(x) = 1 - \sin x \geq 0$ for all $x$, with equality exactly when
$\sin x = 1$, \emph{i.e.} $x = \frac\pi2 + 2k\pi$. On $\intcc{0}{2\pi}$, $f$
is increasing from $f(0) = 1$ to $f(2\pi) = 2\pi + 1$ (the zero of $f'$ at
$\frac{\pi}{2}$ is isolated, so the increase is even strict); minimum $1$ at
$0$, maximum $2\pi + 1$ at $2\pi$. On $\R$, $f' \geq 0$ everywhere with only
isolated zeros, so $f$ is (strictly) increasing by
\cref{thm:g12:deriv:variations}, although $f'$ vanishes at every
$\frac\pi2 + 2k\pi$.
\end{solution}

\begin{solution}{exo:g12:trigo:8}
\emph{1.} Let $R = \sqrt{A^2 + B^2} > 0$. The point
$\left(\frac AR, \frac BR\right)$ lies on the unit circle, so there exists
$\varphi$ with $\cos\varphi = \frac AR$ and $\sin\varphi = \frac BR$. Then
\[
R\cos(x - \varphi) = R\bigl(\cos x\cos\varphi + \sin x \sin\varphi\bigr)
= A\cos x + B\sin x = f(x).
\]

\emph{2.} Since $\cos$ has range $\intcc{-1}{1}$, $f$ has maximum $R$ and
minimum $-R$. For $\cos x + \sin x$: $A = B = 1$, $R = \sqrt2$,
$\varphi = \frac\pi4$, so the equation becomes
$\sqrt2\cos\left(x - \frac\pi4\right) = 1$, \emph{i.e.}
$\cos\left(x - \frac\pi4\right) = \frac{\sqrt2}{2}$, giving
$x - \frac\pi4 = \pm\frac\pi4 + 2k\pi$: in $\intoc{-\pi}{\pi}$,
$x \in \left\{0, \frac\pi2\right\}$.
\end{solution}

\begin{solution}{exo:g12:trigo:9}
$(\sin)'' = (\cos)' = -\sin$ and $(\cos)'' = (-\sin)' = -\cos$: both solve
$y'' = -y$.

Let $g = y - a\cos - b\sin$. Then $g'' = -g$ (linearity of differentiation),
$g(0) = y(0) - a = 0$ and $g'(0) = y'(0) - b = 0$. Set
$E = g^2 + (g')^2$. Then
\[
E' = 2g g' + 2g' g'' = 2g g' - 2g' g = 0,
\]
so $E$ is constant equal to $E(0) = 0$. A sum of two squares vanishing
identically forces $g = 0$, hence $y = a\cos + b\sin$.
\end{solution}

\section*{Chapter \ref{ch:g12:integ} --- Integration}

\begin{solution}{exo:g12:integ:1}
\[
\int_1^2 \left(3x^2 - \frac{1}{x^2}\right)\dd x
= \left[x^3 + \frac1x\right]_1^2 = \left(8 + \tfrac12\right) - 2 = \frac{13}{2}.
\]
\[
\int_0^{\pi/2} \cos t\,\dd t = \bigl[\sin t\bigr]_0^{\pi/2} = 1 .
\]
With $u = 2t+1$, $\frac{1}{2t+1} = \frac12\,\frac{u'}{u}$:
\[
\int_0^1 \frac{\dd t}{2t+1} = \frac12\bigl[\ln(2t+1)\bigr]_0^1
= \frac{\ln 3}{2}.
\]
\end{solution}

\begin{solution}{exo:g12:integ:2}
$x\,\eu^{x^2} = \frac12\,(x^2)'\,\eu^{x^2}$ has \omterm{def:g12:integ:primitive}{primitives}
$\frac12\eu^{x^2} + c$. The condition $F(0) = 1$ gives $\frac12 + c = 1$,
so
\[
F(x) = \frac{\eu^{x^2} + 1}{2}.
\]
\end{solution}

\begin{solution}{exo:g12:integ:3}
\[
\mu = \frac{1}{\pi}\int_0^\pi \sin t\,\dd t
= \frac{1}{\pi}\bigl[-\cos t\bigr]_0^\pi = \frac{2}{\pi} \approx 0.64 .
\]
Graphically, the rectangle of base $\intcc{0}{\pi}$ and height
$\frac{2}{\pi}$ has the same area as the arch of the \omterm{def:g12:trigo:cossin}{sine} curve.
\end{solution}

\begin{solution}{exo:g12:integ:4}
Take $u' = 1$, $v = \ln t$, so $u = t$, $v' = \frac1t$:
\[
\int_1^{\eu} \ln t\,\dd t = \bigl[t\ln t\bigr]_1^{\eu} - \int_1^{\eu} 1\,\dd t
= \eu - (\eu - 1) = 1 .
\]
Take $u' = \sin t$, $v = t$, so $u = -\cos t$, $v' = 1$:
\[
\int_0^\pi t\sin t\,\dd t = \bigl[-t\cos t\bigr]_0^\pi + \int_0^\pi \cos t\,\dd t
= \pi + 0 = \pi .
\]
\end{solution}

\begin{solution}{exo:g12:integ:5}
The curves intersect where $x^2 = x + 2$, \emph{i.e.} $x \in \{-1, 2\}$, and
on $\intcc{-1}{2}$ the line is above the \omterm{def:g11:quad:parabola}{parabola}. The area is
\[
\int_{-1}^{2} \bigl(x + 2 - x^2\bigr)\dd x
= \left[\frac{x^2}{2} + 2x - \frac{x^3}{3}\right]_{-1}^{2}
= \frac{10}{3} - \left(-\frac{7}{6}\right) = \frac{9}{2}.
\]
\end{solution}

\begin{solution}{exo:g12:integ:6}
\emph{1.} $I = \bigl[\ln(1+t)\bigr]_0^1 = \ln 2$.

\emph{2.}
\[
I_n + I_{n+1} = \int_0^1 \frac{t^n(1 + t)}{1+t}\,\dd t
= \int_0^1 t^n \,\dd t = \frac{1}{n+1}.
\]
On $\intcc{0}{1}$, $0 \leq \dfrac{t^n}{1+t} \leq t^n$, so by comparison
$0 \leq I_n \leq \frac{1}{n+1}$.

\emph{3.} From the recurrence, by a telescoping induction,
\[
I_n = (-1)^n\left(I_0 - \left(1 - \frac12 + \dots + \frac{(-1)^{n-1}}{n}\right)\right),
\]
\emph{i.e.}
$\displaystyle\sum_{k=1}^{n} \frac{(-1)^{k-1}}{k} = \ln 2 - (-1)^n I_n$
(recall $I_0 = \ln 2$). Since $\abs{I_n} \leq \frac{1}{n+1} \to 0$, the
alternating harmonic sum tends to $\ln 2$.
\end{solution}

\begin{solution}{exo:g12:integ:7}
Distance:
\[
d = \int_0^{100} 30\left(1 - \eu^{-t/50}\right)\dd t
= 30\left[t + 50\,\eu^{-t/50}\right]_0^{100}
= 30\left(100 + 50\eu^{-2} - 50\right)
= 1500 + 1500\,\eu^{-2}.
\]
Numerically $d \approx 1500 + 203 = 1703$ m. Average speed:
$\frac{d}{100} = 15\bigl(1 + \eu^{-2}\bigr) \approx 17.0$ m/s.
\end{solution}

\begin{solution}{exo:g12:integ:8}
\emph{1.} $S_n$ is the total area of $n$ rectangles of width $\frac1n$ and
heights $f\left(\frac kn\right)$, $k = 1, \dots, n$, where
$f(t) = \frac{1}{1+t}$: it is the ``right endpoint'' rectangle \omterm{def:g10:numbers:approx}{approximation}
of the area under $f$ on $\intcc{0}{1}$.

\emph{2.} $f$ is \omterm{def:g10:functions:variations}{decreasing}, so on each \omterm{def:g10:numbers:interval}{interval}
$\intcc{\frac{k-1}{n}}{\frac kn}$,
$f\left(\frac kn\right) \leq f(t) \leq f\left(\frac{k-1}{n}\right)$;
integrating and summing over $k$:
\[
S_n \leq \int_0^1 f(t)\,\dd t \leq S_n + \frac1n\bigl(f(0) - f(1)\bigr)
= S_n + \frac{1}{2n}.
\]
(The middle term of the exercise's display, with
$f(0) - f(1) = 1 - \frac12$, is exactly this bound.) Hence
$0 \leq \ln 2 - S_n \leq \frac{1}{2n} \to 0$, so $S_n \to \ln 2$.
\end{solution}

\begin{solution}{exo:g12:integ:9}
\emph{1.} $W_0 = \int_0^{\pi/2} \dd t = \frac{\pi}{2}$ and
$W_1 = \bigl[-\cos t\bigr]_0^{\pi/2} = 1$.

\emph{2.} Integrate by parts with $u' = \sin t$, $v = \sin^{n+1} t$, so
$u = -\cos t$, $v' = (n+1)\sin^n t \cos t$:
\[
W_{n+2} = \bigl[-\cos t \sin^{n+1} t\bigr]_0^{\pi/2}
+ (n+1)\int_0^{\pi/2} \cos^2 t\,\sin^n t\,\dd t
= (n+1)\int_0^{\pi/2} (1 - \sin^2 t)\sin^n t\,\dd t,
\]
the bracket vanishing at both ends. Hence
$W_{n+2} = (n+1)(W_n - W_{n+2})$, \emph{i.e.}
$W_{n+2} = \frac{n+1}{n+2} W_n$.

\emph{3.} $W_2 = \frac12 W_0 = \frac{\pi}{4}$,
$W_3 = \frac23 W_1 = \frac23$,
$W_4 = \frac34 W_2 = \frac{3\pi}{16}$.
On $\intoo{0}{\frac\pi2}$, $0 < \sin t < 1$, so
$\sin^{n+1} t \leq \sin^n t$ with strict inequality on the interior;
integrating, $0 < W_{n+1} \leq W_n$ (in fact $<$): $(W_n)$ is \omterm{def:g10:functions:variations}{decreasing} and
positive.
\end{solution}

\begin{solution}{pb:g12:integ:1}
\textbf{1.} $\left[x^3 - x^2 + x\right]_0^1 = 1$; \quad
$\left[\ln x\right]_1^{\eu} = 1$; \quad
$\left[\sin x\right]_0^{\pi/2} = 1$.

\textbf{2.} $\left[\frac12 \eu^{x^2}\right]_0^1 =
\frac{\eu - 1}{2}$; \quad
$\left[\ln(x^2 + 1)\right]_0^1 = \ln 2$.

\textbf{3.} $\int_0^1 x\eu^x = \left[x\eu^x\right]_0^1 -
\int_0^1 \eu^x = \eu - (\eu - 1) = 1$. And
$\int_1^{\eu} \ln x = \left[x\ln x\right]_1^{\eu} -
\int_1^{\eu} 1 = \eu - (\eu - 1) = 1$.

\textbf{4.} $\frac{1}{\pi}\int_0^\pi \sin x\,\dd x =
\frac{2}{\pi} \approx 0.64$: the average height of the \omterm{def:g12:trigo:cossin}{sine}
arch is $\frac2\pi$, noticeably more than one half --- the
arch is fat near its top.

\textbf{5.} $\int_0^1 (x - x^2)\,\dd x = \frac12 - \frac13 =
\frac16$.

\textbf{6.} $L_n = \frac1n \sum_{k=0}^{n-1}
\left(\frac kn\right)^2$ and
$U_n = \frac1n \sum_{k=1}^{n} \left(\frac kn\right)^2$:
rectangles of width $\frac1n$ with heights read at the left
(below the curve) or right (above) endpoints.

\textbf{7.} True for $n = 1$ ($1 = \frac{1 \cdot 2 \cdot
3}{6}$). If it holds for $n$, then adding $(n+1)^2$:
\[
\frac{n(n+1)(2n+1)}{6} + (n+1)^2
= \frac{(n+1)\left(2n^2 + n + 6n + 6\right)}{6}
= \frac{(n+1)(n+2)(2n+3)}{6},
\]
the formula at $n + 1$: heredity, done.

\textbf{8.} $U_n = \frac{1}{n^3} \cdot
\frac{n(n+1)(2n+1)}{6} = \frac{(1 + \frac1n)(2 + \frac1n)}{6}
\to \frac26 = \frac13$, and $L_n = U_n - \frac1n \to \frac13$:
squeezed between its staircases, the area is exactly
$\frac13$.

\textbf{9.} $\int_0^1 x^2\,\dd x =
\left[\frac{x^3}{3}\right]_0^1 = \frac13$: one line. The
fundamental theorem converts an infinite squeeze into a single
\omterm{def:g12:deriv:derivative}{anti-derivative} evaluation --- that is why it is called
fundamental.

\textbf{10.} Segment: $\int_{-1}^{1}(1 - x^2)\,\dd x =
2 - \frac23 = \frac43$. Triangle with base the chord (length
$2$, at height $1$) and apex $(0,0)$: area
$\frac12 \times 2 \times 1 = 1$. Ratio: $\frac43$ --- exactly
Archimedes' theorem, verified by the machine he lacked.

\textbf{11.} $T\left(1 + \frac14 + \frac{1}{16} + \dots\right)
= T \cdot \frac{1}{1 - \frac14} = \frac43 T$ (geometric series
of ratio $\frac14$): with $T = 1$, the segment's area
$\frac43$ again. Archimedes summed the series by a purely
geometric argument --- the quarter-bites of the chocolate-bar
kind, two millennia before limits were words.

\textbf{12.} Exhaustion needs a clever geometric
self-similarity and gives exact areas case by case; Riemann
needs only monotonicity and gives a universal definition at
the price of computing sums; the fundamental theorem needs the
whole edifice of \omterm{def:g12:deriv:derivative}{derivatives} and repays it by making areas
one-line computations.

\textbf{13.} $A(1) = 0$ and $A'(x) = \frac1x$ (fundamental
theorem): $A$ is the \omterm{def:g12:integ:primitive}{primitive} of $\frac1x$ on
$\intoo{0}{+\infty}$ vanishing at $1$ --- which is precisely
the \omterm{def:g12:exp:ln}{natural logarithm}.

\textbf{14.} $g'(x) = \frac{a}{ax} - \frac1x = 0$: $g$ is
constant; at $x = 1$: $g(1) = A(a) - A(1) = A(a)$. Hence
$A(ax) = A(a) + A(x)$ for all $x$, and with $x = b$:
$\ln(ab) = \ln a + \ln b$. The \omterm{def:g11:func:reference}{hyperbola}'s area does not care
about horizontal scale: stretching $t$ by $a$ squashes
$\frac1t$ by $a$, and the area survives --- addition of
\omterm{def:g12:exp:ln}{logarithms} is a scaling symmetry.

\textbf{15.} Iterating: $\ln(a^n) = n\ln a$ (induction on
$n$). And $0 = \ln 1 = \ln\left(a \cdot \frac1a\right) =
\ln a + \ln\frac1a$: $\ln\frac1a = -\ln a$.

\textbf{16.} $\frac{1}{11} + \frac{1}{12} + \dots +
\frac{1}{20} \approx 0.669$, against $\ln 2 \approx 0.693$:
ten rectangles under $\frac{1}{1+t}$, already close. The
harmonic series diverges, but its \emph{slices} from $n$ to
$2n$ settle on $\ln 2$: divergence at a crawl, measured by an
area.

\textbf{17.} Distance: $\int_0^{10} 3t^2\,\dd t = 1000$ m.
\omterm{def:g11:stat:mean}{Mean} speed: $100$ m/s. Instant with $v = 100$:
$3t^2 = 100$: $t = \sqrt{\frac{100}{3}} \approx 5.77$ s.

\textbf{18.} Mass: $\int_0^4 (2 + x)\,\dd x = 8 + 8 = 16$ kg.
Moment: $\int_0^4 (2x + x^2)\,\dd x = 16 + \frac{64}{3} =
\frac{112}{3}$. Center of mass:
$\frac{112/3}{16} = \frac73 \approx 2.33$ m --- pushed beyond
the middle by the heavier right end, as intuition demands.

\textbf{19.} $W_0 = \frac\pi2$; $W_1 = \left[-\cos
t\right]_0^{\pi/2} = 1$; and
$W_2 = \int_0^{\pi/2} \frac{1 - \cos 2t}{2}\,\dd t =
\frac\pi4$.

\textbf{20.} Area: the definition, squeezed by staircases.
Accumulation: distances from speeds, masses from densities ---
the \omterm{def:g12:integ:area}{integral} as a running total. \omterm{def:g12:deriv:derivative}{Anti-derivative}: the
fundamental theorem, Newton's and Leibniz's shortcut through
Archimedes' patience. The pearl: the $\ln$ button is the area
under $\frac1t$, and that area's scaling symmetry \emph{is}
the identity $\ln(ab) = \ln a + \ln b$.
\end{solution}

\section*{Chapter \ref{ch:g12:diffeq} --- Differential Equations}

\begin{solution}{exo:g12:diffeq:1}
\emph{(a)} $y = C\eu^{3x}$; $y(0) = 2$ gives $y = 2\eu^{3x}$.

\emph{(b)} $y' = -\frac12 y$, so $y = C\eu^{-x/2}$; $y(0) = -1$ gives
$y = -\eu^{-x/2}$.

\emph{(c)} Equilibrium $y_p = 5$; $y = C\eu^{-x} + 5$; $y(0) = 0$ gives
$C = -5$: $y = 5\left(1 - \eu^{-x}\right)$.
\end{solution}

\begin{solution}{exo:g12:diffeq:2}
$N' = kN$, so $N(t) = N_0 \eu^{kt}$. Doubling in $3$ hours means
$\eu^{3k} = 2$, \emph{i.e.}
\[
k = \frac{\ln 2}{3} \approx 0.231\ \text{h}^{-1}.
\]
\end{solution}

\begin{solution}{exo:g12:diffeq:3}
$y_p'(x) = \eu^x + x\eu^x = y_p(x) + \eu^x$: $y_p$ is a particular solution
(this is the resonant case of \cref{met:g12:diffeq:particular}). By
\cref{prop:g12:diffeq:general}, the solutions are
$y = C\eu^{x} + x\eu^{x} = (C + x)\,\eu^x$, $C \in \R$.
\end{solution}

\begin{solution}{exo:g12:diffeq:4}
Try $y_p = \beta\eu^x$: $\beta\eu^x = -2\beta\eu^x + \eu^x$ gives
$3\beta = 1$, so $y_p = \frac13\eu^x$. General solution
$y = C\eu^{-2x} + \frac13\eu^{x}$; the condition $y(0) = 1$ gives
$C = \frac23$:
\[
y(x) = \frac{2}{3}\,\eu^{-2x} + \frac{1}{3}\,\eu^{x}.
\]
\end{solution}

\begin{solution}{exo:g12:diffeq:5}
$N(t) = N_0\,\eu^{-\lambda t}$ with $\lambda = \frac{\ln 2}{5730}$
(\cref{ex:g12:exp:halflife}). We solve $\eu^{-\lambda t} = 0.6$:
\[
t = \frac{\ln(1/0.6)}{\lambda} = 5730\,\frac{\ln(5/3)}{\ln 2}
\approx 5730 \times \frac{0.5108}{0.6931} \approx 4220 \text{ years}.
\]
\end{solution}

\begin{solution}{exo:g12:diffeq:6}
\emph{1.} Salt flows in at $0.2 \times 5 = 1$ kg/min. The outflow carries
concentration $\frac{m}{100}$ kg/L at $5$ L/min, \emph{i.e.} $\frac{m}{20}$
kg/min. Hence $m' = 1 - \frac{m}{20}$.

\emph{2.} Equilibrium $m_p = 20$; $m(t) = 20 + C\eu^{-t/20}$, and
$m(0) = 0$ gives $C = -20$:
\[
m(t) = 20\left(1 - \eu^{-t/20}\right) \xrightarrow[t\to+\infty]{} 20 \text{ kg}.
\]
In the long run the tank's concentration equals that of the incoming brine:
$0.2$ kg/L $\times$ $100$ L $= 20$ kg.
\end{solution}

\begin{solution}{exo:g12:diffeq:7}
\emph{1.} $v' = g - \frac{k}{m} v$ with $\frac km = \frac{16}{80} = 0.2$.
Equilibrium $v_\infty = \frac{mg}{k} = \frac{9.8}{0.2} = 49$ m/s;
$v(t) = 49 + C\eu^{-0.2t}$, and $v(0) = 0$ gives
\[
v(t) = 49\left(1 - \eu^{-0.2 t}\right).
\]

\emph{2.} Terminal velocity $49$ m/s ($\approx 176$ km/h). We want
$1 - \eu^{-0.2t} = 0.95$, \emph{i.e.} $\eu^{-0.2t} = 0.05$:
\[
t = \frac{\ln 20}{0.2} \approx \frac{3.00}{0.2} \approx 15 \text{ s}.
\]
\end{solution}

\begin{solution}{exo:g12:diffeq:8}
\emph{1.} $z = \frac1y$ gives $z' = -\frac{y'}{y^2}
= -\frac{y(1-y)}{y^2} = -\frac{1-y}{y} = -\frac1y + 1 = -z + 1$.

\emph{2.} Equilibrium $z_p = 1$, so $z(t) = 1 + C\eu^{-t}$. From
$y(0) = \frac{1}{10}$, $z(0) = 10$, so $C = 9$ and
\[
y(t) = \frac{1}{1 + 9\,\eu^{-t}} .
\]

\emph{3.} As $t \to +\infty$, $9\eu^{-t} \to 0$ and $y(t) \to 1$. The curve
is the classic S-shaped (\emph{sigmoid}) logistic curve: slow growth at
first ($y$ small, $y' \approx y$), fastest growth when $y = \frac12$
(where $y' = y(1-y)$ is maximal), then saturation towards the carrying
capacity $1$.
\end{solution}

\begin{solution}{exo:g12:diffeq:9}
By \cref{thm:g12:diffeq:affine}, $y(x) = C\eu^{ax} - \frac ba$. Hence
\[
u_{n+1} = C\eu^{a(n+1)} - \frac ba
= \eu^{a}\left(C\eu^{an} - \frac ba\right) + \frac ba\left(\eu^a - 1\right)
= q\,u_n + r
\]
with $q = \eu^{a}$ and $r = \frac{b}{a}\left(\eu^{a} - 1\right)$.
The condition $\abs q < 1$ means $\eu^a < 1$, \emph{i.e.} $a < 0$: exactly
the condition under which the solutions of the \omterm{def:g12:diffeq:ode}{differential equation}
converge to the equilibrium $-\frac ba$ --- and indeed the fixed point of
the recurrence is $\frac{r}{1 - q} = -\frac{b}{a}$.
\end{solution}

\section*{Chapter \ref{ch:g12:comb} --- Combinatorics and Counting}

\begin{solution}{exo:g12:comb:1}
Multiplication principle:
$26^2 \times 10^3 \times 26^2 = 26^4 \times 10^3 = 456\,976\,000$.

Without repeated characters, the four letters must be distinct
($26 \times 25 \times 24 \times 23$ ways, filling the letter positions in
order) and the three digits distinct ($10 \times 9 \times 8$):
\[
26 \times 25 \times 24 \times 23 \times 10 \times 9 \times 8
= 358\,800 \times 720 = 258\,336\,000 .
\]
\end{solution}

\begin{solution}{exo:g12:comb:2}
$\dbinom83 = \dfrac{8 \times 7 \times 6}{3!} = 56$;
$\dbinom{10}{8} = \dbinom{10}{2} = \dfrac{10 \times 9}{2} = 45$;
\[
\frac{\binom n2}{\binom{n+1}2}
= \frac{n(n-1)/2}{(n+1)n/2} = \frac{n-1}{n+1}.
\]
\end{solution}

\begin{solution}{exo:g12:comb:3}
Choose the committee: $\binom{30}{4}$ ways. Then choose president and
treasurer among the 4, in order: $4 \times 3 = 12$ ways. Total
\[
\binom{30}{4} \times 12 = 27\,405 \times 12 = 328\,860 .
\]
\end{solution}

\begin{solution}{exo:g12:comb:4}
\[
(x+2)^5 = x^5 + 10x^4 + 40x^3 + 80x^2 + 80x + 32 ,
\]
\[
(1-x)^6 = 1 - 6x + 15x^2 - 20x^3 + 15x^4 - 6x^5 + x^6 .
\]
In $(2x+3)^7$, the term in $x^3$ is
$\binom{7}{3}(2x)^3\,3^4 = 35 \times 8 \times 81\, x^3$: the coefficient is
$22\,680$.
\end{solution}

\begin{solution}{exo:g12:comb:5}
\emph{1.} $\dbinom{52}{5} = 2\,598\,960$.

\emph{2.} Exactly one ace: choose it ($4$ ways) and complete with $4$
non-aces: $4 \times \binom{48}{4} = 4 \times 194\,580 = 778\,320$.
At least one ace: complementary counting,
$\binom{52}{5} - \binom{48}{5} = 2\,598\,960 - 1\,712\,304 = 886\,656$.

\emph{3.} Choose the rank of the three-of-a-kind ($13$), its suits
($\binom43 = 4$), the rank of the pair ($12$ remaining), its suits
($\binom42 = 6$):
$13 \times 4 \times 12 \times 6 = 3744$.
\end{solution}

\begin{solution}{exo:g12:comb:6}
MATH has 4 distinct letters: $4! = 24$ anagrams.

BANANA has 6 letters: three A's, two N's, one B. Choose the positions of the
A's ($\binom63$), then of the N's among the rest ($\binom32$), the B takes
the last spot:
\[
\binom{6}{3}\binom{3}{2} = 20 \times 3 = 60 .
\]
(Equivalently $\frac{6!}{3!\,2!\,1!} = 60$.)
\end{solution}

\begin{solution}{exo:g12:comb:7}
\emph{Algebraically:}
\[
k\binom nk = \frac{k\,n!}{k!(n-k)!} = \frac{n!}{(k-1)!\,(n-k)!}
= n\,\frac{(n-1)!}{(k-1)!\bigl((n-1)-(k-1)\bigr)!} = n\binom{n-1}{k-1}.
\]
\emph{By double counting:} count pairs (committee of $k$, president in it).
Either choose the committee ($\binom nk$) then its president ($k$):
$k\binom nk$ pairs. Or choose the president first ($n$ options) then the
other $k-1$ members among the remaining $n-1$: $n\binom{n-1}{k-1}$ pairs.
\end{solution}

\begin{solution}{exo:g12:comb:8}
A path consists of exactly $m + n$ steps, of which $m$ are East and $n$ are
North; it is entirely determined by the set of instants (among the $m+n$) at
which one steps East. There are $\binom{m+n}{m}$ such choices.
\end{solution}

\begin{solution}{exo:g12:comb:9}
Split a set of $m + n$ people into a group $A$ of $m$ and a group $B$ of
$n$. A $k$-element subset contains some number $j$ of members of $A$
($0 \leq j \leq k$) and $k - j$ members of $B$; for fixed $j$ there are
$\binom mj \binom{n}{k-j}$ such subsets, and the addition principle over $j$
gives Vandermonde's identity.

With $m = n = k$:
\[
\binom{2n}{n} = \sum_{j=0}^n \binom nj \binom{n}{n-j}
= \sum_{j=0}^n \binom nj^{2},
\]
using the symmetry $\binom{n}{n-j} = \binom nj$.
\end{solution}

\begin{solution}{exo:g12:comb:10}
\emph{Via \cref{exo:g12:comb:7}:}
\[
\sum_{k=1}^n k\binom nk = \sum_{k=1}^n n\binom{n-1}{k-1}
= n\sum_{j=0}^{n-1}\binom{n-1}{j} = n\,2^{n-1},
\]
by \cref{cor:g12:comb:sums}. \emph{Via differentiation:} differentiating
$(1+x)^n = \sum_k \binom nk x^k$ gives
$n(1+x)^{n-1} = \sum_k k \binom nk x^{k-1}$; evaluate at $x = 1$.
\end{solution}

\begin{solution}{pb:g12:comb:1}
\textbf{1.} $26^2 \times 10^3 = 676\,000$ plates. BANANA: $6$
letters with A tripled and N doubled:
$\frac{6!}{3!\,2!} = 60$ anagrams.

\textbf{2.} $\binom{32}{5} = 201\,376$ hands;
$\binom42 \binom{28}{3} = 6 \times 3\,276 = 19\,656$ with
exactly two aces.

\textbf{3.} A path is a word with $4$ R's and $3$ U's: choose
the U positions: $\binom73 = 35$.

\textbf{4.} $(1+x)^4 = 1 + 4x + 6x^2 + 4x^3 + x^4$. At
$x = 1$: $\sum_k \binom nk = 2^n$; at $x = -1$:
$\sum_k (-1)^k \binom nk = 0$ --- row sums and alternating row
sums of \omterm{prop:g11:binom:pascal}{Pascal's triangle}.

\textbf{5.} Committees of $k$ people with a chair, from $n$:
choose the committee then its chair ($\binom nk \times k$), or
the chair then the other members
($n \times \binom{n-1}{k-1}$): equal. Summing over $k$: the
right side sums to $n \sum_j \binom{n-1}{j} = n\,2^{n-1}$.

\textbf{6.} A row of $10$ stars and $3$ bars encodes the order
(scoops of flavor 1 before the first bar, etc.); the row has
$13$ symbols and is determined by the bar positions:
$\binom{13}{3} = 286$ orders.

\textbf{7.} $12$ stars, $2$ bars: $\binom{14}{2} = 91$.

\textbf{8.} With $x', y', z' \geq 0$ and
$x' + y' + z' = 9$: $\binom{11}{2} = 55$.

\textbf{9.} A monomial $a^i b^j c^k$ with $i + j + k = 5$:
$\binom72 = 21$.

\textbf{10.} Formula: $3$ stars, $1$ bar: $\binom41 = 4$;
list: $(3,0)$, $(2,1)$, $(1,2)$, $(0,3)$: agreement.

\textbf{11.} ``Identical'' entered when an order was declared
to be nothing but the \emph{counts} per flavor --- the stars
carry no names. If the scoops are eaten in order, each of the
$10$ distinct positions picks a flavor freely:
$4^{10} = 1\,048\,576$ \omterm{def:g12:seq:sequence}{sequences} --- a different model and a
different world (\cref{met:g12:comb:model}: always ask
\emph{ordered? distinct? repetition allowed?}).

\textbf{12.} $D_1 = 0$; $D_2 = 1$ (swap); $D_3 = 2$ (the two
$3$-cycles); $D_4 = 9$.

\textbf{13.} Guest 1 gets hat $k \neq 1$: $n - 1$ choices. If
guest $k$ gets hat 1, the remaining $n - 2$ guests derange
their own hats: $D_{n-2}$ ways. If guest $k$ does \emph{not}
get hat 1, relabel hat 1 as guest $k$'s forbidden hat: the
$n - 1$ remaining guests derange: $D_{n-1}$ ways. Hence
$D_n = (n-1)(D_{n-1} + D_{n-2})$. Check:
$D_4 = 3(2 + 1) = 9$; and $D_5 = 4(9 + 2) = 44$.

\textbf{14.} Of the $3! = 6$ assignments, subtract those
fixing at least one hat: three fix a given hat ($2!$ each,
$3 \times 2 = 6$), overcounting the pairs ($3$ pairs, $1!$
each) which must return, and re-subtracting the identity
($1$): $D_3 = 6 - 6 + 3 - 1 = 2$, i.e.\
$3!\left(1 - 1 + \frac12 - \frac16\right) = 2$. In general
$D_n = n!\sum_{k=0}^{n} \frac{(-1)^k}{k!}$.

\textbf{15.} $\frac{D_5}{120} = \frac{44}{120} \approx
0.3667$, already close to $\frac1\eu \approx 0.3679$: the
alternating sum $1 - 1 + \frac{1}{2!} - \frac{1}{3!} + \dots$
marches to $\eu^{-1}$. A large party's hats derange about
$36.8\,\%$ of the time --- the lottery's and the secretary's
constant, third sighting.

\textbf{16.} $\P(\text{valid}) = \frac{D_{10}}{10!} \approx
0.368$. Each redraw succeeds with \omterm{def:g10:proba:distribution}{probability}
$\approx \frac1\eu$, so the expected number of draws is about
$\eu \approx 2.7$: budget three hat-passings.

\textbf{17.} Each handshake contributes $2$ to the total
degree count, so the sum of all guests' handshake-numbers is
even. A sum of \omterm{def:g10:numbers:sets}{integers} is even only if the number of odd
terms is even: odd-shakers come in even numbers. (At three
guests: possible handshake profiles never have exactly one or
three odd entries --- check the four possible \omterm{def:g10:functions:graph}{graphs}.)

\textbf{18.} $n = 1$: $1 = 1$. If
$1^3 + \dots + n^3 = \left(\frac{n(n+1)}{2}\right)^2$, adding
$(n+1)^3$:
\[
\frac{n^2(n+1)^2}{4} + (n+1)^3
= \frac{(n+1)^2\left(n^2 + 4n + 4\right)}{4}
= \left(\frac{(n+1)(n+2)}{2}\right)^{\!2} :
\]
heredity. For $n = 3$: $1 + 8 + 27 = 36 = 6^2$.

\textbf{19.} A path to $(n, n)$ makes $2n$ steps and crosses
the anti-diagonal $x + y = n$ at exactly one lattice point
$(j, n - j)$; the first half is a path with $j$ R's among $n$
steps ($\binom nj$ choices), the second half, read backwards,
likewise ($\binom nj$ again, by symmetry). Summing over the
crossing point: $\binom{2n}{n} = \sum_j \binom nj^2$ ---
Vandermonde's identity, drawn.

\textbf{20.} Multiply stages, add cases: plates and poker
hands. Encode: paths as RU-words, orders as stars and bars.
Count twice: committees-with-chair, handshakes, mid-cut paths.
Subtract and correct: the deranged hats, with $\frac1\eu$ as
the residue. Next stops: these counts under \omterm{def:g10:proba:distribution}{probability}'s
fractions, and the path-counting powers of the adjacency
matrices two chapters ahead.
\end{solution}

\section*{Chapter \ref{ch:g12:complex} --- Complex Numbers}

\begin{solution}{exo:g12:complex:1}
$(3-2\iu)(1+4\iu) = 3 + 12\iu - 2\iu - 8\iu^2 = 11 + 10\iu$.

$\dfrac{2+\iu}{1-\iu} = \dfrac{(2+\iu)(1+\iu)}{(1-\iu)(1+\iu)}
= \dfrac{1 + 3\iu}{2} = \dfrac12 + \dfrac32\iu$.

$\iu^{2027}$: since $\iu^4 = 1$ and $2027 = 4 \times 506 + 3$,
$\iu^{2027} = \iu^3 = -\iu$.
\end{solution}

\begin{solution}{exo:g12:complex:2}
\emph{(a)} $\Delta = 16 - 52 = -36$:
$z = \dfrac{4 \pm 6\iu}{2} = 2 \pm 3\iu$. \omterm{def:g12:complex:conjugate}{Conjugate} pair; product
$= 4 + 9 = 13 = \frac ca$. \emph{(b)} $\Delta = 1 - 4 = -3$:
$z = \dfrac{-1 \pm \iu\sqrt3}{2}$ (the cube roots of unity $j$ and
$\conj j$); product $= \frac{1 + 3}{4} = 1$.
\end{solution}

\begin{solution}{exo:g12:complex:3}
$\abs{z_1} = \sqrt2$, \omterm{def:g12:complex:argument}{argument} $\frac{3\pi}{4}$ (second quadrant):
$z_1 = \sqrt2\,\eu^{3\iu\pi/4}$.
$\abs{z_2} = 2$, \omterm{def:g12:complex:argument}{argument} $-\frac\pi6$:
$z_2 = 2\,\eu^{-\iu\pi/6}$. Hence
\[
\frac{z_1}{z_2} = \frac{\sqrt2}{2}\,\eu^{\iu\left(\frac{3\pi}{4} + \frac{\pi}{6}\right)}
= \frac{\sqrt2}{2}\,\eu^{11\iu\pi/12}.
\]
Algebraically,
\[
\frac{z_1}{z_2} = \frac{-1+\iu}{\sqrt3 - \iu}
= \frac{(-1+\iu)(\sqrt3+\iu)}{4}
= \frac{-\sqrt3 - 1 + \iu(\sqrt3 - 1)}{4}.
\]
Identifying \omterm{def:g12:complex:def}{real parts}:
$\frac{\sqrt2}{2}\cos\frac{11\pi}{12} = \frac{-\sqrt3-1}{4}$, so
\[
\cos\frac{11\pi}{12} = -\frac{\sqrt6 + \sqrt2}{4}.
\]
\end{solution}

\begin{solution}{exo:g12:complex:4}
\emph{(a)} $\abs{z - 2} = \abs{z - (-\iu)}$ means $M$ is equidistant from
$A(2, 0)$ and $B(0, -1)$: the \emph{perpendicular bisector} of $[AB]$.

\emph{(b)} Distance $3$ from the point $(1,1)$: the \emph{circle} of center
$1 + \iu$ and radius $3$.

\emph{(c)} By \cref{prop:g12:complex:geometry}(4), $\frac{z-1}{z+1}$ purely
imaginary means the \omterm{def:g11:vect:vector}{vectors} from $A(1,0)$ to $M$ and from $B(-1,0)$ to $M$
are \omterm{def:g11:scal:orthogonal}{orthogonal}: $M$ sees the segment $[AB]$ under a right angle. The set is
the \emph{circle of diameter $[AB]$} (the unit circle) \emph{minus the
points $A$ and $B$} themselves.
\end{solution}

\begin{solution}{exo:g12:complex:5}
De Moivre: $\cos3\theta + \iu\sin3\theta = (\cos\theta + \iu\sin\theta)^3$.
\omterm{def:g10:algebra:expand}{Expanding} by the binomial theorem with $c = \cos\theta$, $s = \sin\theta$:
\[
(c + \iu s)^3 = c^3 + 3\iu c^2 s - 3 c s^2 - \iu s^3
= (c^3 - 3cs^2) + \iu(3c^2 s - s^3).
\]
Identifying parts and using $s^2 = 1 - c^2$, $c^2 = 1 - s^2$:
\[
\cos3\theta = 4\cos^3\theta - 3\cos\theta,
\qquad
\sin3\theta = 3\sin\theta - 4\sin^3\theta .
\]
\end{solution}

\begin{solution}{exo:g12:complex:6}
\[
\cos^3\theta = \left(\frac{\eu^{\iu\theta} + \eu^{-\iu\theta}}{2}\right)^{\!3}
= \frac{\eu^{3\iu\theta} + 3\eu^{\iu\theta} + 3\eu^{-\iu\theta} + \eu^{-3\iu\theta}}{8}
= \frac{\cos3\theta + 3\cos\theta}{4}.
\]
Hence
\[
\int_0^{\pi/2}\cos^3\theta\,\dd\theta
= \frac14\left[\frac{\sin3\theta}{3} + 3\sin\theta\right]_0^{\pi/2}
= \frac14\left(-\frac13 + 3\right) = \frac{2}{3}.
\]
(Consistent with $W_3 = \frac23$ in \cref{exo:g12:integ:9}, since
$\int_0^{\pi/2}\cos^3 = \int_0^{\pi/2}\sin^3$ by the symmetry
$\theta \mapsto \frac\pi2 - \theta$.)
\end{solution}

\begin{solution}{exo:g12:complex:7}
$c = a + \eu^{\pm\iu\pi/3}(b - a)$ with $b - a = 2 + \iu$ and
$\eu^{\pm\iu\pi/3} = \frac12 \pm \iu\frac{\sqrt3}{2}$:
\[
\eu^{\iu\pi/3}(2+\iu) = \left(\tfrac12 + \iu\tfrac{\sqrt3}{2}\right)(2+\iu)
= \left(1 - \tfrac{\sqrt3}{2}\right) + \iu\left(\tfrac12 + \sqrt3\right),
\]
so $c_1 = \left(2 - \frac{\sqrt3}{2}\right) +
\iu\left(\frac32 + \sqrt3\right)$, and similarly with the angle
$-\frac\pi3$:
$c_2 = \left(2 + \frac{\sqrt3}{2}\right) + \iu\left(\frac32 - \sqrt3\right)$.
\end{solution}

\begin{solution}{exo:g12:complex:8}
$-4 = 4\,\eu^{\iu\pi}$, so the solutions of $z^4 = -4$ are
\[
z_k = \sqrt2\,\eu^{\iu\left(\frac\pi4 + \frac{k\pi}{2}\right)},
\quad k = 0, 1, 2, 3,
\]
\emph{i.e.} $1 + \iu$, $-1 + \iu$, $-1 - \iu$, $1 - \iu$: the vertices of a
square of circumradius $\sqrt 2$. Grouping \omterm{def:g12:complex:conjugate}{conjugate} \omterm{def:g11:quad:discriminant}{roots},
\[
z^4 + 4 = \bigl(z - (1{+}\iu)\bigr)\bigl(z - (1{-}\iu)\bigr)
\bigl(z + (1{-}\iu)\bigr)\bigl(z + (1{+}\iu)\bigr)
= (z^2 - 2z + 2)(z^2 + 2z + 2).
\]
\end{solution}

\begin{solution}{exo:g12:complex:9}
With $q = \eu^{\iu\theta} \neq 1$:
\[
\sum_{k=0}^{n} \eu^{\iu k\theta} = \frac{\eu^{\iu(n+1)\theta} - 1}{\eu^{\iu\theta} - 1}
= \frac{\eu^{\iu(n+1)\theta/2}}{\eu^{\iu\theta/2}}
\cdot\frac{\eu^{\iu(n+1)\theta/2} - \eu^{-\iu(n+1)\theta/2}}{\eu^{\iu\theta/2} - \eu^{-\iu\theta/2}}
= \eu^{\iu n\theta/2}\,
\frac{\sin\frac{(n+1)\theta}{2}}{\sin\frac{\theta}{2}},
\]
using $\eu^{\iu\alpha} - \eu^{-\iu\alpha} = 2\iu\sin\alpha$ (the
\emph{half-angle \omterm{def:g10:algebra:expand}{factoring}}). Taking \omterm{def:g12:complex:def}{real parts}:
\[
S = \frac{\sin\frac{(n+1)\theta}{2}}{\sin\frac{\theta}{2}}\,
\cos\frac{n\theta}{2}.
\]
\end{solution}

\begin{solution}{exo:g12:complex:10}
\emph{1.} $\omega$ is a $5$-th \omterm{def:g11:quad:discriminant}{root} of unity different from $1$, so the sum
of all five \omterm{def:g11:quad:discriminant}{roots} vanishes (\cref{thm:g12:complex:rootsofunity}), which reads
$1 + \omega + \omega^2 + \omega^3 + \omega^4 = 0$.

\emph{2.} $u + v = \omega + \omega^2 + \omega^3 + \omega^4 = -1$ by point 1.
For the product, using $\omega^5 = 1$:
\[
uv = (\omega + \omega^4)(\omega^2 + \omega^3)
= \omega^3 + \omega^4 + \omega^6 + \omega^7
= \omega^3 + \omega^4 + \omega + \omega^2 = -1 .
\]
So $u, v$ are the \omterm{def:g11:quad:discriminant}{roots} of $X^2 + X - 1 = 0$:
$\left\{\frac{-1+\sqrt5}{2}, \frac{-1-\sqrt5}{2}\right\}$. Now
$u = \omega + \conj\omega = 2\cos\frac{2\pi}{5} > 0$ (since
$\frac{2\pi}{5} < \frac{\pi}{2}$), so $u = \frac{-1 + \sqrt5}{2}$.

\emph{3.} $\cos\frac{2\pi}{5} = \frac u2 = \frac{\sqrt5 - 1}{4}$.
\end{solution}

\section*{Chapter \ref{ch:g12:arith} --- Arithmetic}

\begin{solution}{exo:g12:arith:1}
$17 \times 119 = 2023$, so $2026 = 17 \times 119 + 3$: quotient $119$,
remainder $3$. For $-2026$: $-2026 = 17\times(-120) + 14$ (indeed
$17 \times 120 = 2040$ and $2040 - 2026 = 14$): quotient $-120$, remainder
$14$ (the remainder must lie in $\intco{0}{17}$, so it is \emph{not} $-3$).
\end{solution}

\begin{solution}{exo:g12:arith:2}
Modulo $10$: $7^2 = 49 \equiv 9 \equiv -1$. Hence
$7^{100} = \left(7^2\right)^{50} \equiv (-1)^{50} = 1 \pmod{10}$: the last
digit of $7^{100}$ is $1$.
\end{solution}

\begin{solution}{exo:g12:arith:3}
Euclid: $1071 = 2\times462 + 147$; $462 = 3\times147 + 21$;
$147 = 7\times21 + 0$. So $\gcd = 21$.

Back-substitution: $21 = 462 - 3\times147 = 462 - 3(1071 - 2\times462)
= 7\times462 - 3\times1071$. Thus $u = -3$, $v = 7$:
$1071\times(-3) + 462\times7 = 21$.
\end{solution}

\begin{solution}{exo:g12:arith:4}
Every \omterm{def:g10:numbers:sets}{integer} is $\equiv 0, 1, 2$ or $3 \pmod 4$, and squaring:
$0^2 \equiv 0$, $1^2 \equiv 1$, $2^2 = 4 \equiv 0$, $3^2 = 9 \equiv 1$. So
$n^2 \equiv 0$ or $1 \pmod 4$. A sum of two squares is then congruent to
$0 + 0$, $0 + 1$ or $1 + 1$, \emph{i.e.} to $0$, $1$ or $2 \pmod 4$ ---
never to $3$.
\end{solution}

\begin{solution}{exo:g12:arith:5}
\omterm{def:g12:arith:divides}{Divisibility} by $2$: among $n$ and $n + 1$, one is even. \omterm{def:g12:arith:divides}{Divisibility} by
$3$: if $n \equiv 0$, then $3 \mid n$; if $n \equiv 1 \pmod 3$, then
$2n + 1 \equiv 3 \equiv 0$; if $n \equiv 2$, then $n + 1 \equiv 0$. In all
cases $3$ \omterm{def:g12:arith:divides}{divides} the product. Since $\gcd(2,3) = 1$,
\cref{cor:g12:arith:coprimeprod} gives $6 \mid n(n+1)(2n+1)$.
(This also re-proves that $\frac{n(n+1)(2n+1)}{6}$, the sum of squares of
\cref{exo:g12:seq:1}, is an \omterm{def:g10:numbers:sets}{integer}.)
\end{solution}

\begin{solution}{exo:g12:arith:6}
We look for the inverse of $5$ modulo $11$: testing (or Bézout),
$5 \times 9 = 45 = 44 + 1 \equiv 1 \pmod{11}$. Multiplying the \omterm{def:g12:arith:congruence}{congruence} by
$9$:
\[
x \equiv 9 \times 3 = 27 \equiv 5 \pmod{11}.
\]
The solutions are the \omterm{def:g10:numbers:sets}{integers} $x = 5 + 11k$, $k \in \Z$. (Check:
$5\times5 = 25 \equiv 3 \pmod{11}$.)
\end{solution}

\begin{solution}{exo:g12:arith:7}
$\gcd(17, 40) = 1$, so solutions exist. Euclid:
$40 = 2\times17 + 6$; $17 = 2\times6 + 5$; $6 = 5 + 1$. Back-substituting:
$1 = 6 - 5 = 6 - (17 - 2\times6) = 3\times6 - 17
= 3(40 - 2\times17) - 17 = 3\times40 - 7\times17$. Hence
$17\times(-7) - 40\times(-3) = 1$: the particular solution
$(x_0, y_0) = (-7, -3)$.

General solution of $17x - 40y = 1$: subtracting the particular relation,
$17(x + 7) = 40(y + 3)$; since $\gcd(17, 40) = 1$, Gauss gives
$40 \mid x + 7$, so $x = -7 + 40k$ and then $y = -3 + 17k$, $k \in \Z$
(all of which check).

For $17x - 40y = 6$, multiply the particular solution by $6$:
$(x_1, y_1) = (-42, -18)$, and the same reasoning gives
\[
x = -42 + 40k, \qquad y = -18 + 17k, \qquad k \in \Z .
\]
(E.g.\ $k = 2$: $x = 38$, $y = 16$; indeed $17\times38 - 40\times16
= 646 - 640 = 6$.)
\end{solution}

\begin{solution}{exo:g12:arith:8}
Suppose $\sqrt2 = \frac ab$ with $a, b \in \N^*$; then $a^2 = 2b^2$. In the
\omterm{def:g12:arith:prime}{prime} factorization of a square, every exponent is even; so the exponent of
$2$ in $a^2$ is even, while in $2b^2$ it is odd (one more than an even
number). Two factorizations of the same \omterm{def:g10:numbers:sets}{integer} with different exponents of
$2$ contradict the uniqueness in \cref{thm:g12:arith:fta}. Hence no such
fraction exists: $\sqrt2 \notin \Q$.
\end{solution}

\begin{solution}{exo:g12:arith:9}
\emph{1.} From $k\binom pk = p \binom{p-1}{k-1}$, $p$ \omterm{def:g12:arith:divides}{divides}
$k\binom pk$. For $1 \leq k \leq p-1$, $p \nmid k$ and $p$ \omterm{def:g12:arith:prime}{prime} give
$\gcd(p, k) = 1$, so Gauss's lemma yields $p \mid \binom pk$.

\emph{2.} Induction on $a$. For $a = 0$: $0^p \equiv 0$. Assume
$a^p \equiv a \pmod p$. By the binomial theorem,
\[
(a+1)^p = \sum_{k=0}^{p} \binom pk a^k
\equiv a^p + 1 \pmod p,
\]
all middle terms vanishing modulo $p$ by point 1. By the induction
hypothesis, $(a+1)^p \equiv a + 1 \pmod p$. This proves
$a^p \equiv a$ for all $a \in \N$, and the case $a < 0$ follows by writing
$a \equiv a + kp$ for a suitable positive representative.
\end{solution}

\begin{solution}{exo:g12:arith:10}
$n \equiv 2 \pmod 3$ and $n \equiv 3 \pmod 5$: write $n = 2 + 3s$; then
$2 + 3s \equiv 3 \pmod 5$, \emph{i.e.} $3s \equiv 1 \pmod 5$. The inverse of
$3$ modulo $5$ is $2$ ($3\times2 = 6 \equiv 1$), so $s \equiv 2 \pmod 5$,
say $s = 2 + 5t$, and $n = 8 + 15t$: the first two conditions \omterm{def:g11:stat:mean}{mean}
$n \equiv 8 \pmod{15}$.

Adding $n \equiv 2 \pmod 7$: $8 + 15t \equiv 2 \pmod 7$, and
$15 \equiv 1 \pmod 7$, so $t \equiv -6 \equiv 1 \pmod 7$, say
$t = 1 + 7u$. Hence $n = 23 + 105u$:
\[
n \equiv 23 \pmod{105}.
\]
(Check: $23 = 3\times7 + 2 = 5\times4 + 3 = 7\times3 + 2$.)
\end{solution}

\section*{Chapter \ref{ch:g12:matrix} --- Matrices and Graphs}

\begin{solution}{exo:g12:matrix:1}
\[
A + B = \begin{pmatrix} 3 & 2\\ 1 & 2\end{pmatrix}, \quad
AB = \begin{pmatrix} 4 & 2\\ 1 & 1\end{pmatrix}, \quad
BA = \begin{pmatrix} 2 & 4\\ 1 & 3\end{pmatrix}, \quad
A^2 = \begin{pmatrix} 1 & 4\\ 0 & 1\end{pmatrix}.
\]
Note $AB \neq BA$.
\end{solution}

\begin{solution}{exo:g12:matrix:2}
$\det A = 6 - 5 = 1 \neq 0$:
$A^{-1} = \begin{pmatrix} 3 & -5\\ -1 & 2 \end{pmatrix}$.
$\det B = 12 - 12 = 0$: $B$ is not invertible.
\end{solution}

\begin{solution}{exo:g12:matrix:3}
The system is $AX = Y$ with $A$ as in \cref{exo:g12:matrix:2} and
$Y = \begin{pmatrix} 1\\ 2\end{pmatrix}$:
\[
X = A^{-1}Y = \begin{pmatrix} 3 & -5\\ -1 & 2\end{pmatrix}
\begin{pmatrix} 1\\ 2\end{pmatrix}
= \begin{pmatrix} -7\\ 3\end{pmatrix}:
\qquad x = -7,\ y = 3 .
\]
\end{solution}

\begin{solution}{exo:g12:matrix:4}
\emph{1.} $N^2 = \begin{pmatrix} 0&1\\0&0\end{pmatrix}
\begin{pmatrix} 0&1\\0&0\end{pmatrix} = 0$, and $I$ commutes with every
matrix.

\emph{2.} Since the two terms commute, the binomial theorem applies and all
terms containing $N^2$ vanish:
\[
A^k = (2I + N)^k = 2^k I + k\,2^{k-1} N
= \begin{pmatrix} 2^k & k\,2^{k-1}\\ 0 & 2^k \end{pmatrix}.
\]
Check for $k = 2$:
$A^2 = \begin{pmatrix} 2&1\\0&2\end{pmatrix}^2
= \begin{pmatrix} 4&4\\0&4\end{pmatrix}$, and the formula gives
$2^2 = 4$, $2 \times 2 = 4$. \checkmark
\end{solution}

\begin{solution}{exo:g12:matrix:5}
\emph{Induction.} For $n = 1$:
$A^1 = \begin{pmatrix} 0&1\\1&1\end{pmatrix}
= \begin{pmatrix} F_0 & F_1\\ F_1 & F_2\end{pmatrix}$. Assume the formula
for $n$; then
\[
A^{n+1} = A^n A
= \begin{pmatrix} F_{n-1} & F_n\\ F_n & F_{n+1}\end{pmatrix}
\begin{pmatrix} 0 & 1\\ 1 & 1\end{pmatrix}
= \begin{pmatrix} F_n & F_{n-1} + F_n\\ F_{n+1} & F_n + F_{n+1}\end{pmatrix}
= \begin{pmatrix} F_n & F_{n+1}\\ F_{n+1} & F_{n+2}\end{pmatrix}.
\]
\emph{Identity.} For $2\times2$ matrices, expanding shows
$\det(MN) = \det M \det N$; hence
$\det(A^n) = (\det A)^n = (-1)^n$, and
$\det A^n = F_{n-1}F_{n+1} - F_n^2$. (This is \emph{Cassini's identity}.)
\end{solution}

\begin{solution}{exo:g12:matrix:6}
\emph{1.} Ordering vertices $1, 2, 3$:
\[
M = \begin{pmatrix} 0&1&1\\ 0&0&1\\ 1&0&0\end{pmatrix}, \quad
M^2 = \begin{pmatrix} 1&0&1\\ 1&0&0\\ 0&1&1\end{pmatrix}, \quad
M^3 = \begin{pmatrix} 1&1&1\\ 1&0&1\\ 1&0&1 \end{pmatrix}.
\]

\emph{2.} $\bigl(M^3\bigr)_{11} = 1$: exactly one closed walk of length $3$
at vertex $1$, namely $1 \to 2 \to 3 \to 1$. (The walk
$1 \to 3 \to 1$ has length $2$ only, and $1 \to 3$ then $3\to1$ then
$1\to3$ ends at $3$.)
\end{solution}

\begin{solution}{exo:g12:matrix:7}
\emph{1.} $a_{n+1} = 0.8a_n + 0.3b_n$, $b_{n+1} = 0.2a_n + 0.7b_n$:
$M = \begin{pmatrix} 0.8 & 0.3\\ 0.2 & 0.7\end{pmatrix}$.

\emph{2.} $MX = X$ gives $0.8a + 0.3b = a$, \emph{i.e.} $0.3b = 0.2a$, so
$b = \frac23 a$; with $a + b = 1$: $a = 0.6$, $b = 0.4$.

\emph{3.} Using $b_n = 1 - a_n$:
$a_{n+1} = 0.8a_n + 0.3(1 - a_n) = 0.5a_n + 0.3$, so
\[
c_{n+1} = a_{n+1} - 0.6 = 0.5a_n + 0.3 - 0.6 = 0.5(a_n - 0.6) = 0.5\,c_n .
\]
Hence $c_n = 0.5^n c_0 \to 0$: $a_n \to 0.6$ and $b_n \to 0.4$, whatever the
initial distribution.
\end{solution}

\begin{solution}{exo:g12:matrix:8}
\emph{1.} $\det P = -2$, so
$P^{-1} = -\frac12\begin{pmatrix} -1 & -1\\ -1 & 1\end{pmatrix}
= \frac12\begin{pmatrix} 1 & 1\\ 1 & -1\end{pmatrix}$. Then
\[
AP = \begin{pmatrix} 4 & 2\\ 4 & -2 \end{pmatrix}, \qquad
D = P^{-1}AP = \frac12\begin{pmatrix} 1&1\\1&-1\end{pmatrix}
\begin{pmatrix} 4&2\\4&-2\end{pmatrix}
= \begin{pmatrix} 4 & 0\\ 0 & 2\end{pmatrix}.
\]

\emph{2.} From $A = PDP^{-1}$, an immediate induction gives
$A^n = PD^nP^{-1}$ with
$D^n = \begin{pmatrix} 4^n & 0\\ 0 & 2^n\end{pmatrix}$, so
\[
A^n = P D^n P^{-1}
= \begin{pmatrix} 4^n & 2^n\\ 4^n & -2^n\end{pmatrix}\cdot
\frac12\begin{pmatrix} 1&1\\1&-1\end{pmatrix}
= \frac12\begin{pmatrix} 4^n + 2^n & 4^n - 2^n\\
4^n - 2^n & 4^n + 2^n\end{pmatrix}.
\]
Applying $A^n$ to $X_0 = \begin{pmatrix} u_0\\v_0\end{pmatrix}$ reproduces
exactly the formulas of \cref{ex:g12:matrix:coupled}.
\end{solution}

\section*{Chapter \ref{ch:g12:space} --- Vectors, Lines and Planes in Space}

\begin{solution}{exo:g12:space:1}
$\vect{AB}(2, 1, -1)$ and $\vect{AC}(1, -1, 1)$ are not proportional
($\frac21 \neq \frac{1}{-1}$), so the points are not aligned. Line $(AB)$:
\[
(x, y, z) = (1 + 2t,\; t,\; 2 - t), \qquad t \in \R .
\]
\end{solution}

\begin{solution}{exo:g12:space:2}
$2(x - 1) - (y - 1) + 3(z - 0) = 0$, \emph{i.e.}
\[
2x - y + 3z - 1 = 0 .
\]
For $B(0,2,1)$: $0 - 2 + 3 - 1 = 0$: yes, $B$ lies on the plane.
\end{solution}

\begin{solution}{exo:g12:space:3}
$\vect{AB}(1,1,0)$, $\vect{AC}(1,0,1)$:
\[
\cos\widehat{A}
= \frac{\vect{AB}\cdot\vect{AC}}{\norm{\vect{AB}}\,\norm{\vect{AC}}}
= \frac{1}{\sqrt2\,\sqrt2} = \frac12,
\]
so $\widehat A = 60^\circ$.
\end{solution}

\begin{solution}{exo:g12:space:4}
Points of $\mathcal D$: $(1 + t,\, 2 - t,\, 2t)$. Substituting into the
\omterm{def:g10:algebra:equation}{equation} of $\mathcal P$:
\[
2(1+t) + (2 - t) - 2t + 1 = 5 - t = 0 \iff t = 5 .
\]
Unique \omterm{def:g10:numbers:interunion}{intersection} point: $(6, -3, 10)$.
\end{solution}

\begin{solution}{exo:g12:space:5}
The \omterm{def:g12:space:normal}{normal vectors} $(1, -1, 2)$ and $(2, 1, 1)$ are not \omterm{def:g12:space:collinear}{collinear}, so the
planes intersect in a line. Adding the two \omterm{def:g10:algebra:equation}{equations}: $3x + 3z = 5$, so
$x = \frac53 - z$; then the first \omterm{def:g10:algebra:equation}{equation} gives
$y = x + 2z - 1 = \frac23 + z$. With $z = t$:
\[
(x, y, z) = \left(\frac53 - t,\; \frac23 + t,\; t\right), \qquad t \in \R,
\]
a line with \omterm{def:g11:vect:direction}{direction vector} $(-1, 1, 1)$. (Both \omterm{def:g10:algebra:equation}{equations} check
identically.)
\end{solution}

\begin{solution}{exo:g12:space:6}
Directions $\vec u_1(1, 2, -1)$ and $\vec u_2(1, 1, 2)$ are not \omterm{def:g12:space:collinear}{collinear},
so the lines are not parallel. Intersecting would require
\[
1 + t = 2 + s, \qquad 2t = 1 + s, \qquad 3 - t = 1 + 2s .
\]
The first two give $t = 1 + s$ and $2(1+s) = 1 + s$, so $s = -1$, $t = 0$;
but then the third reads $3 = -1$, absurd. No common point and not
parallel: the lines are skew.
\end{solution}

\begin{solution}{exo:g12:space:7}
\emph{1.} $\operatorname{dist} = \dfrac{\abs{2\cdot3 - 1 + 2\cdot1 - 3}}
{\sqrt{4 + 1 + 4}} = \dfrac{4}{3}$.

\emph{2.} $H = M_0 + t\,\vec n$ with $\vec n(2, -1, 2)$, choosing $t$ so
that $H \in \mathcal P$:
\[
2(3 + 2t) - (1 - t) + 2(1 + 2t) - 3 = 4 + 9t = 0 \iff t = -\frac49 .
\]
Hence $H\left(\frac{19}{9}, \frac{13}{9}, \frac19\right)$, and
\[
M_0H = \norm{t\,\vec n} = \frac49 \times 3 = \frac43,
\]
matching the distance formula.
\end{solution}

\begin{solution}{exo:g12:space:8}
\omterm{def:g10:coordgeom:system}{Coordinates}: $G(1,1,1)$, and the plane $(BDE)$ passes through $B(1,0,0)$,
$D(0,1,0)$, $E(0,0,1)$.

\emph{1.} The plane $(BDE)$ has \omterm{def:g10:algebra:equation}{equation} $x + y + z = 1$ (satisfied by the
three points, and they are not aligned), so $\vec n(1,1,1)$ is normal to
it. Since $\vect{AG} = (1,1,1) = \vec n$, the diagonal $(AG)$ is \omterm{def:g11:scal:orthogonal}{orthogonal}
to the plane $(BDE)$.

\emph{2.} Distance from $A(0,0,0)$:
$\frac{\abs{0 + 0 + 0 - 1}}{\sqrt3} = \frac{1}{\sqrt3} = \frac{\sqrt3}{3}$.
The line $(AG)$ is $(t, t, t)$; it meets the plane when $3t = 1$: at the
point $\left(\frac13, \frac13, \frac13\right)$, the centroid of the
triangle $BDE$.
\end{solution}

\begin{solution}{exo:g12:space:9}
\emph{1.} $\vect{AB}(1, 0, -1)$, $\vect{AC}(-1, 1, 0)$,
$\vect{AD}(1, 1, 1)$. If $\vect{AD} = a\vect{AB} + b\vect{AC}$, the
\omterm{def:g10:coordgeom:system}{coordinates} give $a - b = 1$, $b = 1$, $-a = 1$: the first two force
$a = 2$, contradicting $a = -1$. Not \omterm{def:g12:space:collinear}{coplanar}.

\emph{2.} $\vect{BC}(-2, 1, 1)$, $\vect{BD}(0, 1, 2)$. Solving
$\vec n \cdot \vect{BC} = \vec n \cdot \vect{BD} = 0$: from
$b + 2c = 0$, $b = -2c$; then $-2a - 2c + c = 0$ gives $c = -2a$. With
$a = 1$: $\vec n(1, 4, -2)$. Plane through $B(2,1,0)$:
\[
x + 4y - 2z - 6 = 0 .
\]

\emph{3.} Distance from $A(1,1,1)$:
$\dfrac{\abs{1 + 4 - 2 - 6}}{\sqrt{21}} = \dfrac{3}{\sqrt{21}}$.
Area of $BCD$: $\norm{\vect{BC}}^2 = 6$, $\norm{\vect{BD}}^2 = 5$,
$\vect{BC}\cdot\vect{BD} = 3$, so
\[
\text{Area} = \frac12\sqrt{6 \times 5 - 9} = \frac{\sqrt{21}}{2}.
\]

\emph{4.}
\[
V = \frac13 \times \frac{\sqrt{21}}{2} \times \frac{3}{\sqrt{21}}
= \frac12 .
\]
\end{solution}

\section*{Chapter \ref{ch:g12:condprob} --- Conditional Probability and Independence}

\begin{solution}{exo:g12:condprob:1}
There are $12$ face cards, of which $4$ are kings:
$\pcond{\text{face}}{\text{king}} = \frac{4}{12} = \frac13$.

\omterm{def:g12:condprob:indep}{Independence}: $\P(\text{king} \cap \text{heart}) = \frac{1}{52}$ (the king
of hearts), and
$\P(\text{king})\,\P(\text{heart}) = \frac{4}{52}\times\frac{13}{52}
= \frac{1}{52}$. Equal: the \omterm{def:g11:prob:model}{events} are \omterm{def:g12:condprob:indep}{independent}.
\end{solution}

\begin{solution}{exo:g12:condprob:2}
\emph{1.} First branch: red $\frac58$, blue $\frac38$; second branches
(without replacement): after red, red $\frac47$ / blue $\frac37$; after
blue, red $\frac57$ / blue $\frac27$.

\emph{2.} $\P(RR) = \frac58 \times \frac47 = \frac{5}{14}$. By total
\omterm{def:g10:proba:distribution}{probability},
\[
\P(\text{2nd red}) = \frac58\cdot\frac47 + \frac38\cdot\frac57
= \frac{20 + 15}{56} = \frac58 .
\]
(The same as the first draw --- by symmetry, the second ball is a uniformly
random ball of the urn.)
\end{solution}

\begin{solution}{exo:g12:condprob:3}
$\P(A) = \frac{6}{36} = \frac16$ (six pairs sum to $7$),
$\P(B) = \frac16$, and $A \cap B = \{(3,4)\}$ has \omterm{def:g10:proba:distribution}{probability}
$\frac{1}{36} = \P(A)\P(B)$: $A$ and $B$ are \omterm{def:g12:condprob:indep}{independent}. (The sum $7$ is
special: whatever the first die, exactly one value of the second gives it.)

$\P(C) = \frac{5}{36}$ and $B \cap C = \{(3,3)\}$:
$\P(B \cap C) = \frac{1}{36} \neq \frac16 \times \frac{5}{36}
= \frac{5}{216}$. Not \omterm{def:g12:condprob:indep}{independent}.
\end{solution}

\begin{solution}{exo:g12:condprob:4}
\emph{1.} Total \omterm{def:g10:proba:distribution}{probability} with the partition by machine:
\[
\P(D) = 0.5\times0.01 + 0.3\times0.02 + 0.2\times0.04
= 0.005 + 0.006 + 0.008 = 0.019 = 1.9\% .
\]

\emph{2.} Bayes:
$\pcond{D}{M_3} = \dfrac{0.008}{0.019} = \dfrac{8}{19} \approx 0.42$. The
machine producing only a fifth of the output accounts for over $40\%$ of
the defects.
\end{solution}

\begin{solution}{exo:g12:condprob:5}
The condition is
\[
\frac{0.99\,p}{0.99\,p + 0.05\,(1-p)} \geq 0.9 .
\]
The denominator is positive, so this reads
$0.99p \geq 0.891p + 0.045 - 0.045p$, \emph{i.e.} $0.144\,p \geq 0.045$,
\emph{i.e.}
\[
p \geq \frac{0.045}{0.144} = 0.3125 .
\]
The test alone reaches $90\%$ certainty only if the disease already affects
more than $31\%$ of the tested population --- which is why mass screening
of rare diseases requires confirmation tests.
\end{solution}

\begin{solution}{exo:g12:condprob:6}
\emph{1.} Three \omterm{def:g12:seq:sequence}{sequences} realize $E$ (HHT, HTH, THH), each of \omterm{def:g10:proba:distribution}{probability}
$p^2(1-p)$ by \omterm{def:g12:condprob:indep}{independence}: $\P(E) = 3p^2(1 - p)$.

\emph{2.} $f(p) = 3p^2 - 3p^3$ has $f'(p) = 6p - 9p^2 = 3p(2 - 3p)$,
positive then negative on $\intoo{0}{1}$: \omterm{def:g10:functions:extrema}{maximum} at $p = \frac23$, where
$\P(E) = 3\cdot\frac49\cdot\frac13 = \frac49$.
\end{solution}

\begin{solution}{exo:g12:condprob:7}
Say you picked door 1 and the host opened door 3 (\omterm{def:g11:prob:model}{event} $H_3$). With
$B_i$ = ``prize behind door $i$'', $\P(B_i) = \frac13$ and
\[
\pcond{B_1}{H_3} = \frac12, \qquad
\pcond{B_2}{H_3} = 1, \qquad
\pcond{B_3}{H_3} = 0
\]
(if the prize is behind your door, the host chooses between doors 2 and 3
at random; if it is behind door 2, he is forced to open 3). Bayes:
\[
\pcond{H_3}{B_1}
= \frac{\frac13\cdot\frac12}{\frac13\cdot\frac12 + \frac13\cdot1 + 0}
= \frac{1/6}{1/2} = \frac13,
\qquad
\pcond{H_3}{B_2} = \frac{\frac13}{\frac12} = \frac23 .
\]
Staying wins with \omterm{def:g10:proba:distribution}{probability} $\frac13$, switching with \omterm{def:g10:proba:distribution}{probability}
$\frac23$.
\end{solution}

\begin{solution}{exo:g12:condprob:8}
\emph{1.} The majority is wrong when at least two of the three copies are
flipped:
\[
3\varepsilon^2(1-\varepsilon) + \varepsilon^3
= 3(0.01)(0.9) + 0.001 = 0.028,
\]
much smaller than $\varepsilon = 0.1$: repetition coding works.

\emph{2.} If $1$ was sent (as $111$), receiving $101$ requires exactly one
flip: \omterm{def:g10:proba:distribution}{probability} $\varepsilon(1-\varepsilon)^2 = 0.081$. If $0$ was sent
($000$), two flips are needed: $\varepsilon^2(1-\varepsilon) = 0.009$.
Bayes with equal priors:
\[
\pcond{101}{\,\text{sent }1} = \frac{0.081}{0.081 + 0.009} = 0.9 .
\]
The majority decoding ($101 \mapsto 1$) is indeed the more probable guess.
\end{solution}

\begin{solution}{pb:g12:condprob:1}
\textbf{1.} $\P(\text{red}) = \frac12 \cdot \frac23 + \frac12
\cdot \frac14 = \frac13 + \frac18 = \frac{11}{24}$.

\textbf{2.} $\P(\text{I} \mid \text{red}) =
\dfrac{\frac12 \cdot \frac23}{\frac{11}{24}} =
\dfrac{1/3}{11/24} = \dfrac{8}{11}$.

\textbf{3.} $\P(\text{even} \cap \leq 4) = \P(\{2, 4\}) =
\frac13$ and $\P(\text{even})\P(\leq 4) = \frac12 \cdot
\frac23 = \frac13$: \omterm{def:g12:condprob:indep}{independent} (unexpectedly!). For
$\leq 3$: $\P(\{2\}) = \frac16$ against
$\frac12 \cdot \frac12 = \frac14$: dependent.

\textbf{4.} The \omterm{def:g10:proba:operations}{complements} of
\omterm{def:g12:condprob:indep}{independent}
\omterm{def:g11:prob:model}{events} are themselves
\omterm{def:g12:condprob:indep}{independent}, so
\[
\P(\text{none}) = 0.1 \times 0.1 = 0.01,
\qquad
\P(\text{at least one}) = 0.99 .
\]

\textbf{5.} $\frac{4}{52} \times \frac{3}{51} =
\frac{1}{221}$.

\textbf{6.} Sick: $100$; healthy: $99\,900$. Positives:
$99$ true (99\,\% of $100$) and $4\,995$ false (5\,\% of
$99\,900$): $5\,094$ positives in all.

\textbf{7.} $\P(\text{sick} \mid +) = \frac{99}{5\,094}
\approx 1.9\,\%$: of the people the test flags, fewer than
one in fifty is actually sick.

\textbf{8.} $\dfrac{0.001 \times 0.99}{0.001 \times 0.99 +
0.999 \times 0.05} = \dfrac{0.00099}{0.05094} \approx 0.019$:
the counts and the formula agree.

\textbf{9.} Confused: $\P(+ \mid \text{sick}) = 99\,\%$
(the test's advertised quality) with
$\P(\text{sick} \mid +)$ (the patient's actual question). With
a base rate of $\frac{1}{1000}$, the sick are so rare that
even a $5\,\%$ leak from the huge healthy majority swamps the
true positives fifty to one.

\textbf{10.} $\dfrac{0.019 \times 0.99}{0.019 \times 0.99 +
0.981 \times 0.05} \approx 28\,\%$. Each \omterm{def:g12:condprob:indep}{independent} positive
multiplies the odds by the same factor: evidence accumulates,
and a third positive would push past $90\,\%$ --- which is
why confirmatory tests exist.

\textbf{11.} $\dfrac{0.05 \times 0.99}{0.05 \times 0.99 +
0.95 \times 0.05} \approx 51\,\%$: on a high-risk group the
same test is informative. Moral: screen where the base rate
is substantial; mass-screening a rare condition manufactures
false alarms by the thousand.

\textbf{12.} Expected innocent matches:
$10^7 \times 10^{-6} = 10$. Among the roughly $11$ matching
people ($10$ innocent, $1$ guilty --- if the culprit is in
the city), the suspect found \emph{by the match alone} is
innocent with \omterm{def:g10:proba:distribution}{probability} about $\frac{10}{11} \approx
91\,\%$.

\textbf{13.} The prosecutor quoted
$\P(\text{match} \mid \text{innocent}) = 10^{-6}$; the court
needed $\P(\text{innocent} \mid \text{match})$ --- question
9's confusion under oath. (And in the real case that
popularized these errors, a second blunder compounded it:
see question 14.)

\textbf{14.} Naive squaring: $\left(\frac{1}{8500}\right)^2
\approx \frac{1}{72 \text{ million}}$. With the second
tragedy five times likelier given the first:
$\frac{1}{8500} \times \frac{5}{8500} \approx
\frac{1}{14.5 \text{ million}}$ --- five times larger, and
still ignoring that the \emph{alternative} hypothesis must be
weighed by Bayes too. Lesson: multiplying probabilities
requires \omterm{def:g12:condprob:indep}{independence}, and shared causes (genetics,
environment) destroy it.

\textbf{15.} Bayes multiplies the likelihood ratios of
\emph{all} the evidence into one posterior: the match, the
alibi, the fibers, each shifting the odds. The juror's
sentence: \emph{ask how probable the evidence is under both
hypotheses, and never mistake the one conditional for its
reverse}.

\textbf{16.} $\dfrac{0.4 \times 0.15}{0.4 \times 0.15 +
0.6 \times 0.005} = \dfrac{0.060}{0.063} \approx 95\,\%$
spam: one word nearly settles it.

\textbf{17.} Five copies fail when $3$, $4$ or $5$ flip:
$\binom53 \varepsilon^3(1-\varepsilon)^2 + \binom54
\varepsilon^4 (1 - \varepsilon) + \varepsilon^5 \approx
0.0086$ --- against $0.028$ for three copies: repetition
buys reliability at the price of bandwidth.

\textbf{18.} Even flips: $\binom40 0.9^4 + \binom42 0.9^2
0.1^2 + \binom44 0.1^4 = 0.6561 + 0.0486 + 0.0001 = 0.7048$;
and $\frac{1 + 0.8^4}{2} = \frac{1 + 0.4096}{2} = 0.7048$:
the trick formula tracks the bias $(1 - 2\varepsilon)$
through each relay.

\textbf{19.} Spammy words travel in packs --- a message with
``lottery'' very likely contains ``win'', so their joint
evidence is weaker than the product pretends. The filter
still works because the exaggeration usually pushes in the
\emph{right direction}: classification needs the correct side
of $50\,\%$, not the correct posterior.

\textbf{20.} Clinic: a rare condition turns a strong test
into a weak verdict --- update from the base rate, then
accumulate. Courtroom: a rare match in a large city indicts
the city, not the man --- likelihoods are not verdicts.
Inbox: many weak clues, multiplied, sort the mail --- even
naively. The three traps: forget the prior and be fooled;
fake \omterm{def:g12:condprob:indep}{independence} and be very fooled; read evidence piecemeal
and never know what you know.
\end{solution}

\section*{Chapter \ref{ch:g12:randvar} --- Random Variables and the Binomial Distribution}

\begin{solution}{exo:g12:randvar:1}
$\E(X) = -0.3 + 0 + 0.8 + 0.5 = 1$.
$\E(X^2) = 0.3 \times 1 + 0 + 0.4\times4 + 0.1\times25 = 4.4$, so by
König--Huygens $\V(X) = 4.4 - 1 = 3.4$ and
$\sigma(X) = \sqrt{3.4} \approx 1.84$.
\end{solution}

\begin{solution}{exo:g12:randvar:2}
\emph{1.} The $10$ questions are independent Bernoulli trials with
$p = \frac14$: $X \sim \mathcal B\left(10, \frac14\right)$,
$\E(X) = 2.5$,
$\sigma(X) = \sqrt{10 \times \frac14 \times \frac34} = \sqrt{1.875}
\approx 1.37$.

\emph{2.} $\P(X = 0) = \left(\frac34\right)^{10} \approx 0.056$;
\[
\P(X = 5) = \binom{10}{5}\left(\frac14\right)^5\left(\frac34\right)^5
\approx 0.058;
\quad
\P(X \geq 1) = 1 - \left(\tfrac34\right)^{10} \approx 0.944 .
\]
\end{solution}

\begin{solution}{exo:g12:randvar:3}
Independent identical trials: $X \sim \mathcal B(400,\ 0.05)$, so
$\E(X) = 20$ claims and
\[
\sigma(X) = \sqrt{400 \times 0.05 \times 0.95} = \sqrt{19} \approx 4.4 .
\]
\end{solution}

\begin{solution}{exo:g12:randvar:4}
The payment $P$ received satisfies
$\E(P) = \frac{5 + 6}{6} = \frac{11}{6}$, so the fair price is
$c = \frac{11}{6} \approx 1.83$ euros. The fair gain $G = P - \frac{11}6$
takes values $-\frac{11}{6}, \frac{19}{6}, \frac{25}{6}$ with probabilities
$\frac46, \frac16, \frac16$:
\[
\V(G) = \E(G^2)
= \frac{4 \times 121 + 361 + 625}{6 \times 36} = \frac{1470}{216}
= \frac{245}{36} \approx 6.8,
\qquad \sigma(G) \approx 2.6 .
\]
A fair game has zero \emph{average} gain but its outcomes still fluctuate:
fair is not riskless.
\end{solution}

\begin{solution}{exo:g12:randvar:5}
$X \sim \mathcal B(8,\ 0.7)$.
\[
\P(X = 6) = \binom86 (0.7)^6(0.3)^2 \approx 0.296 ;
\]
\[
\P(X \geq 6) = \P(6) + \P(7) + \P(8)
\approx 0.296 + 8(0.7)^7(0.3) + (0.7)^8
\approx 0.296 + 0.198 + 0.058 = 0.552 ;
\]
$\P(X \geq 1) = 1 - (0.3)^8 \approx 0.99993$. Mode:
$(n+1)p = 6.3$, so the most probable value is $k^* = 6$
(\cref{exo:g12:randvar:7}).
\end{solution}

\begin{solution}{exo:g12:randvar:6}
The number of passengers showing up is $X \sim \mathcal B(104,\ 0.9)$; the
flight is overbooked when $X \geq 101$:
\[
\P(X \geq 101) = \sum_{k=101}^{104}\binom{104}{k}(0.9)^k(0.1)^{104-k}
\approx 0.006 .
\]
Selling $4\%$ more tickets than seats causes an incident on only about
$0.6\%$ of flights --- the economics behind overbooking.
\end{solution}

\begin{solution}{exo:g12:randvar:7}
\[
\frac{\P(X = k+1)}{\P(X = k)}
= \frac{\binom{n}{k+1}}{\binom nk}\cdot\frac{p}{1-p}
= \frac{n - k}{k + 1}\cdot\frac{p}{1-p},
\]
using $\binom{n}{k+1} = \binom nk \frac{n-k}{k+1}$. This ratio is $\geq 1$
iff $(n-k)p \geq (k+1)(1-p)$ iff $np - k p \geq k - kp + 1 - p$ iff
$k \leq (n+1)p - 1$. So the probabilities increase strictly while
$k + 1 \leq (n+1)p$ and decrease afterwards: the maximum is attained at
$k^* = \floor{(n+1)p}$ (shared with $k^* - 1$ when $(n+1)p$ is an
integer).
\end{solution}

\begin{solution}{exo:g12:randvar:8}
\emph{1.} Heads first at toss $k$ means $k-1$ tails then heads:
probability $\left(\frac12\right)^{k-1}\cdot\frac12 = 2^{-k}$, for
$1 \leq k \leq 10$. Total winning probability
$\sum_{k=1}^{10} 2^{-k} = 1 - 2^{-10} = \frac{1023}{1024}$ (the game is
lost only on ten consecutive tails).

\emph{2.} Expected payoff:
\[
\sum_{k=1}^{10} 2^k \cdot 2^{-k} = \sum_{k=1}^{10} 1 = 10 \text{ euros}.
\]
Without the cap, the sum $\sum_{k\geq1} 1$ diverges: the expected payoff is
infinite, although the game almost always pays a small amount --- the
famous \emph{Saint Petersburg paradox}, showing that expectation alone does
not measure the value of a game.
\end{solution}

\section*{Chapter \ref{ch:g12:sums} --- Sums of Random Variables and the Law of Large Numbers}

\begin{solution}{exo:g12:sums:1}
Each die has \omterm{def:g12:randvar:exp}{expectation}
$\frac{1 + 2 + \dots + 6}{6} = \frac{7}{2}$, so by linearity
$\E(S) = \frac72 + \frac72 = 7$. The dice are \omterm{def:g12:sums:indep}{independent}, so
$\V(S) = \frac{35}{12} + \frac{35}{12} = \frac{35}{6} \approx 5.83$.
\end{solution}

\begin{solution}{exo:g12:sums:2}
$\E(X) = 50$ and $\V(X) = 100 \times 0.25 = 25$.

Markov ($X \geq 0$): $\P(X \geq 75) \leq \frac{50}{75} = \frac23$.

Bienaymé--Chebyshev:
$\P(\abs{X - 50} \geq 25) \leq \frac{25}{25^2} = 0.04$ --- and this even
bounds the \emph{two-sided} \omterm{def:g11:prob:model}{event}, of which $X \geq 75$ is half. Chebyshev
is far sharper here because it exploits the \omterm{def:g12:randvar:exp}{variance}, not just the \omterm{def:g11:stat:mean}{mean}.
(The true value of $\P(X \geq 75)$ is smaller than $10^{-6}$: both bounds
are crude.)
\end{solution}

\begin{solution}{exo:g12:sums:3}
$F_n$ is the \omterm{prop:g12:sums:mean}{sample mean} of $n$ Bernoulli$\left(\frac12\right)$ variables,
of \omterm{def:g12:randvar:exp}{variance} $\sigma^2 = \frac14$. The bound
\[
\P\left(\abs{F_n - 0.5} \geq 0.05\right)
\leq \frac{1/4}{n \times 0.05^2} = \frac{100}{n}
\]
is $\leq 0.05$ as soon as $n \geq 2000$.
\end{solution}

\begin{solution}{exo:g12:sums:4}
By \cref{prop:g12:sums:mean}, $\E(M_n) = 5\%$ (diversification does not
change the expected return) and
$\sigma(M_n) = \frac{20\%}{\sqrt n}$. Requiring
$\frac{20}{\sqrt n} < 4$ gives $\sqrt n > 5$, \emph{i.e.} $n \geq 26$.
This is the principle of \emph{diversification}: spreading capital over
\omterm{def:g12:sums:indep}{independent} risks \omterm{def:g12:arith:divides}{divides} the risk (\omterm{def:g11:stat:variance}{standard deviation}) by $\sqrt n$
without reducing the expected return.
\end{solution}

\begin{solution}{exo:g12:sums:5}
\emph{1.} Counting the $9$ equally likely pairs: $S$ takes the values
$2, 3, 4, 5, 6$ with probabilities
$\frac19, \frac29, \frac39, \frac29, \frac19$. Hence
$\E(S) = \frac{2 + 6 + 12 + 10 + 6}{9} = 4$ and
$\E(S^2) = \frac{4 + 18 + 48 + 50 + 36}{9} = \frac{156}{9} = \frac{52}{3}$,
so $\V(S) = \frac{52}{3} - 16 = \frac43$.

\emph{2.} One variable: $\E(X) = 2$,
$\E(X^2) = \frac{1 + 4 + 9}{3} = \frac{14}{3}$,
$\V(X) = \frac{14}{3} - 4 = \frac23$. Then $\E(S) = 2 + 2 = 4$ and, by
\omterm{def:g12:condprob:indep}{independence}, $\V(S) = \frac23 + \frac23 = \frac43$. Same values.
\end{solution}

\begin{solution}{exo:g12:sums:6}
With $Y = X$: $\V(X + Y) = \V(2X) = 4\V(X)$, while
$\V(X) + \V(Y) = 2\V(X)$. The two agree only when $\V(X) = 0$, \emph{i.e.}
when $X$ is constant --- so additivity genuinely requires \omterm{def:g12:condprob:indep}{independence}
(here $X$ is maximally dependent on itself).
\end{solution}

\begin{solution}{exo:g12:sums:7}
Under fairness, $F_n$ is the \omterm{def:g11:stat:mean}{mean} of $n = 1200$ Bernoulli variables with
$p = \frac16$, $\sigma^2 = \frac16\cdot\frac56 = \frac{5}{36}$:
\[
\P\left(\abs{F_n - \tfrac16} \geq 0.05\right)
\leq \frac{5/36}{1200 \times 0.0025} = \frac{5}{108} \approx 0.046 .
\]
The observed deviation is exactly $0.05$: an \omterm{def:g11:prob:model}{event} that a fair die produces
with \omterm{def:g10:proba:distribution}{probability} at most $4.6\%$ --- and since Chebyshev is very
conservative, the true \omterm{def:g10:proba:distribution}{probability} is far smaller. The fairness hypothesis
is not plausible; the die is very likely loaded.
\end{solution}

\begin{solution}{exo:g12:sums:8}
\emph{1.} $p(1-p) = \frac14 - \left(p - \frac12\right)^2 \leq \frac14$.

\emph{2.} $F_n$ has \omterm{def:g12:randvar:exp}{expectation} $p$ and \omterm{def:g12:randvar:exp}{variance}
$\frac{p(1-p)}{n} \leq \frac{1}{4n}$; Bienaymé--Chebyshev gives
\[
\P\bigl(\abs{F_n - p} \geq \varepsilon\bigr)
\leq \frac{p(1-p)}{n\varepsilon^2} \leq \frac{1}{4n\varepsilon^2},
\]
valid whatever the unknown $p$.

\emph{3.} With $\varepsilon = 0.03$ and level $0.05$:
\[
\frac{1}{4n(0.03)^2} \leq 0.05
\iff n \geq \frac{1}{4 \times 0.0009 \times 0.05} \approx 5556 .
\]
About $5600$ people suffice by this crude bound (the classical
\omterm{def:g10:numbers:approx}{normal-approximation} answer is nearer $1100$, see
\cref{ch:g12:contdist}).
\end{solution}

\begin{solution}{pb:g12:sums:1}
\textbf{1.} $V(X + Y) = V(X) + V(Y) = \frac{35}{12} +
\frac{35}{12} = \frac{35}{6}$ (\omterm{def:g12:condprob:indep}{independence}).

\textbf{2.} $\E(\bar X) = 3.5$;
$\sigma(\bar X) = \frac{1.71}{\sqrt{100}} \approx 0.17$: the
hundred-dice average hugs $3.5$ within a fifth of a point.

\textbf{3.} $\P(X \geq 10) \leq \frac{2}{10} = 0.2$.

\textbf{4.} $\P \leq \frac{V(\bar X)}{0.5^2} =
\frac{35/1200}{0.25} \approx 0.117$: at most about
$12\,\%$.

\textbf{5.} With $Y = -X$: $V(X + Y) = V(0) = 0$, while
$V(X) + V(Y) = 2V(X) > 0$: perfectly anti-correlated variables
cancel --- \omterm{def:g12:randvar:exp}{variance} addition is an \omterm{def:g12:condprob:indep}{independence} privilege.

\textbf{6.} $\E = \frac{18}{37} - \frac{19}{37} =
-\frac{1}{37} \approx -0.027$;
$\sigma = \sqrt{1 - \left(\frac{1}{37}\right)^2} \approx
1.00$.

\textbf{7.} $\E(G) = -2.70$, $\sigma(G) = 10 \times 1.00 =
10$: the drift is only $0.27\sigma$. The night's noise dwarfs
the edge --- a large minority of gamblers (roughly $40\,\%$,
says the bell) walk out ahead, which is precisely what keeps
them coming back.

\textbf{8.} Casino's take: \omterm{def:g12:randvar:exp}{expectation} $+27\,027$ euros,
\omterm{def:g11:stat:variance}{standard deviation} $\approx 1\,000$: the profit sits $27$
\omterm{def:g11:stat:variance}{standard deviations} above zero. At $27\sigma$, ``the casino
might lose this year'' is not a risk, it is a rounding error.

\textbf{9.} $\P(\text{loss}) = \P(G \geq 27\,027$ for the
gamblers$)$ $\leq \frac{10^6 \times 1}{27\,027^2} \approx
0.0014$: even the bluntest inequality in the book guarantees
the house at $99.86\,\%$ --- the truth is astronomically
stronger.

\textbf{10.} Each euro staked, whenever and however chosen,
has expected return $-\frac{1}{37}$ of itself; by linearity
the expected total gain is $-\frac{1}{37} \times
(\text{total staked})$, negative for every strategy that
stakes anything. A winning system would violate the linearity
of \omterm{def:g12:randvar:exp}{expectation} --- no clever sequencing of bad bets makes a
good one. (Doubling systems merely trade many small wins for
rare catastrophic losses.)

\textbf{11.} The theorem: the \emph{\omterm{def:g10:stats:series}{frequency}} of red over
$n$ spins \omterm{def:g12:seq:limit}{converges} (in \omterm{def:g10:proba:distribution}{probability}) to $\frac{18}{37}$. It
says nothing about spin $n + 1$: the wheel has no memory, and
after ten reds the chance of red is still $\frac{18}{37}$.
Believing otherwise is the \emph{gambler's fallacy}, and
question 12 shows what actually happens to streaks.

\textbf{12.} The surplus of $100$ reds is not repaid ---
future spins are fair copies, expected surplus stays $100$.
But $\frac{100}{10^6} = 0.0001$: one hundredth of a point.
The fallacy's error in one sentence: \emph{the law of large
numbers dilutes past accidents in an ocean of new trials; it
never sends the wheel to collect debts.}

\textbf{13.} $n \geq \frac{1}{4 \times 0.05 \times 0.03^2}
\approx 5\,556$ people.

\textbf{14.} $\frac{1.96^2 \times 0.25}{0.03^2} \approx
1\,067$: five times fewer. Chebyshev holds for \emph{every}
\omterm{def:g12:randvar:rv}{distribution} and pays for its universality with slack; the
pollsters' constant $1.96$ comes from the actual bell shape
of the fluctuations, which concentrates far harder.

\textbf{15.} Margin $\propto \frac{1}{\sqrt n}$: halving it
quadruples the \omterm{def:g10:proba:sample}{sample} --- about $4\,400$ people for
$\pm 1.5$ points. Precision is bought at quadratic prices.

\textbf{16.} The model is $n$ \omterm{def:g12:sums:indep}{independent} draws with success
\omterm{def:g10:proba:distribution}{probability} $p$ --- the population size $N$ never appears
(sampling a tiny fraction of a large, well-mixed population
is what \omterm{def:g12:condprob:indep}{independence} encodes). The pollsters' analogy: to
taste the soup, one well-stirred spoonful suffices ---
whether the pot serves ten or ten thousand. The catch is the
stirring --- a biased sampling frame, as in the Middle School
volume's poll that fooled a country --- not the pot.

\textbf{17.} $315\sqrt n < 30n \iff \sqrt n > 10.5 \iff
n > 110$: beyond a hundred policies the drift outruns the
noise, and each further policy widens the gap by the
$\sqrt n$ law --- pooling \emph{is} the business model.

\textbf{18.} \omterm{def:g12:sums:indep}{Independent} equal shares: $\sigma_{\text{pf}} =
\frac{20\,\%}{\sqrt{25}} = 4\,\%$: same expected return,
one-fifth the fluctuation. The lunch: risk reduction at zero
cost in \omterm{def:g12:randvar:exp}{expectation} --- averaging \omterm{def:g12:sums:indep}{independent} randomness.

\textbf{19.} Perfectly correlated: the portfolio is one stock
in $25$ costumes: $\sigma = 20\,\%$, no reduction at all.
Diversification rents \emph{\omterm{def:g12:condprob:indep}{independence}}; when a crisis
correlates everything (2008: all housing bets were one bet),
the $\sqrt n$ protection evaporates exactly when it is
needed.

\textbf{20.} Sums drift like $n$ and fluctuate like
$\sqrt n$, so means settle like $\frac{1}{\sqrt n}$: the
casino banks the drift over a million spins; the pollster
buys $\frac{1}{\sqrt{1100}}$ of noise for a phone budget;
the insurer outgrows its own claims at $n > 110$; the
investor \omterm{def:g12:arith:divides}{divides} risk by $\sqrt{25}$. Abusers: the gambler
who believes in debts (there is only dilution), and the
financier who believes in \omterm{def:g12:condprob:indep}{independence} (there is sometimes
only one bet). The fluctuations' universal shape --- the
bell --- is the \omterm{def:g12:limcont:continuity}{continuous} chapter's star, and its theorem,
the central limit theorem, is the summit of the university
volumes' \omterm{def:g10:proba:distribution}{probability}.
\end{solution}

\section*{Chapter \ref{ch:g12:contdist} --- Continuous Random Variables}

\begin{solution}{exo:g12:contdist:1}
$\P(X \leq 5) = \frac{5}{15} = \frac13$;
$\P(4 \leq X \leq 10) = \frac{6}{15} = \frac25$;
$\E(X) = \frac{15}{2} = 7.5$ min;
$\sigma(X) = \frac{15}{\sqrt{12}} = \frac{5\sqrt3}{2} \approx 4.3$ min.
\end{solution}

\begin{solution}{exo:g12:contdist:2}
$f \geq 0$ and
$\int_0^2 \frac38 t^2\,\dd t = \frac38\left[\frac{t^3}{3}\right]_0^2
= \frac38 \times \frac83 = 1$: a \omterm{def:g12:contdist:density}{density}.
\[
\E(X) = \int_0^2 \frac38 t^3\,\dd t = \frac38 \times \frac{16}{4} = \frac32,
\qquad
\P(X \geq 1) = \int_1^2 \frac38 t^2\,\dd t = \frac{8 - 1}{8} = \frac78 .
\]
\end{solution}

\begin{solution}{exo:g12:contdist:3}
\emph{1.} $\E(X) = \frac{1}{0.2} = 5$ years;
$\P(X > 5) = \eu^{-0.2\times5} = \eu^{-1} \approx 0.37$.

\emph{2.} By memorylessness (\cref{thm:g12:contdist:memoryless}),
$\pcond{X > 3}{X > 8} = \P(X > 5) = \eu^{-1} \approx 0.37$: the three years
of service change nothing.
\end{solution}

\begin{solution}{exo:g12:contdist:4}
$\P(X > m) = \eu^{-\lambda m} = \frac12$ gives
$m = \frac{\ln 2}{\lambda} \approx \frac{0.69}{\lambda}$, smaller than
$\E(X) = \frac1\lambda$. The exponential \omterm{def:g12:contdist:density}{density} has a long right tail: a
few atypically long lifetimes pull the \emph{\omterm{def:g11:stat:mean}{mean}} above the
\emph{\omterm{def:g11:stat:median}{median}}, the value that half the population exceeds. (This is the
half-life of \cref{ex:g12:exp:halflife}: half the atoms survive it, even
though the average lifetime is longer.)
\end{solution}

\begin{solution}{exo:g12:contdist:5}
$168 = \mu - \sigma$ and $182 = \mu + \sigma$: about $68\%$.
$189 = \mu + 2\sigma$: above it lies half of the remaining
$100 - 95.4 = 4.6\%$, so about $2.3\%$.
$154 = \mu - 3\sigma$: about $\frac{100 - 99.7}{2} = 0.15\%$.
\end{solution}

\begin{solution}{exo:g12:contdist:6}
$\P(X < 500) \leq 2.5\%$ means $500$ must sit at least two \omterm{def:g11:stat:variance}{standard
deviations} below the \omterm{def:g11:stat:mean}{mean} (the normal leaves $2.3\% \approx 2.5\%$ below
$\mu - 2\sigma$; with the conventional $1.96$, the reasoning is identical):
\[
\mu - 2\sigma \geq 500 \iff \mu \geq 500 + 8 = 508 \text{ g}.
\]
The machine must be set to $508$ g on average --- the price of the
guarantee is $8$ g of ``free'' product per bag.
\end{solution}

\begin{solution}{exo:g12:contdist:7}
\emph{1.} With $p = \frac16$, $n = 1200$:
$\sqrt{\frac{p(1-p)}{n}} = \sqrt{\frac{5/36}{1200}} \approx 0.0108$, so the
$95\%$ fluctuation \omterm{def:g10:numbers:interval}{interval} is
\[
\frac16 \pm 1.96 \times 0.0108 \approx \intcc{0.146}{0.188}.
\]

\emph{2.} The observed \omterm{def:g10:stats:series}{frequency} $\frac{260}{1200} \approx 0.217$ lies far
outside: the fairness hypothesis is rejected at the $5\%$ level. The
deviation is about $4.6$ \omterm{def:g11:stat:variance}{standard deviations} --- for a normal
\omterm{def:g10:numbers:approx}{approximation}, a \omterm{def:g10:proba:distribution}{probability} of order $10^{-6}$, far more conclusive than
the $\leq 4.6\%$ bound from Bienaymé--Chebyshev
(\cref{exo:g12:sums:7}).
\end{solution}

\begin{solution}{exo:g12:contdist:8}
\emph{1.} Margin $\frac{1}{\sqrt n} \leq 0.02$ requires
$n \geq 2500$ people.

\emph{2.} To distinguish scores $1$ point apart, the margin must be well
under $0.5$ point, requiring $n \geq \left(\frac{1}{0.005}\right)^2 =
40\,000$ by the simplified formula --- already impractical for most polls.
Worse, the statistical margin only accounts for \emph{sampling} error;
systematic biases (unrepresentative \omterm{def:g10:proba:sample}{samples}, non-response, last-minute
swings) do not shrink as $n$ grows and typically exceed one point. A
$1$-point race is genuinely too close to call.
\end{solution}

\begin{solution}{exo:g12:contdist:9}
\emph{1.} Since $\alpha > 0$, $c \leq \alpha X + \beta \leq d
\iff \frac{c - \beta}{\alpha} \leq X \leq \frac{d - \beta}{\alpha}$, so
\[
\P(c \leq Y \leq d)
= \int_{(c-\beta)/\alpha}^{(d-\beta)/\alpha} f(t)\,\dd t .
\]
The substitution $y = \alpha t + \beta$ (i.e.\ reading the area in the $y$
variable, $t = \frac{y - \beta}{\alpha}$, $\dd t = \frac{\dd y}{\alpha}$)
turns this into
$\int_c^d \frac{1}{\alpha} f\left(\frac{y-\beta}{\alpha}\right)\dd y$: the
\omterm{def:g11:func:function}{function} $g(y) = \frac1\alpha f\left(\frac{y-\beta}{\alpha}\right)$ is a
\omterm{def:g12:contdist:density}{density} of $Y$.

\emph{2.} With $f = \varphi$, $\alpha = \sigma$, $\beta = \mu$:
\[
g(y) = \frac{1}{\sigma}\,\varphi\!\left(\frac{y - \mu}{\sigma}\right)
= \frac{1}{\sigma\sqrt{2\pi}}\,
\exp\left(-\frac{(y-\mu)^2}{2\sigma^2}\right),
\]
which is therefore the \omterm{def:g12:contdist:density}{density} of
$\mathcal N(\mu, \sigma^2) = \mu + \sigma\,\mathcal N(0,1)$.
\end{solution}



\backmatter
\printindex

\end{document}
